\documentclass[../main.tex]{subfiles}
\begin{document}
%==========================================TIÊU ĐỀ TẬP========================================
\begin{table}[h]
  \begin{adjustwidth}{-1cm}{}
    \begin{tabular}{>{\centering\arraybackslash}p{6cm}|>{\raggedright\arraybackslash}p{14cm}}
      \multirow{5}{6cm}
      % Cột bên trái: ÉPISODE và số tập
      {\\[-40pt]
        \flaregothic\fontsize{20pt}{20pt}\selectfont \centering
        \color{eptype}HISTOIRE\color{black}\\
        \brian\fontsize{72pt}{72pt}\selectfont
        \color{epnum}4
        \\[18pt]
        % \flaregothic\fontsize{14pt}{14pt}\selectfont
        % \color{epattr}BIENVENUE, \uline{\color{Banpire} Mel} !
        % \flaregothic\fontsize{18pt}{18pt}\selectfont
        % \color{epattr}ÉPISODE PERDU
        \color{black}
      }
       &
      % Dòng thứ nhất: tên tập tiếng Nhật
      {\fotskip\fontsize{18pt}{18pt}\selectfont パーティー前の\ruby{準}{じゅん}\ruby{備}{び}大\ruby{混}{こん}\ruby{乱}{らん}！} \\[6pt]
      % Dòng thứ hai: tên tập tiếng Anh
       & {\cafeteria\fontsize{18pt}{18pt}\selectfont Preparation Chaos Before the Party!}                  \\[4pt]
      % Dòng thứ ba: tên tập tiếng Việt
       & {\vnmsans\fontsize{18pt}{18pt}\selectfont Buổi chuẩn bị hỗn loạn trước giờ G!}                    \\[8pt]
      % Dòng thứ tư: Nhân vật xuất hiện
       & {\sen Track Used: Birthday Party (\ruby{誕}{たん}\ruby{生}{じょう}パーティー)}
    \end{tabular}
  \end{adjustwidth}
\end{table}
%===========================================NỘI DUNG==========================================
\color{black}

\pov{Hikawa Kuon}

Mẹ tôi hôm nay vẫn bận đi giao hàng cho kho. Dù là thứ Bảy hay Chủ Nhật thì với bà cũng chẳng khác gì ngày thường, công việc vẫn cứ tiếp diễn như một thói quen đã ăn sâu vào cuộc sống.

Còn em gái tôi, vừa về đến nhà sau khu tan học ca chiều sớm, vả lại cũng khá mệt mỏi. Thế nên, chỉ còn lại hai bố con tôi cùng nhau đi chợ.

Thật ra, tôi cũng không có gì để phàn nàn. Những buổi đi chợ với bố chẳng có gì đặc biệt, nhưng lại mang một cảm giác quen thuộc đến lạ. Không cần tới siêu thị, cũng chẳng cần đến cửa hàng tiện lợi, mà chỉ ghé qua chợ dân sinh cách nhà không xa, nơi hai bố con vẫn hay ghé qua mỗi khi cần mua đồ tươi sống. Tôi chẳng có yêu cầu gì đặc biệt về bữa ăn, lúc nào cũng chỉ đơn giản: ``Ăn gì cũng được,'' tất nhiên là ngoại trừ rau củ ra.

Và vì cái tính ``dễ ăn dễ nuôi'' ấy, theo lời bố tôi, mỗi lần đi chợ chỉ mất đâu đó năm đến mười phút là xong.

Hôm nay cũng vậy. Chúng tôi lướt qua mấy sạp hàng quen, chọn vài lát thịt bò tươi, một ít rau xanh, vài gói nấm và đậu phụ. Mấy thứ này tối nay sẽ được nấu thành một nồi lẩu nghi ngút khói, một món ăn hơi đặc biệt và khác thường trong mùa hè oi ả giữa tháng Bảy. Một bữa tối chỉ có tôi, em gái tôi, bố mẹ tôi, và nếu có thể, mời thêm cả bác Hijaki-san cùng bà nội tôi có thể chung vui với gia đình.

Không cần quá xa hoa, cũng không cần quá đông đúc. Chỉ là một bữa tối nhỏ đơn giản và đấm ấm.

Khi tôi nhìn quanh, chợ không quá đông, nhưng vẫn nhộn nhịp với những tiếng rao hàng, tiếng cười nói của mấy bà cô bán rau, hay tiếng xe cộ vụt qua trên con đường gần đó. Một không khí bình thường, chẳng có gì nổi bật, nhưng lại đem đến một sự yên bình khó tả.

Mua đồ xong, hai bố con tôi quay về. Tôi ngồi sau xe máy, tai vẫn bịt tai nghe như mọi khi, nghe một bài nhạc quen thuộc mà mình hay bật mỗi khi đi đường.

Khi đến trước cửa nhà, bố tôi dừng xe, khẽ nghiêng đầu ra sau gọi tôi: ``\vie{Con xách đồ lên nhà trước đi.}''

Tôi bỏ tai nghe xuống một bên, đáp gọn lỏn: ``\vie{Dạ.}''

Hai túi đồ được tôi xách lên một cách gọn gàng, không để ý gì nhiều đến xung quanh. Chỉ là một buổi đi chợ bình thường như bao ngày khác - ít nhất, đó là những gì tôi nghĩ vào lúc này.

\ct

Tôi xách hai túi đồ lên nhà, bước vào cửa rồi đi thẳng lên cầu thang. Căn nhà này hiện có bảy tầng, mà bếp lại nằm tận tầng sáu, nên việc mang đồ lên chẳng khác gì một bài tập thể lực. Nhưng tôi đã quen rồi, mỗi lần đi chợ, đi học hay đi làm về cũng đều như vậy.

Trên đường lên, tôi ghé qua tầng hai, dừng lại trước cánh cửa phòng của Aloe-chan. Cô nàng không có nhà, chắc vẫn còn đang làm việc ở đâu đó xa lắm, tận sáng mai mới về. Dám chắc là đi hát ở một chỗ nào đó mà mẹ tôi chỉ đây.

Tôi gõ cửa theo thói quen, dù biết rõ bên trong chẳng có ai, rồi nói bâng quơ: ``Aloe-chan, tối mai nhớ về sớm một chút nhé.''

Không có tiếng trả lời. Dĩ nhiên rồi. Tôi nhún vai, quay lưng tiếp tục đi lên.

Đến tầng bốn, tôi dừng lại trước cửa phòng của bác Hijaki và bà nội. Hai người họ lúc nào cũng ở nhà, ngoại trừ những khi bác phải ra ngoài làm việc.

Tôi mở cửa cánh cửa trượt, rồi cất giọng nhẹ nhàng: ``\vie{Bác Hijaki-san, bà. Tối nay có tiệc sinh nhật của cháu, cháu mời hai người ạ.}''

Bác nằm ở trên giường, trông thấy tôi, mỉm cười đáp: ``\vie{Ừ, ừ, tao biết rồi. Bọn tao không có quên đâu mà sợ.}''

Bà nội tôi cũng ngó ra, gật đầu: ``\vie{Lát bà bảo bác con chuẩn bị quà.}''

Tôi bật cười đáp lại: ``\vie{Bà cứ đến ăn là được rồi, không cần quà đâu ạ.}''

Tôi thì vốn dĩ không quý bà nội lắm, có lẽ bởi vì khoảng cách thế hệ mà tư tưởng của hai người chúng tôi khác nhau. Với cả, gần như đêm nào bà cũng lọ mọ lên trên tầng, tới mức làm phiền con cái - hoặc ít nhất, tôi cho là vậy.

Thành ra, việc mời bà nội xuống nhà ăn sinh nhật với tôi làm cho tôi cảm tưởng đó là một nghĩa vụ hơn là một lời mời từ tâm. Nhưng biết sao được? Tôi thì cũng không quá bận tâm việc đó, miên là ai đó dỗ dành bà ăn cơm là được.

Sau khi mời xong, tôi lại tiếp tục lên trên nhà, đặt đống đồ ăn lên bếp, rồi ngó sang tầng bảy, phòng của Coco-chan.

Cô nàng cũng không có nhà, chắc lại bận một công việc nào đó. Nghe phong thanh thì có vẻ như nàng Rồng đang chuẩn bị ký hợp đồng với một công ty VTuber mới sau khi rời khỏi \hll, mà ký lại ở tận bên Hoa Kỳ cơ. Tôi cũng chẳng hỏi kỹ làm gì, chỉ biết em ấy đang cố gắng theo cách của mình.

Tôi nghĩ rằng em ấy hôm nay cũng sẽ không về nhà tối nay đâu. Chán thật.

Không có gì để nói thêm, tôi nhanh chóng quay xuống tầng năm, cũng là phòng của hai bố con tôi. Đẩy cửa bước vào, tôi nằm sấp xuống đệm và bật máy tính lên. Công việc vẫn còn đó, mấy email cần xử lý, một số tài liệu cần kiểm tra.

Tôi đeo tai nghe vào, mở một danh sách nhạc quen thuộc, để nhạc chạy nền trong khi làm việc.

Tôi liếc sang phải, nhìn thấy Ryuuhou đang ngồi xem vô tuyến, hoặc đúng hơn, lướt Y**Tube. Cô em gái đang xem chương trình dành cho trẻ con kia, dù đã mười sáu, mười bảy tuổi rồi.

``\vie{Vẫn xem cái chương trình cổ lỗ sĩ đó sao?}'' tôi cất giọng, tay vẫn gõ bàn phím.

``\vie{Vâng. Nhưng mà nó hay mà,}'' em tôi đáp lại, thi thoảng còn bật cười.

Tôi thở dài một tiếng. ``\vie{Thế bài tập bài tủng thế nào rồi?}''

``\vie{Làm ở trên trường hết rồi,}'' Ryuuhou thản nhiên đáp.

Tôi lắc đầu, lầm bầm: ``\vie{\hl{Cẩn thận không lại ăn thêm một cái tin nhắn thông báo thiếu bài tập đấy\dots}}'' và tiếp tục với công việc. Thi thoảng, hai anh em chúng tôi cứ thế mà nói qua nói lại vài câu vô thưởng vô phạt, rồi ai làm việc nấy.

Tôi cũng chẳng mảy may để ý xem bên dưới đang có chuyện gì. Vì với tôi, đây chỉ là một ngày bình thường như bao ngày khác.

\dots Có lẽ vậy.

\pov{-}

Kuon vừa khuất khỏi tầm mắt trên cầu thang, Ryujin cũng vừa rút chìa khóa xe, chuẩn bị lên nhà cùng cậu con trai.

Nhưng đúng lúc đó, một giọng nói vang lên từ phía sau: ``A-nô\dots\ Chú là Hikawa Ryujin-san phải không ạ?''

Người đàn ông cơ khí nhướng mày, quay đầu lại. Trước mắt ông là một cảnh tượng kỳ lạ: có khoảng ba mươi cô gái đang đứng ngay trước cửa nhà mình, với những người quen thuộc, nhưng lại cũng có người mà ông chưa gặp bao giờ. Hoặc có lẽ ông vẫn đang hơi choáng ngợp vì số lượng lớn các cô gái ở đây.

Nhưng trước khi kịp hỏi, ông nhanh chóng nhận ra phía sau họ là bốn chiếc xe - bốn chiếc ``xe điên'' vừa làm náo loạn cả Kawachi trên đường đến đây. Mỗi chiếc xe, theo cách nào đó, đều đã ``hóa trang'' để trốn tránh sự chú ý.

Miosha, vốn hóa thân của Mio, đã lặng lẽ biến mất, để lại một Mio bình thường. Không gì khác, là một cô nàng Sói đen đang đứng tươi cười giữa nhóm.

Chiếc Lambo*ghini của Okayu thì sau màn phóng đi như tên lửa đã hoàn toàn biến mất trên bầu trời, như thể chưa từng tồn tại.

Xe quân sự của nhóm NePoLaBo đã được ngụy trang thành một bụi cây bên đường. Nhìn thoáng qua, người ngoài chắc chắn sẽ không nhận ra sự khác biệt.

Còn chiếc xe cảnh sát của Fubuki? Nó đã bị thu nhỏ thành một món đồ chơi, gọn lỏn trong tay nàng Cáo trắng, như thể chẳng có gì đáng nghi cả.

Ryujin nhìn cảnh tượng trước mặt, trong lòng ông thoáng một chút quen thuộc. Ông đưa tay gãi cằm, nghi hoặc hỏi: ``Các cháu là\dots?''

Từ trong nhóm, nàng SÓi đen và nàng Dị giới bước lên trước. Họ nhìn người đàn ông cơ khí kia, mỉm cười: ``Chú có nhớ bọn cháu không ạ?''

Người đàn ông nheo mắt, nhìn kỹ họ một lúc lâu, cố gắng lục lại trí nhớ. Sau vài giây, ông bất giác thốt lên: ``À, đúng rồi! Hai đứa này\dots\ với cả cháu tóc xanh nào đó mấy tháng trước đến nhà chú uống rượu đây mà!''

Nghe đến hai chữ ``uống rượu'', Lamy lập tức giơ tay lên hào hứng: ``Dạ, là cháu! Cháu vẫn chưa chịu thua chú sau cái buổi đó đâu!''

Ryujin gật gù: ``Hôm đó cháu uống rượu như điên, cuối cùng là vẫn gục đấy thôi. Hôm nay lại có gan tái đấu à?''

Nàng Yêu tinh tuyết phồng má lên giận dỗi: ``Đ-Đấy là vì cháu áp dụng sai chiến thuật thôi! Cháu tin lần này cháu sẽ thắng!''

Người đàn ông bật cười. ``Được rồi, để xem hôm nay cháu giở bài vở gì đây.'' Ông quay về phía hai cô nàng vừa nãy, ngờ ngợ hỏi: ``Hình như là Mio-chan với cháu tóc vàng nào đấy cũng sang nhà chú hôm đó đúng không?''

Mio và Flare đồng loạt giơ tay xác nhận. Thậm chí, cô nàng Yêu tinh vàng còn tự giới thiệu bản thân. ``Chú còn bảo lần sau đến chơi nhớ báo trước mà\dots\ Ừm, nhưng mà bọn cháu quên mất.''

Ông tiếp tục nhìn các cô gái còn lại trong đoàn, rồi tiếp tục kể lại: ``Chú nhớ ra rồi. Các cháu có phải là những người đã đón Tết năm ngoái không?''

Ngay lập tức, một loạt giọng nói náo loạn từ đa số các thành viên đều cùng chen chúc nhau nói:

``Chú nhớ ra bọn cháu rồi sao-nanodesu!?'' từ Rushia, ``Bọn cháu đã sang đây ăn Tết từ năm ngoái ạ!'' từ Mel, và thậm chí là giọng nói trong trẻo từ Haato: ``Chính cháu đã bắt Roboco-senpai với Noel-chan đền ổ điện nhà chú đấy ạ!''

Tất cả những thành viên đã từng đến nhà Kuon dịp năm mới năm ngoái đều đồng loạt nhao nhao lên, vô cùng phấn khích khi Ryujin vẫn còn nhớ. Ai nấy đều cười nói rôm rả, tạo nên một khung cảnh vô cùng ồn ào khiến cho những người đi đường tỏ vẻ khó hiểu, thậm chí có người còn tưởng đó là đám trẻ con nào đó tới nhà hàng xóm chơi.

Giữa đám đông đang hào hứng đó, Suisei đứng khoanh tay, phồng má một cách đầy ghen tị: ``Mồ\~{}! Tại sao mình lại không được mời chứ!?''

Bên cạnh, Marine cũng thể hiện sự bất mãn không kém: ``Em cũng có được mời đâu! Bất công quá!''

``Thì lúc đó hai đứa đã lên văn phòng đâu,'' A-chan lên tiếng, mỉm cười giải thích. Dù vậy, nàng Ca sĩ và nàng Thuyền trưởng vẫn tỏ vẻ không vui.

Một số thành viên khác từ Thế hệ thứ Tư và Thế hệ thứ Năm cũng có cảm xúc tương tự. Một số thì cảm thấy có chút ghen tị, trong khi một số lại tỏ ra ngạc nhiên vì không hề biết chuyện này.

Nhưng ngay sau đó, Luna đột nhiên lên tiếng với giọng trẻ con đặc trưng: ``Cơ mà, Sui-chan cũm từng sang nhà Kuonsa-senpai ăn mì ramen dịp lễ Vu Lan năm ngoái\footnote{Xem lại phần {\abv Truyện phụ}, {\flaregothic ÉPISODE 66}, quyển số Năm.} rùi nà, đúng hơm-nora\~{}?''

Suisei chớp mắt, khẽ ngạc nhiên: ``Hả? Thật à?''

Luna nghiêng đầu. ``Nhớ mà\~{}! Hôm đó chị còn nói là `Mì mẹ của Kuonsa-senpai ngon ghê á!' nữa đó-nanora.''

Nàng Ca sĩ chớp mắt thêm lần nữa, rồi trầm ngâm suy nghĩ. Có vẻ như cô đã quá mệt vào ngày hôm đó nên chẳng còn nhớ nổi. Nhưng dù không nhớ, cô vẫn cảm thấy không được mời ăn Tết là một sự thiệt thòi lớn.

Marine nghe vậy thì càng thêm ghen tị, bĩu môi lầm bầm: ``Hừm\dots\ Mọi người ai cũng từng đến nhà Kuon-senpai ít nhất một lần rồi\dots\ Còn mình thì không\dots''

Ryujin nhếch mép cười khi thấy các cô gái nhà \hll\ phấn khích với chuyện Tết. Ông khoanh tay lại, gật gù nói: ``Đúng rồi đấy. Chắc có mấy đứa ở đây không biết đâu, nhưng cái Tết năm đó đúng là một bộ phim tài liệu sống động! Cả Kawachi, Fushuu và Neihei đều được chứng kiến sự náo loạn do mấy đứa gây ra!''

Những ai từng tham gia năm ngoái bật cười, trong khi những người mới thì tò mò hơn. Với các thành viên như Sora, Aki hay Miko, họ làm sao quên được những tình huống trớ trêu khi lau dọn nhà cửa, chơi những trò dân gian, hay thậm chí là đi chùa lễ hội.

Nàng Thần tượng tươi cười, nhẩm lại những sự kiện đã qua thời điểm đó: ``Quả thật chúng cháu đã có một cái Tết `đặc biệt' ở nhà chú. Bọn cháu đã giúp dọn dẹp nhà cửa, làm bánh chưng, đi chúc Tết và còn ghi hình lại nữa.''

``Và bộ phim được đánh giá rất cao đó, ne!'' nàng Vu nữ chiêm thêm đầy hào hứng.

``Khoan, khoan,'' Towa thốt lên, ``làm bánh chưng ư?'' Nhưng rồi, cô nàng, cùng với các thành viên còn lại của thế hệ thứ Tư đều cùng đồng loạt nghiêng đầu: ``Mà, bánh chưng là gì vậy?''

``Nói ra thì cũng dài dòng lắm đấy\dots'' Mel lên tiếng. ``Để khi nào chị kể lại nhé.''

Nàng Dị giới vàng thở dài, tỏ vẻ tiếc nuối: ``Thế mà cháu không có mặt! Lần đó chắc vui lắm nhỉ?''

Người đàn ông bật cười, rồi tiếp tục kể lại những sự việc oái oăm trong mùa Tết năm đó: ``Vui thì vui, nhưng cũng mệt lắm đấy! Đặc biệt là lúc gói bánh chưng, có đứa gói vuông vức như trong sách, có đứa thì\dots\ thôi khỏi nói. Chắc trong đám này cũng có đứa gói bánh mà nhìn như viên gạch ấy nhỉ?''

Nghe vậy, Ayame, Roboco và một số thành viên khác đỏ bừng mặt. Đúng là những người đã gói những chiếc bánh chưng độc lạ.

Ông tiếp tục: ``Chú còn nhớ hồi sáng mùng Một ấy, hai bố con chú còn đi chúc Tết nữa, cơ mà chú nhớ rất rõ là Fubu-chan cứ lăm lăm lì xì thôi! Thằng Cò cứ làu bàu mãi.''

Nàng Cáo trắng nghe vậy liền đỏ mặt, ngượng ngùng đáp lại: ``Đ-Đúng là cháu có ham lì xì thật\dots''

``Nhưng ít nhất thì bà cũng lì xì lại cho người ta rồi chứ?'' nàng Ma cà rồng cười nhẹ.

Và cứ vậy, Ryujin cùng vài thành viên đã tham dự năm ngày Tết náo loạn kia đầy nhiệt huyết. Những thành viên cũ thì trầm ngâm nhớ lại những kỷ niệm vừa đáng nhớ vừa đáng quên, còn những thành viên mới (như Marine, Suisei và các thành vien Thế hệ thứ Tư và thứ Năm) thì có người bày tỏ sự tiếc nuối, có người lại càng  lúc càng tò mò về cái Tết có một không hai này, và trong lòng dần xuất hiện một mong muốn nho nhỏ: rằng có lẽ họ cũng nên thử trải nghiệm nó một lần.

``Giờ cháu mới biết Kuon-senpai mời các senpai đến nhà ăn Tết đó ạ,'' nàng Cừu cười mỉm, dù trong lòng cô dấy lên một chút ghen tị.

``Em không biết nhà anh ấy như thế nào, nhưng mà nghe có vẻ vui đấy!'' Polka thảng thốt.

Lamy cũng gật đầu, rên rỉ: ``Tôi cũng muốn như vậy nữa\~{}''

Botan vỗ vai cô bạn của mình, giọng bằng bằng: ``Bình tĩnh nào, Yukihana. Tôi nghĩ nó cũng giống với ngày đón năm mới ở chỗ ta thôi.'', dù rằng đằng sau cũng ẩn giấu một chút tò mò.

Ryujin hướng mặt về phía Senchou, người vẫn đang phồng má phụng phịu: ``Mà, cháu tóc đỏ kia thì\dots\ Marine-chan đúng không? Hình như chú từng gặp cháu rồi nhỉ?''

``Đúng rồi đó chú!'' nàng Thuyền trưởng phấn khởi. ``Lúc mà kỷ niệm 1 năm \textit{hologra}\footnote{Xem lại {\flaregothic SPÉCIAL} {\ab VII}, quyển số Bốn.} đó ạ!''

Với các thành viên còn lại mà ông chưa biết đến, ông nhẹ nhàng hỏi từng người một, và được các thành viên giới thiệu một cách sôi nổi, toát lên tính cách và cá tính của con người họ, như cách mà họ đã thể hiện khi họ lần đầu lên văn phòng.

Quả thật, cảnh tượng này trông như thể những đứa con xa quê được về nhà sau một khoảng thời gian rất dài.

Giữa lúc cả nhóm đang trò chuyện rôm rả, Ryujin chậm rãi tiến đến gần A-chan, một người bình thường mà ông đã quá quen thuộc. Với giọng điệu trầm ấm nhưng vẫn mang nét vui vẻ thường ngày, ông hỏi: ``Dạo này cháu thế nào rồi, A-chan?''

Nàng Đeo kính khẽ cúi đầu đáp lễ phép: ``Cháu khỏe ạ. Cảm ơn chú đã hỏi ạ.''

Người đàn ông cơ khí gật gù, ánh mắt sắc sảo lướt qua mọi người trước khi tiếp tục: ``Mọi người ở công ty vẫn ổn chứ?''

A-chan mỉm cười: ``Vâng ạ! Mọi người đều làm rất tốt. Ai cũng đang nỗ lực để trở thành idol thực thụ.''

Nghe vậy, ông bật cười đầy ẩn ý: ``Chú cứ tưởng các cháu định theo đuổi sự nghiệp diễn viên hài cơ đấy? Con nhà chú cứ than vãn suốt.''

A-chan chớp mắt, hơi lúng túng trước câu châm chọc đầy bất ngờ: ``À thì\dots\ cháu\dots''

Nhưng rồi, ông chỉ cười xòa: ``Chú đùa thôi. Dù gì thì chú cũng đoán được lý do vì sao các cháu lại đến đây rồi. Hình như các cháu có gọi điện cho chú từ hồi sáng, đúng không?''

Nàng Thần tượng tươi cười gật đầu: ``Vâng ạ! Bọn cháu đã trao đổi với chú một chút về kế hoạch của tụi cháu cho ngày hôm nay.''

Ryujin nhướng mày, khoanh tay trước ngực, ánh mắt ông gợn lên một chút hiếu kỳ: ``Thế mấy đứa định làm thế nào?''

Nàng Thần tượng kể lại kế hoạch bí mật mà cả nhóm đã vạch ra. Nhưng dường như Ryujin đã nắm rõ ngay từ đầu, bởi khi A-chan còn chưa kịp nói hết, ông đã kết luận ngay với một nụ cười đầy ẩn ý: ``Vậy là mấy đứa định tạo bất ngờ cho con trai chú đúng không?''

Tất cả đều há hốc miệng vì ông ấy nắm bắt rất nhanh, và rồi nhanh chóng gật đầu lia lịa. Ông lại bật cười: ``Chú biết con trai chú thế nào mà.''

Ông tiếp tục: ``Vậy mấy đứa đã chuẩn bị gì rồi?''

Đáp lại câu hỏi ấy, hàng loạt túi đồ nặng trịch được giơ lên, từ nguyên liệu làm bánh cho tới những đồ tươi sống làm cho người đàn ông nhà Hikawa phải choáng ngợp.

``Nhiều thế! Chỗ này chắc phải ăn được cả tháng đấy chứ đùa!'' Ryujin khẽ trố mắt, suýt bật ngửa xuống đất.

Noel phấn khích hô lớn: ``Những món liên quan đến bò ạ!''

Pekora lập tức chỉnh lại: ``Cụ thể hơn là lẩu bò và thịt nướng ạ, peko.''

Ryujin nhếch mép cười, nụ cười của một người đã quá quen với những màn náo loạn của đám nhóc này. Ông đáp với giọng điệu đầy tự nhiên: ``Trùng hợp ghê, nhà chú cũng đang định ăn lẩu bò đây! Mấy đứa có thể làm chung cho vui.''

Cả nhóm đồng loạt reo lên đầy phấn khích.

Ngay lập tức, Sora và A-chan đảm nhận vai trò phân công nhiệm vụ. Những ai có kinh nghiệm nấu nướng được xếp vào nhóm bếp núc, trong khi những người còn lại lo phần trang trí. Khi đến lượt Ayame và Haato xung phong tham gia bếp, A-chan và Sora lập tức từ chối không thương tiếc, vì lý do gì thì ai cũng rõ.

``Không được đâu, Nakiri-san. Cô quên mất một điều mà cô không được phép làm trong công ty sao?''

``Còn Haato-chan thì\dots\ bữa ăn của em quá `đặc biệt' đến mức không ai dám thử.''

Nàng Song tính phồng má giận dỗi, trong khi nàng Quỷ sừng chỉ cười trừ.

Sau khi nhóm nấu ăn lên bếp cùng Ryujin và Kaien (vừa trở về sau một chuyến giao hàng kiêm đi chợ), nhóm còn lại bắt tay vào trang trí phòng tiệc. Theo lời của người đàn ông, tầng ba là nơi rộng nhất có thể chứa hơn bốn mươi người, đó cũng chính là phòng mà mọi người tập trung vào dịp Tết năm ngoái, nên đó sẽ là địa điểm lý tưởng nhất cho bữa tiệc sinh nhật.

Trước khi ai nấy bắt đầu công việc, cả nhóm tập trung lại để nghe lời động viên từ ``trưởng ban'' A-chan và Sora. Ryujin, dù không có trong kế hoạch ban đầu, cũng bị lôi kéo vào buổi khởi động tinh thần này.

Sora hít một hơi thật sâu, rồi hào hứng hô lớn: ``Mọi người ơi! Hôm nay là ngày gì?''

Tất cả đồng thanh đáp: ``\textbf{Ngày 17 tháng 7!}''

Cô nàng tiếp tục: ``Vậy hôm nay là ngày gì nào!?''

Mọi người càng phấn khích hơn, đồng loạt reo lên: ``\textbf{Là sinh nhật của anh Tổng chúng ta!!}''

Và rồi, cô giơ nắm tay lên đầy quyết tâm: ``Vậy chúng ta hãy cùng nhau làm nên một bữa tiệc sinh nhật đáng nhớ cho anh ấy nào!''

Tất cả nhất loạt đồng ý, bầu không khí trở nên sôi động hơn bao giờ hết. Ryujin nhìn khung cảnh đó cũng chỉ bật cười: ``\vie{\textit{Làm mình liên tưởng đến hồi thanh niên sôi nổi ngày xưa.}}''

Và như thế, bước thứ tư trong kế hoạch ``Sinh nhật bất ngờ'' của \hll\ chính thức bắt đầu.

\ct

Để đảm bảo bí mật tuyệt đối, nhóm ba mươi hai cô gái sẽ âm thầm di chuyển lên tầng hai bằng cầu thang bên phải. Trong khi đó, Ryujin sẽ lên bằng cầu thang bên trái, là lối dẫn trực tiếp đến phòng của Kuon. Ông sẽ giữ phong thái bình thường, tránh gây bất cứ nghi ngờ nào tới cậu con trai của mình.

Sau khi tách khỏi nhóm, người đàn ông chậm rãi bước lên cầu thang. Ông không quên cúi chào mẹ mình - bà nội của Kuon - rồi ngay lập tức lên phòng.

Khi bước vào, ông thấy cảnh tượng quen thuộc: cậu quý tử của ông vẫn cặm cụi bên máy tính, chăm chú vào công việc như mọi ngày, dường như không để tâm tới những sự tình xung quanh cậu, còn Ryuuhou thì đang nằm dài trên đệm, mắt dán chặt vào màn hình ti-vi.

Nghe tiếng bước chân, Kuon chỉ ngẩng đầu lên một chút rồi hỏi: ``\vie{Bố có việc gì mà lên lâu thế?}''

Ryujin đáp lại bằng một giọng điệu thản nhiên, giả vờ không biết gì: ``\vie{Bố vừa mới xuống nói chuyện với bác Hijaki.}''

Cậu Tổng nghe vậy chỉ gật đầu nhẹ, mắt dán vào máy tính: ``\vie{Ra thế. Chắc cũng chỉ mấy chuyện việc vặt thôi nhỉ?}''

Người đàn ông gật đầu. Ông liếc mắt quanh phòng một lượt, rồi hỏi tiếp, cố giữ phong thái bình thường: ``\vie{Đã cất đồ lên trên gác chưa?}''

``\vie{Rồi bố. Con để sẵn trong tủ lạnh rồi ạ.}''

Như thường lệ, ông lại quay sang hỏi thăm cô em gái của cậu Tổng bằng vài câu xã giao: ``\vie{Trưa nay ăn hết cơm chưa?}''

Ryuuhou đáp mà không rời mắt khỏi màn hình: ``\vie{Ăn hết rồi ạ.}''

``\vie{Làm bài xong chưa?}''

``\vie{Rồi ạ.}''

Nghe vậy, ông Ryujin chỉ gật gù, rồi nhanh chóng rời khỏi phòng, tiến thẳng lên bếp để hỗ trợ việc chuẩn bị bữa tối với các cô nàng \hll.

Trong phòng khách, cô em gái vẫn đang mải mê xem ti-vi, trong khi anh trai của cô tiếp tục gõ bàn phím lách cách, hoàn toàn không tỏ vẻ nghi ngờ gì.

Nhưng rồi, giọng của cô bất chợt vang lên: ``\vie{Onii-san, nhà mình tính ăn gì tối nay vậy?}''

``\vie{Lẩu bò,}'' cậu đáp gọn, mắt vẫn liếc về đống tài liệu.

``\vie{Chỉ nhà mình thôi sao?}''

``\vie{Ừ, chắc vậy. Chỉ bốn người nhà mình thôi. Cùng lắm rủ thêm bác Hijaki-san với bà nội lên ăn chung,}'' cậu dừng lại một chút, rồi lẩm bẩm: `Anh thì ăn ít người hay nhiều người cũng chẳng quan trọng, miễn vui là được.''

Dứt lời, cậu tiếp tục công việc, vừa gõ bàn phím, vừa tranh thủ lướt tin tức trên báo điện tử trên một trang tính khác.

Bỗng một tiêu đề giật gân thu hút sự chú ý của cậu: ``\emph{Bốn chiếc xe điên gây náo loạn thành phố!}''

Tò mò, Kuon nhấp vào bài báo. Một đoạn văn chi tiết hiện ra trước mắt cậu:

\txtbx{Ngày 17/7/2021, một ngày bất thường đã khiến thủ đô Kawachi rơi vào hỗn loạn khi bốn chiếc xe với hành vi kỳ lạ đồng loạt xuất hiện trên đường phố, gây ra hàng loạt sự cố giao thông, phá hủy hàng loạt hạ tầng và khiến người dân không khỏi kinh ngạc.

  \

  Theo các nhân chứng, cả bốn chiếc xe được ghi nhận đã xuất phát từ một điểm chung, sau đó di chuyển theo các hướng khác nhau. Mỗi chiếc xe mang đặc điểm và hành động riêng, tạo nên những cảnh tượng khó tin giữa lòng thành phố.

  \

  Chiếc xe đầu tiên nổi bật với khả năng biến thành robot hai chân, nhảy qua cây cầu Eisui. Anh S.T., một người đi đường chứng kiến cho hay: ``Tôi cứ nghĩ mình đang xem phim viễn tưởng, không thể tin nổi cảnh đó là thật.''

  \

  Chiếc xe thứ hai gây sốc khi lao xuống hồ Kanken, nổi trên mặt nước rồi bay lên chạy trên đường ray tàu hỏa, làn náo loạn khu khố cổ. Bà H.K., cư dân gần đó chia sẻ: ``Tôi chưa từng thấy thứ gì kỳ lạ như vậy, như thể xe biết bay!''

  \

  Chiếc xe thứ ba, một chiếc bán tải quân sự mà chúng tôi chưa xác định được đã đối đầu với lực lượng an ninh, gây ra vụ nổ lớn làm sập một cây cầu quan trọng. Bà Y.R., kể lại: ``Tiếng nổ lớn đến mức rung chuyển cả khu phố, thật đáng sợ.''

  \

  Chiếc xe cuối cùng, mang dáng vẻ xe cảnh sát, tham gia rượt đuổi một nhóm tội phạm, dẫn đến nhiều vụ va chạm và vụ nổ trên cao tốc. Ông T.J., tài xế vô tình gây tai nạn chia sẻ, giọng run rẩy: ``Tôi chỉ muốn yên ổn làm việc, ai ngờ lại gặp cảnh phim hành động ngoài đời!''

  \

  Sự việc đã gây tắc nghẽn giao thông nghiêm trọng, làm hư hỏng nhiều phương tiện và công trình công cộng. Lực lượng chức năng đã nhanh chóng vào cuộc để kiểm soát tình hình, nhưng người dân vẫn chưa hết bàng hoàng trước những gì xảy ra. Một nhân chứng đã trả lời với phóng viên: ``Tôi đã hơn 70 tuổi mà tôi chưa gặp trường hợp hỗn loạn nào như vậy giữa lòng Thủ đô cả. Không lẽ đó là người ngoài hành tinh!?''

  \

  Theo báo cáo của Đội Cảnh sát Trung ương Thủ đô, đã có ít nhất 32 người thiệt mạng, trong đó có 12 chiến sĩ công an đã hy sinh sau vụ nổ ở cây cầu Shuouyou, 20 người còn lại đã thiệt mạng sau những vụ nổ trên cầu cao tốc trên cao thuộc vành đai ngoài của thủ đô, tất cả đều là những kẻ nằm trong một tổ chức xã hội đen mà cả nước đang truy nã.

  \

  Tại nút giao gần siêu thị GO Dragon, một số lượng lớn người dân cũng đã tiếp tay xúm vào đánh hội đồng một nhóm người được cho là cầm đầu của băng đảng xã hội đen khét tiếng. Lực lượng công an đã phải rất vất vả thuyết phục người dân, họ mới chịu buông tha cho nhóm người và giao cho phía công an nhằm phục vụ công tác điều tra.

  \

  Hiện các cơ quan chức năng đang điều tra nguyên nhân và trách nhiệm liên quan. Trong khi đó, sự kiện này chắc chắn sẽ còn là đề tài bàn tán của người dân Kawachi trong thời gian dài.}

Lướt qua những bức ảnh đi kèm và những dòng bài báo, ánh mắt Kuon chợt sững lại vài giây.

Trong những bức ảnh mà bài báo đưa ra, cậu nhìn thấy bốn chiếc xe: một chiếc xe màu đen viền đỏ mang dáng dấp của Mio, một chiếc xe hình tên lửa giống với chiếc của Okayu, một chiếc xe bán tải quân sự, và một chiếc xe cảnh sát nhìn khá tương đồng với chiếc của Fubuki.

``Sao bốn chiếc xe này\dots\ trông quen thuộc ấy nhỉ?''

Cậu cau mày, cố nhớ xem đã từng thấy chúng ở đâu. Nhưng chỉ vài giây sau, Kuon lắc đầu, quyết định gạt bỏ suy nghĩ đó. ``Thôi kệ đi. Cứ tập trung vào công việc. Nghĩ làm gì cho nhức đầu.''

Cậu quay trở lại với bản báo cáo của công ty, để mặc những dấu hiệu đáng ngờ trôi vào quên lãng.

Không lâu sau, tiếng mở cửa vang lên từ phía dưới nhà. Kaien đã về, tay xách đầy nguyên liệu cho bữa tối, bao gồm cả những bìa đậu để ăn kèm với lẩu bò mà do ``quý tử'' của bà yêu cầu.

Vừa bước vào, bà cất tiếng hỏi: ``\vie{Bố con đâu rồi con?}''

Kuon chỉ đáp lại ngắn gọn mà không quay đầu, ngón tay chỉ lên trên: ``\vie{Ở trên bếp ạ.}''

Nghe vậy, bà Kaien gật đầu rồi nhanh chóng đi lên chuẩn bị bữa tối sớm cùng chồng, cẩn thận đóng cửa trượt để tránh gió điều hòa bay mất.

Mọi thứ vẫn diễn ra bình thường, theo góc nhìn của cậu Tổng. Cậu hoàn toàn không hay biết rằng, ngay lúc này, một nhóm ba mươi hai cô gái đang âm thầm bận rộn ở tầng sáu, chuẩn bị cho một kế hoạch mà cậu không hề ngờ tới\dots

\ct

Kaien vừa bước chân lên đến bếp thì đã thấy chồng mình cùng năm cô gái trẻ đang bận rộn chuẩn bị nguyên liệu. Bà khựng lại đôi chút, ánh mắt lướt qua từng gương mặt quen quen. Sau một lúc quan sát, bà nhận ra bốn trong số họ.

``Ơ kìa, chẳng phải là Mel-chan, Okayu-chan, Choco-san với Shion-chan sao?''

Cả năm người họ đồng hoạt hướng mặt về phía bà, đồng thanh lễ phép: ``Chúng cháu chào cô ạ!''

Người nội trợ mỉm cười đáp lại, rồi ánh mắt bà dừng lại ở cô gái nhỏ nhắn, mái tóc hồng nhạt, trên đầu đội một chiếc vương miện, có phần quen thuộc nhưng nhất thời bà chưa nhớ ra.

``Còn cháu là ai? Trông cô thấy quen lắm\dots''

Chưa kịp để bà đoán, chồng bà đã lên tiếng giải thích: ``À, đó là Luna-chan. Con bé cũng tới giúp tổ chức tiệc sinh nhật cho thằng Kuon nhà mình.''

% $\forall X \in transaction, buys(X, item_1) \wedge buys(X, item_2) \Rightarrow buys(X, item_3), [s, c]$

Luna lập tức nhoẻn miệng cười, đôi mắt long lanh ngước lên và nói bằng giọng non nớt đặc trưng: ``Cháu chào cô nạ-nora\~{}''

Nghe câu chào ngọng nghịu ấy, Kaien khẽ bật cười, một chút trìu mến xen lẫn ngạc nhiên hiện rõ trên gương mặt bà.

Bà lên tiếng, tỏ vẻ hơi ngạc nhiên: ``À, cô nhớ rồi! Cháu với Sui-chan từng sang nhà cô ăn mì ramen dịp lễ Vu Lan, đúng không?''

Luna hớn hở gật đầu lia lịa: ``Dạ đứn ạ-nora!''

Kaien lại đảo mắt nhìn quanh, rồi hỏi tiếp: ``Thế hôm nay mấy đứa tính nấu món gì vậy?''

Cả năm cô gái đồng thanh, không hẹn mà cùng reo lên như đã chuẩn bị từ trước: ``Nướng ạ!''

Bà ngạc nhiên, nhướng mày nhìn Ryujin và năm cô gái, thắc mắc: ``Ơ, chẳng phải nhà mình định ăn lẩu bò hay sao?''

Mel cười khúc khích, liền giải thích: ``Bọn cháu cũng có làm lẩu nữa ạ! Nhưng nhóm lẩu thì đang chuẩn bị ở dưới tầng ba rồi. Bọn cháu là nhóm phụ trách nướng.''

Nghe nhắc đến tầng ba, Kaien lập tức hiểu ra vấn đề. Bà nhìn sang Ryujin, nửa cười nửa trách yêu: ``\vie{Thế mà em không hề nghĩ rằng họ sẽ đến đấy!}''

Người thợ cơ khí gật đầu, giọng nhẹ nhàng nhưng thoáng chút tự hào: ``\vie{Ừ. Nhóm này chia ra làm hai: một nhóm nấu ăn, một nhóm trang trí. Còn tầng ba thì tụi nó chọn làm nơi tổ chức tiệc.}''

Kaien khẽ gật đầu, ánh mắt lấp lánh ý cười. Bà không hề ngạc nhiên hay thắc mắc thêm nữa, dường như tất cả đã trở nên hợp lý trong lòng bà. Bà khẽ nói: ``\vie{Vậy mà anh không nói sớm, làm em cứ tưởng có khách bất ngờ!}''

Ryujin nhún vai, cười thoải mái: ``\vie{A-chan sắp xếp cả đấy. Bọn này vừa mới đến thôi, em xuống tầng ba phụ chúng chuẩn bị lẩu đi, bên phòng của anh Hijaki.}''

``\vie{Aloe-chan với Coco-chan chắc là không về đâu nhỉ?}''

``\vie{Chúng nó có việc của chúng nó, kệ đi. Chắc hôm nay không về đâu.}''

Kaien cầm lấy túi nguyên liệu vừa mang về, thoáng lắc đầu hơi thất vọng, nhưng trên môi vẫn giữ nụ cười: ``\vie{Thế là bày vẽ ra được hẳn một kế hoạch bất ngờ cơ đấy\dots\ Mấy đứa này thật biết cách làm trò.}''

Ryujin vẫy tay, không quên dặn thêm: ``\vie{Đi từ cầu thang bên kia, đừng để thằng Kuon nó nghi ngờ.}''

Kaien gật đầu, rồi xách đồ quay xuống tầng ba như lời dặn, để lại bầu không khí nhộn nhịp trên bếp.

Sau khi bà rời đi, Ryujin chống tay lên bàn, đảo mắt nhìn năm cô gái đang háo hức chờ đợi, rồi cười hỏi: ``Thế, mấy đứa đã sẵn sàng rồi chứ?''

Không ai bảo ai, cả năm người đồng thanh đáp lại, tràn đầy khí thế: ``\textbf{Rồi ạ, (nanora)!}''

Và thế là, những con người tại nhà gia đình Kuon - ba mươi hai cô gái trẻ cùng ông bố cơ khí và người đàn bà chủ kho kiêm nội trợ - bắt đầu công việc nấu nướng, chuẩn bị cho một bữa tiệc sinh nhật không giống bất kỳ bữa tiệc nào khác.

Những tiếng cười râm ran vang lên trong những căn bếp nhỏ, báo hiệu cho một buổi chuẩn bị đầy ắp sự háo hức và niềm vui vừa mới bắt đầu.

\il

\subsubsection*{\phantomsection\label{P4.1a}}
\addcontentsline{toc}{subsubsection}{Nhóm 1 - Nhóm làm đồ nướng}

\begin{center}
  {\color{T1} Nhóm 1 - Ryujin, Luna, Okayu, Shion, Choco, Mel\\Nhóm làm đồ nướng}
\end{center}

Trong căn bếp sáng ánh đèn, không khí trở nên rộn ràng một cách lặng lẽ. Nhóm đầu tiên - nhóm chuẩn bị món nướng - đã sẵn sàng bắt tay vào việc.

Nhưng điều mà người đàn ông của gia đình Hikawa ngạc nhiên là số lượng thịt bò và thịt lợn mà năm cô gái chuẩn bị. Nó nhiều đến mức khi chất đống chúng lên còn chạm cả trần nhà.

Ryujin trố mắt nhìn đống thịt khổng lồ, không giấu nổi sự kinh ngạc tới ngỡ ngàng: ``Nhiều thế này, mấy đứa định mở tiệc cho cả khu phố à?''

Mel cười khúc khích, giải thíc cho ông: ``Không có đâu chú ơi, chỗ này là để phục vụ cho ba mươi hai người lận đó ạ!''

Shion nhoẻn miệng cười gian: ``Bọn cháu cũng không phải tới đây chỉ để chuẩn bị rồi đi về đâu chứ!''

Okayu gật đầu, thêm vào: ``Mỗi người ăn một ít thôi mà, meo\~{}.''

Luna ngây thơ nghiêng đầu, lí nhí nói: ``Cháu nhĩ thế nì là vừa đủ nạ-nora\~{}.''

Nghe vậy, Ryujin chỉ có thể thở dài, mỉm cười trừ: ``Thôi được rồi, chú chịu mấy đứa luôn.''

Ryujin đứng cạnh chiếc bàn gỗ quen thuộc, nhẹ nhàng mở từng lớp màng bọc thực phẩm, để lộ ra hai bát lớn - một bát đầy ắp thịt bò cắt lát dày, một bát khác là thịt lợn tươi rói, mỡ nạc cân đối.

Ông đưa tay nhấc bát thịt bò lên, nhìn lướt qua năm cô gái trẻ đang háo hức đứng vây quanh. ``Trước tiên, chúng ta sẽ thái và ướp thịt đã,'' ông nói, giọng đều đều nhưng thân thiện.

Rồi, như một người thầy trong căn bếp nhỏ, ông hỏi: ``Bình thường các cháu ăn thịt nướng kiểu gì?''

Năm cô gái đồng thanh, không chút do dự: ``Yakiniku ạ-(nora)!''

Nghe câu trả lời, Ryujin khẽ cười: ``Chuẩn kiểu Tokyo rồi đấy nhỉ. Chú cũng từng nghe nó rồi, nhưng chưa bao giờ ăn. Chắc là nó ngon lắm chứ nhể?''

Cả nhóm cũng bật cười theo lời nói tếu táo của ông, không khí thoải mái hơn hẳn.

Nhưng ngay sau đó, ông đặt bát thịt xuống bàn, ánh mắt hơi ánh lên vẻ trêu chọc: ``Nhưng hôm nay, chú muốn giới thiệu cho các cháu một kiểu nướng khác một chút. Không chắc nó hợp khẩu vị với mấy đứa, nhưng chú cá là nó sẽ khiến các cháu đổi vị đấy.''

``Ồ\dots?'' cả nhóm đồng loạt nghiêng đầu, tò mò. Shion là người lên tiếng đầu tiên: ``Kiểu nướng đó là gì vậy?''

Người đàn ông nhìn cô, khóe miệng nhếch lên: ``Thịt áp chảo. Món đó nhà chú hay làm mỗi khi rảnh.''

Năm cô gái đưa mắt nhìn nhau, gương mặt thoáng vẻ ngạc nhiên. Mel nhíu mày: ``Áp chảo\dots\ nghe lạ nhỉ. Có khác gì thịt nướng không ạ?''

Choco chống cằm, đoán thử: ``Nghe thì giống thịt nướng bên mình, nhưng chắc cách làm có điểm khác biệt. Hình như là ta dùng chảo thay vì dùng vỉ nướng chăng?''

Luna nghiêng đầu, đôi mắt long lanh: ``Nướng với thịt nợn thì đúng nạ đó-nora\~{}''

Okayu cầm cằm, trầm ngâm: ``Hay đây là cách ăn đặc biệt của nhà anh Kuon-senpai chăng?''

Nghe những cô gái bàn luận, người thì đoán mò, người thì suy luận, Ryujin cười khẽ, đôi tay khoanh lại, nét mặt bình thản. Ông cần mẫn đáp: ``Đó chính là điểm khác biệt so với kiểu nướng Yakiniku của các cháu, ngoài việc có thể nướng vỉ ra, người ta cũng có thể dùng chảo để nướng nữa - gọi là rán. Ở nơi đây, cái gì có thể nướng được thì đều cho lên vỉ cả, không cứ phải là thịt bò. Thịt lợn cũng vậy thôi, quan trọng là cách tẩm ướp.''

Năm cô gái gật gù, ánh mắt càng thêm háo hức.

Ông gật đầu, rồi chậm rãi mở ngăn kéo, rút ra một con dao sắc. ``Trước tiên, phải cắt những miếng vừa ăn đã. Chú nghĩ, chắc các cháu biết dùng dao rồi nhỉ?''

Nàng Phù thủy nhanh nhảu chỉ sang hai người bên cạnh: ``Có Choco-sensei với Mel-senpai giỏi khoản này lắm ạ!''

Ryujin liếc mắt sang, hỏi tiếp: ``Thế ba đứa còn lại thì sao?''

Nàng Mèo đáp với giọng tỉnh bơ quen thuộc: ``Cháu cũng biết dùng dao, chú à\~{}''

Bản thân nàng Phù thủy cũng chỉ cười gượng: ``Cháu biết sơ sơ thôi, chắc không khéo bằng hai người kia.''

Nàng Công chúa xứ Kẹo lắc đầu, thành thật thú nhận: ``Cháu mới biết xơ thôi ạ-nora\dots''

Người đàn ông nghe cả năm người trả lời, gật đầu ra chiều hiểu ý: ``Vậy thì để Choco với Mel thái thịt trước làm mẫu nhé. Xong rồi, ba đứa còn lại sẽ thay phiên nhau làm cho quen tay.''

Ông dứt lời, nhẹ nhàng đẩy bát thịt bò về phía hai cô gái. Hai cô nàng - một người đã từng làm món cà ri\footnote{Xem lại \flaregothic{SPÉCIAL} {\ab XI}, quyển số Năm}, một người thì làm món trứng rán\footnote{Xem lại \flaregothic{ÉPISODE 76}, quyển số Sáu} - không cần đợi nhắc lần hai, đã cầm dao lên, cẩn thận cắt từng lát thịt đều tăm tắp, đúng chuẩn phong cách bếp núc mà họ đã quen thuộc.

Mel vừa cắt vừa quay sang mỉm cười: ``Mấy đứa cũng thử thái đi cho quen tay nhé.''

Choco tiếp lời, giọng vừa dứt khoát vừa thân thiện: ``Làm nhiều thì tay mới cứng, mấy đứa à. Không khéo đến sinh nhật lần sau lại nấu được cho Kuon-sama ăn nữa ấy chứ.''

\ct

Sau khi Mel và Choco hoàn thành phần của mình một cách gọn gàng và nhanh lẹ, dao thớt được chuyển sang tay của ba cô gái còn lại. Lần này đến lượt Okayu, nàng Mèo tím vốn cũng khá biết cơ bản với chuyện bếp núc, nhưng kỹ năng thì vẫn chưa thể sánh với hai người tiền bối.

Nàng Mèo tím cẩn thận cầm dao, hít một hơi rồi bắt đầu cắt từng lát thịt. Những lát đầu tiên trông khá đều và ngay ngắn, nhưng chỉ một lúc sau, do thiếu kinh nghiệm, có vài miếng dày mỏng thất thường, méo mó chỗ này, lẹm khuyết chỗ kia.

Cô nhún vai cười nhẹ nhìn thành quả của mình, có vẻ cũng không lấy làm phiền lòng. ``Chắc em phải luyện thêm nhiều mới thành thạo được,'' cô vừa cắt vừa nói.

Tiếp theo là lượt của Shion. Khác với ba con người ban nãy, nàng Phù thủy tóc xám chẳng buồn nhấc dao lên.

Thay vào đó, cô lặng lẽ đưa tay ra niệm một câu chú, khiến tảng thịt trước mặt bỗng tách ra thành nhiều miếng đều tăm tắp, gọn gàng một cách kỳ lạ.

Ngay khi phép thuật hoàn thành, Luna cùng Okayu lập tức vỗ tay bôm bốp, mắt tròn xoe vì thích thú. Nhưng hai người còn lại - Mel và Choco - thì chỉ khoanh tay, liếc nhìn Shion với ánh mắt nửa trêu nửa không phục. Ryujin thì chỉ có thể nhìn ``tác phẩm'' cô nàng vừa hoàn thành mà không biết phải bình luận như thế nào.

Nàng Y tá chau mày: ``Ý của cô là dùng dao cơ mà! Đừng có chơi phép thuật chứ!''

Nàng Phù thủy bĩu môi, chống nạnh phân bua: ``Thì mọi người bảo thái ra mà, có ai nói cấm dùng phép đâu!''

Nàng Dơi cười khúc khích, thêm vào: ``Làm thế là gian lận đấy, Shion-chan.''

Ryujin đứng cạnh, cũng bật cười, xoa cằm nói đùa: ``Thôi được rồi, chú cứ cho là thế này cũng tính thái đi.''

Cuối cùng, con dao được trao cho Luna, cô ``Công chúa xứ Kẹo'' nhỏ nhắn nhất nhóm. Có lẽ do quen có người khác nấu nướng cho mình, cô nàng khi lần đầu cầm dao đã tỏ ra rất vụng về, hai tay nắm chặt, ánh mắt lúng túng.

Cô cẩn thận đưa dao xuống, nhưng đường cắt đầu tiên lại lệch hẳn sang một bên, làm cho hầu như không cắt được tí nào. Cô hậu đậu cố gắng cắt thành những thớ tiếp theo, nhưng rồi miếng thịt mà cô xắt ra thì chỗ dày, chỗ mỏng, có miếng nhỏ xíu, có miếng lại to oành.

Thậm chí, trong một lần lơ đễnh, cô suýt nữa thì cắt trúng ngón tay, nếu không nhờ người thợ cơ khí và nàng Y tá hô lớn.

Một sự im lặng toát lên cả căn bếp khi cả bốn người nhìn cô thái thịt bằng ánh mắt ái ngại. Ryujin cuối cùng cũng không nhịn được, bật cười hiền hậu: ``Có vẻ như cháu vẫn chưa quen cầm dao đúng không?''

Luna xấu hổ, gãi má: ``Dạ\dots\ Hông ạ-nora\dots''

Ông gật đầu: ``Không sao, để chú chỉ cho.'' Dứt câu, ông tiến lại gần, kiên nhẫn cầm tay cô bé, hướng dẫn đầy tỉ mỉ.

``Quan trọng nhất là phải bảo vệ ngón tay. Nhìn nhé: khum bốn ngón tay lại, ngón cái giấu vào phía sau, rồi để bốn ngón kia tì sát vào lưỡi dao. Cứ từ từ, cắt từng miếng một.''

Luna chăm chú quan sát, đôi mắt long lanh, giọng lí nhí: ``D-Dạ\dots\ zâng.''

Dưới sự hướng dẫn tận tình của Ryujin, từng lát thịt được nàng Công chúa dần dần được cắt ra, tuy chưa thật hoàn hảo nhưng rõ ràng là an toàn và ngay ngắn hơn lúc trước.

Những người đứng quanh lặng lẽ quan sát, và không ai bảo ai, Mel khẽ bật cười, thì thầm với Choco: ``\hl{Nhìn họ như thể hai bố con ghê.}''

Đồng đội gật đầu, mỉm cười, ánh mắt toát lên vẻ ấm áp. Cô không nói một lời gì, chỉ im lặng nhìn khung cảnh trước mặt.

Okayu cũng chăm chú quan sát mọi thứ, miệng lẩm bẩm: ``\hl{Mình cũng muốn học chút ít từ ông ấy để nấu ăn cho Koro-san nữa\dots}''

Khi Luna hoàn thành xong phần của mình, Ryujin lại lên tiếng, lần này quay sang nhắc nhở cả nhóm: ``Thịt lợn các cháu cũng làm y hệt như thế nhé, nhớ cắt vừa miếng thôi.''

Cả năm cô gái đồng thanh, giọng đồng đều, vui vẻ: ``Vâng ạ\~{}!''

Tiếng dao cắt thịt lại vang lên đều đặn, mỗi người đều tập trung vào phần việc của mình. Không khí trong bếp trở nên nhộn nhịp nhưng cũng đầy sự phối hợp ăn ý. Những lát thịt bò và thịt lợn được cắt đều đặn, xếp ngay ngắn trên đĩa, sẵn sàng cho công đoạn tiếp theo.

Shion liền nghiêng đầu hỏi, giọng hồ hởi: ``Sau khi thái xong rồi, bước tiếp theo là gì ạ?''

Choco-sensei - người tạm giữ vai trò ``phó bếp'' - liền đáp ngay: ``Cô nghĩ chắc chúng ta sẽ đem nó đi tẩm ướp chúng lên thôi.''

Người thợ cơ khí gật đầu: ``Choco-chan nói đúng rồi đấy. Muốn có một miếng thịt ngon và dậy mùi, ta phải tẩm ướp nó lên.''

Nàng Phù thủy tò mò tiếp tục: ``Thế chúng ta sẽ ướp bằng gì?''

Okayu, lúc này vẫn đang lau tay, nhanh nhảu chen vào: ``Thường thì là sốt nướng hoặc nước tương thôi ạ, meo\~{}.''

Ryujin nhếch môi cười, ánh mắt ánh lên vẻ trêu chọc: ``ách đó cũng hay. Nhưng ở đây, nhà chúlại hay làm kiểu vừa đậm chất Kawachi, lại mang dáng dấp nhà Hikawa nữa.''

Choco tò mò hỏi: ``Kiểu Hikawa\dots\ là thế nào vậy, Ryujin-sama?''

Ông lấy từ trong chạn ra một túi nhỏ đựng sả, một ít riềng và mấy củ hành khô, đặt xuống bàn: ``Chỉ đơn giản: tiêu, sả, hành khô, tỏi và riềng thôi. Không cầu kỳ sốt, cũng chẳng cần nước tương. Cùng lắm là một ít dầu hào vào để tôn hương vị của nó lên là đủ.''

Mel nghiêng đầu, giọng có chút ngạc nhiên: ``Nghe lạ ghê\dots\ khác hẳn bọn cháu hay ăn.''

Luna tròn mắt, gãi gãi má, ngây ngô hỏi: ``Diềng\dots\ Xả\dots\ là gì vậy-nora?''

Nàng Y tá cười nhẹ, giải thích bằng giọng kiên nhẫn: ``Đó là những loại gia vị thơm cho món ăn đó, Luna-sama. Ở Tokyo thì hơi khó tìm những loại như thế này.''

Mel tiếp lời, khẽ gật đầu: ``Nếu có bán thì giá của chúng cũng đắt lắm á.''

Không để mọi người phải đợi lâu, Ryujin mang hai chiếc bát lớn chứa đầy thịt bò và thịt lợn đã được thái gọn gàng ra, bắt đầu hướng dẫn từng bước: ``Ta sẽ rắc một ít tiêu, xếp sả, hành khô và riềng lên trên. Nhưng lưu ý: tỷ lệ ướp giữa thịt bò và thịt lợn không giống nhau đâu nhé.''

Shion hồn nhiên nảy ra một ý tưởng nghe có vẻ điên rồ: ``Vậy\dots\ mình có thể trộn hai loại thịt này lại luôn không ạ?''

Mel lập tức nhắc khẽ, hơi lắc đầu: ``Không được đâu. Chú ấy vừa nói rồi đấy, tỷ lệ ướp khác nhau mà.''

Luna xụ mặt, lẩm bẩm: ``Chán nhỉ-nora\dots\ Giá như có thể chộn vào thì nhanh hơn bít mấy.''

Ryujin kiên nhẫn giải thích cặn kẽ về tỉ lệ nêm nếm cho từng loại thịt, rồi tiếp tục chỉ dẫn: ``Sau khi rắc đủ gia vị, ta sẽ dùng tay trộn đều cho ngấm. Nhớ là trộn vừa phải thôi, đừng để gia vị dồn cục.''

Mel lễ phép chen vào hỏi: ``Bọn cháu đeo bao tay để trộn cho vệ sinh nhé chú?''

Người đàn ông cơ khí gật đầu: ``Ừ, lấy trong tủ chạn ấy. Chắc chú cũng dùng bao tay thôi không các cháu lại kêu.''

Cả năm cô gái đều phì cười với lời bình của ông.

Thế rồi, từng người một bắt đầu đeo bao tay, chuẩn bị cho công đoạn trộn thịt. Choco và Mel, với kinh nghiệm sẵn có, thực hiện một cách khéo léo, từng miếng thịt được thấm đều gia vị, không quá nhạt cũng không quá đậm. Những tiếng cười nói râm ran vang lên, hòa quyện cùng âm thanh nhào nặn của từng thớ thịt, tạo nên một bầu không khí ấm áp, gần gũi. Ai nấy đều tập trung, nhưng không quên trao đổi những câu chuyện vui vẻ, như thể đang cùng nhau chuẩn bị một bữa cơm gia đình đầy ý nghĩa.

Kể cả Okayu. Dù cho cô nàng trông khá ổn, cô có đôi phần lơi lơ đãng vì mải mê\dots\ ngửi mùi gia vị hơn là trộn, thậm chí còn lén bốc một miếng thịt sống lên ngắm nghía, định giơ lên trước mặt trêu chọc Luna.

``Luna-tan\~{} này, thử cắn sống không\~{}?''

Nàng Công chúa xứ Kẹo tròn mắt hoảng hốt, xua tay liên tục: ``Ẹee\~{}! Hổng dám đâu-nora\~{}!''

Trong khi cả nhóm đang trộn đều, Shion lại tinh ranh lặng lẽ giơ tay niệm phép nhân lúc không ai để ý. Ngay lập tức, toàn bộ các thớ thịt lợn và thịt bò đều được hòa lẫn với nhau, nhưng rồi lại đảo lộn lung tung lên. Shion cười gian, thì thầm: ``\hl{Như thế này có phải nhanh hơn không!}''

Tiếc rằng, hành động của cô không qua mắt được Mel. Nàng Dơi nhíu mày, phát hiện ra liền lắc đầu nhắc: ``Shion-chan, làm vậy không được đâu. Chú Ryujin-san vừa nói rồi đấy, thịt lợn với thịt bò không thể nào trộn lẫn với nhau được. Phải làm riêng ra!''

Nàng Phù thủy chu môi, thở dài: ``Ờ ha\dots\ Quen tay quá, quên mất. Cơ mà trộn từng bát thịt thế này thì lâu lắm á\~{}''

``Muốn ăn thì phải lăn vào bếp thôi, Shion-sama,'' nàng Y tá phản biện. ``Nấu ăn thì phải để cái tâm ở trên, chứ không phải cứ làm là xong đâu.''

Nghe vậy, Shion chỉ có thể rầu rĩ làm theo, tách từng mảng thịt lợn và thịt bò về như cũ và bắt đầu trộn từng bát một - tất nhiên là vẫn bằng ma thuật rồi, cô nàng chẳng muốn động chân động tay một chút nào.

Dẫu cho cả hai cô nàng thạo bếp chỉ có thể ngán ngẩm nhìn cô làm như thể bị ép buộc, còn Ryujin thì cười trừ, họ cũng không muốn làm cho Shion phát cáu thêm.

Cuối cùng là Luna. Dù đã được người bố của Kuon chỉ cẩn thận cách trộn, nhưng cô nàng có phần quá nhiệt tình, tay trộn mạnh đến mức suýt nữa làm bát thịt bật ra khỏi bàn.

Ryujin - với phản xạ của một người làm máy lâu năm - vừa kịp đỡ lấy, bật cười hiền hòa: ``Trộn vừa thôi, mạnh thế là vỡ bát bây giờ.''

Luna đỏ mặt, gãi đầu, giọng lí nhí: ``Zâng-nạ\dots\ tại cháu hông bít nạ-nora\~{}''

Ông lại dịu dàng nhắc lại, chậm rãi chỉ cho cpp cách xoay tròn bát, trộn đều tay mà không làm vỡ bát, vừa hướng dẫn vừa giữ tay cô bé, cho đến khi thịt được trộn xong một cách gọn gàng.

Một lúc sau, khi tất cả các bát thịt đã ướp xong, Ryujin gật gù hài lòng, quay sang cả nhóm: ``Thịt lợn các cháu cũng làm y như thế nhé. Quan trọng là đừng gấp, cứ thong thả thôi.''

Cả năm cô gái đồng thanh đáp, giọng vang lên như một lời hứa hẹn vui vẻ: ``Vâng ạ(-nora!)\~!''

\ct

Ryujin đặt bát thịt đã được trộn ngăn nắp lên bàn, xoa tay, hài lòng nói: ``Chúng ta cứ để thịt tự ướp thế này thôi, không cần động vào nữa.''

Luna tò mò ngó vào bát, nghiêng đầu hỏi bằng giọng non nớt: ``Ủa, thế nà xong rồi nạ\~{}?''

Choco nở nụ cười dịu dàng, nhẫn nại giải thích: ``Ừ, giờ chỉ cần để yên cho gia vị ngấm vào thịt thôi, không cần làm gì thêm đâu.''

Mel tiếp lời, giọng nhẹ tênh nhưng mang đầy kinh nghiệm: ``Thịt để càng lâu thì càng ngấm vị, nướng lên mới ngon được.''

Shion đứng một bên, khoanh tay lém lỉnh cười, rồi bỗng dưng đề xuất: ``Thế thì để tôi tăng tốc giai đoạn lên nhé!''

Ryujin ngạc nhiên nghiêng đầu, mắt tròn xoe: ``Cháu có thể đẩy nhanh cả quá trình ướp thịt được sao?''

Shion cười tự tin, nhún vai một cách đầy ma mị, như thể cô là một nữ phù thủy thực thụ: ``Cháu là phù thủy mà, chỉ cần một lần phép thôi là xong hết!''

Mel lắc đầu cười khúc khích, trêu nhẹ: ``Hiếm khi chị thấy phép thuật của Shion-chan lại dùng đúng chỗ như thế đấy.''

Okayu cũng gật gù, chống cằm: ``Công nhận\~{} Bình thường toàn thấy Shion-chan dùng phép để muốn được lười thôi.''

Nàng Phù thủy nhỏ nghe vậy liền phồng má, nổi cáu: ``Này này! Hai người nói vậy ý gì hả!?''

Không để ai có cơ hội ``phản công'' lại, cô nhanh tay tạo ra một cái vòm ánh sáng nhỏ bao quanh hai bát thịt, rồi thì thầm vài câu thần chú. Trong chớp mắt, bầu không khí trong vòm thay đổi, khiến những thớ thịt bắt đầu ``tua nhanh'' quá trình ngấm gia vị một cách ảo diệu.

``Đấy, thế này là chuẩn rồi. Giờ cứ để thế mà chờ thôi,'' cô nàng hếch cằm, đầy kiêu hãnh.

Ryujin nhìn bát thịt được ướp với tốc độ nhanh, miệng ông khẽ há ra đầy kinh ngạc, nhưng rồi cũng chỉ có thể cười hiền nhìn thành phẩm được đẩy nhanh công đoạn bằng một sức mạnh đặc biệt.

Trong lúc đợi, cả nhóm không tránh khỏi buông mấy câu chuyện phiếm, còn người đàn ông thì tranh thủ lướt báo trên điện thoại. Okayu chống cằm, ánh mắt lơ đãng nhìn bát thịt rồi buột miệng: ``Không biết Koro-san với mọi người dưới tầng kia trang trí tới đâu rồi ha\~{}''

Mel bật cười, nhìn Okayu với vẻ trêu chọc: ``Biết đâu em ấy lại biến phòng khách thành chiến trường cũng nên.''

Choco cười khúc khích: ``Nhưng dù là gì thì có lẽ Kuon-sama sẽ thấy cảm động đến phát khóc đấy.''

Luna thì ngồi lắc lắc chân, gật gù, mắt long lanh như trẻ nhỏ: ``Nếu có bỡi thựt thì zui lắm nora\~{}!''

Shion thì khoanh tay, bĩu môi: ``Không chừng xuống dưới lại thấy phòng khách bị treo ngược cả lên chứ chả đùa.''

Cả nhóm cười rộ lên, ai nấy đều tưởng tượng ra đủ thứ cảnh tượng kỳ quặc mà một thành viên ``chuyên đùa nhây'' của họ có thể bày ra.

Lúc ấy, Ryujin mỉm cười, nhìn đồng hồ trên điện thoại rồi nói: ``Được rồi, chắc cũng sắp xong. Ta xuống xem các đội kia làm thế nào rồi.''

Cả năm cô gái đồng thanh hô vang, giọng nói tràn đầy háo hức: ``Vâng ạ\~{}!''

Thế là cả nhóm theo chân ông, rời khỏi bếp, lần lượt bước xuống tầng hai bằng cầu thang bên phải - nơi mà nhóm trang trí đang âm thầm chuẩn bị một bất ngờ lớn dành cho cậu con trai của ông.

\subsubsection*{\phantomsection\label{P4.1b}}
\addcontentsline{toc}{subsubsection}{Nhóm 2 - Nhóm làm lẩu}

\begin{center}
  {\color{T2} Nhóm 2 - Kaien, Aki, Suisei, Rushia, Noel, Subaru\\Nhóm làm lẩu}
\end{center}

Cùng lúc đó, khi Ryujin đang hướng dẫn năm cô gái ở tầng trên chuẩn bị món nướng, Kaien cũng từ tốn xách giỏ nguyên liệu bước xuống tầng bốn, đi thẳng về phía phòng bếp bên phải - phòng của nhà Hijaki - như lời dặn của chồng bà.

Vừa mở cửa, bà đã thấy không khí trong bếp dưới này cũng náo nhiệt chẳng kém. Bên trong, ngoài những dụng cụ nấu nướng đã bày biện sẵn, có năm cô gái đang tụm lại, chuẩn bị nguyên liệu cho món lẩu. Đập ngay vào mắt của người phụ nữ nhà Hikawa là bốn gương mặt quen thuộc: Aki, Rushia, Subaru và Noel.

Nhưng xen lẫn giữa họ là một cô gái với mái tóc xanh nước biển, gương mặt có phần vừa quen thuộc vừa lạ lẫm đối với bà.

Bà nghiêng đầu hỏi một cách thân thiện: ``Cháu là\dots?''

Cô gái tóc xanh nước biển lễ phép cúi người, trả lời gọn gàng: ``Hoshimachi Suisei ạ!''

Nghe đến cái tên đó, Kaien nhíu mày suy nghĩ vài giây rồi bật cười: ``À, hình như cháu từng đến đây ăn mì ramen cùng với Luna-chan nhà cô  đúng không?''

Suisei gật đầu, nụ cười tươi rói hiện rõ trên gương mặt: ``Dạ đúng ạ! Nước dùng nhà cô nấu ngon lắm á!''

Kaien mỉm cười dịu dàng, ánh mắt đầy vẻ trìu mến: ``Cảm ơn cháu nha. Nhà cô lúc nào cũng có sẵn, nếu rảnh thì mấy đứa thích thì cứ qua, nhà cô nấu cho ăn bất cứ lúc nào.''

Cả năm cô gái gật đầu lia lịa, vẻ mặt ai nấy đều phấn khởi, rồi nhanh chóng bắt tay vào công đoạn đầu tiên của nhóm - chuẩn bị nước lẩu.

Nhưng khi Kaien vừa bước thêm vài bước vào trong phòng, ánh mắt bà lập tức bị hút chặt vào một\dots\ ``ngọn núi thịt'' chất cao ngất, gần như chạm cả lên trần nhà. Bà hơi sững lại, ngẩn ngơ nhìn đống thịt được sắp ngay ngắn trên bàn, dù rõ ràng số lượng cũng chẳng nhiều hơn so với đám thịt mà chồng bà đã chuẩn bị ở tầng trên.

Ánh mắt bà lướt qua từng gương mặt, vừa tò mò vừa buồn cười.

Noel và Suisei thì hớn hở ra mặt, cứ như thể chỉ chờ đúng khoảnh khắc này để cắm đầu vào ``bữa tiệc thịt'' của riêng mình. Hai người cứ ríu rít khoe nhau về việc tối nay ăn bao nhiêu phần, đến mức bà không nhịn được phải bật cười.

Aki đứng bên cạnh, tay chống hông, chỉ biết cười trừ. Có vẻ cô nàng cũng đã quen với kiểu ``tham ăn không lối thoát'' của hai người kia nên chỉ khẽ lắc đầu mà không nói gì.

Subaru thì khoanh tay, liếc nhìn ``núi thịt'' ấy bằng ánh mắt vừa ngán ngẩm vừa bất lực, miệng không quên lầm bầm: ``Thế này chắc thành lẩu thịt luôn mất\dots\ May mà chúng ta có rau củ ăn kèm.''

Rushia cũng chẳng khá hơn là bao, nhìn đống thịt cao sừng sững, mồ hôi tuôn ra không kịp, ngửa mặt than thở nhỏ: ``Đúng kiểu lẩu phiên bản\dots\ cho ba mươi người đây rồi-nanodesu. Đến em nhìn thấy tởn tới già luôn á.''

Kaien chỉ thở dài một tiếng thật khẽ, cuối cùng cũng đành lắc đầu cho qua. Dẫu gì, con trai bà cũng quen với cảnh này rồi, hôm nay có thêm vài đứa ăn khỏe nữa chắc cũng chẳng phải vấn đề lớn.

\ct

Trong lúc cả nhóm loay hoay chuẩn bị lửa, bầu không khí trong căn bếp tầng bốn vẫn tiếp tục rộn ràng, hệt như tầng trên.

Kaien vừa sắp xếp lại nguyên liệu mà mọi người chuẩn bị, cất giọng hỏi: ``Ryujin bảo các cháu xuống đây làm lẩu hả?''

Rushia và Aki đồng thanh đáp: ``Dạ, đúng ạ!''

Bà tò mò hỏi tiếp: ``Vậy, bình thường các cháu hay làm lẩu kiểu gì?''

Subaru cười trừ, nhớ lại một kỷ niệm không mấy dễ chịu, liền kể ngay: ``Hồi đầu năm bọn cháu có làm lẩu vịt ạ\dots\ Nhưng sợ lắm cô!''

Kaien thoáng ngạc nhiên, hỏi lại: ``Sao lại sợ?''

Nàng Vịt bĩu môi, dẩu môi ra than thở: ``Vì cháu suýt nữa bị mọi người lôi vào nồi làm\dots\ vịt lẩu\footnote{Xem lại {\flaregothic ÉPISODE 96}, quyển số Tám} ạ.''

Cả phòng lập tức nổ tung trong tiếng cười. Không khí trở nên rộn ràng đến mức, chính Subaru cũng không nén nổi mà lầm bầm than vãn: ``Mồ\~{}! Không vui chút nào đâu! Đừng cười!''

Sau khi có một tràng cười với tình huống oái oăm của nàng Tomboy, Kaien lại hỏi tiếp, giọng dịu dàng: ``Thế ngoài lẩu vịt, các cháu còn hay ăn lẩu gì nữa?''

Nàng Hiệp sĩ là người đầu tiên đáp, mắt sáng lên: ``Lẩu bò ạ!''

Nàng Pháp sư cười méo xẹo, trêu ngay: ``Bà chỉ thế là nhanh-nanodesu!''

Nàng ``Sao chổi'' chống nạnh, chen vào đầy phấn khích, đôi mắt cô sáng rực lên: ``Lẩu thịt! Thịt, thịt, càng nhiều càng tốt ạ!''

Rushia nhìn Suisei bằng ánh mắt bất lực: ``Lại còn cả chị nữa, Suisei-senpai\dots\ Thịt mãi thôi!''

Kaien bật cười, vỗ nhẹ tay vào thành bàn: ``Nghe mấy đứa nói toàn thấy mỗi thịt, chẳng ai nhắc tới rau củ gì cả nhỉ?''

Suisei lập tức xua hai tay đáp, mặt nghiêm nghị như thể tuyên bố một chân lý: ``Ở đây bọn cháu không làm thế.''

Ngay lập tức, bầu không khí có chút yên lặng khi nàng Ca sĩ toát ra một luồng sát khí\dots\ khiến ai nấy đều bất giác nuốt nước bọt. Duy chỉ có mẹ của Kuon là vẫn điềm nhiên, khẽ cười nhắc nhẹ: ``Con nhà cô cũng y hệt, chẳng chịu ăn rau bao giờ. Thế này nó với cháu là hợp nhau lắm rồi đấy, Hoshimachi-chan.''

Nghe vậy, khuôn mặt của Suisei lập tức đỏ bừng, tai cũng ửng hồng, luống cuống xua tay: ``K-Không có đâu cô ơi! Cô cứ hay đùa!''

Cô nàng quay ngoắt mặt đi, hai tay che mặt lại, trông như một quả cà chua chín mọng. Dù ngoài miệng thì chối đây đẩy, nhưng trong lòng cô vẫn không giấu nổi suy nghĩ lạ lẫm: ``K-Không lẽ thật sự hợp đôi sao\dots?''

Ý nghĩ ấy khiến đầu cô như muốn bốc khói, hai má càng đỏ hơn, còn mấy cô gái khác thì cứ thế cười nghiêng ngả, trêu ghẹo không ngừng.

``Ể\~{}? Sui-chan có bạn trai rồi sao?'' Subaru nhíu mày, nhe răng cười.

``K-Không có!''

``Thế tại sao mặt chị lại ửng lên thế?'' tới lượt Noel cũng góp vui.

``Nhìn mặt chị thế này lừa được ai nữa\~{}'' Aki cũng góp giọng.

``Đ-Đã bảo là không có! Suisei này là của tất cả mọi người! Đ-Đừng có mà ghẹo chị!'' nàng Ca sĩ lấy hai tay che khuôn mặt đang đỏ lên, còn ba cô gái còn lại thì cười lăn lộn.

Riêng chỉ có mỗi Rushia thì thở dài, bó tay với ba cô nàng thích đùa nhây này.

Kaien mỉm cười, nhẹ nhàng xoa dịu: ``Thôi nào, đừng trêu Hoshimachi-chan nữa.''

Cứ như thế, không khí ấm áp của nhóm cũng hòa quyện với tiếng cười rộn ràng từ tầng trên, khi hai nhóm nấu ăn âm thầm chuẩn bị cho một bữa tiệc sinh nhật bất ngờ - mà chủ nhân của bữa tiệc - người vẫn đang cắm đầu vào máy tính - thì hoàn toàn chẳng hay biết gì.

Suisei bất ngờ quay sang hỏi như muốn thoát khỏi chủ đề yêu đương, ánh mắt sáng long lanh đầy tò mò: ``Vậy\dots\ chúng ta sẽ nấu loại lẩu gì vậy, cô ơi?''

Kaien, vẫn đang lựa núi rau củ, ngẩng đầu lên cười dịu dàng: ``Hôm nay tụi mình sẽ cùng nhau làm bò nhúng dấm nha!''

Cả nhóm bốn cô gái thoáng nhìn nhau, ánh mắt xen lẫn chút ngạc nhiên và tò mò khi nghe cái tên món ăn có phần lạ lẫm.

Nàng Cô đồng nhíu mày, tay chống cằm: ``Món gì nghe lạ vậy-nanodesu\dots''

Nàng Tomboy xoa gáy, gật gù tán thành: ``Chị cũng chưa nghe thấy tên nó bao giờ luôn.''

Nữ Hiệp sĩ thì nghiêng đầu, ánh mắt đầy ngây thơ: ``Nghe có vẻ lạ thật\dots\ Ai giải thích cho em biết đi, chứ em mù tịt luôn này!''

Nàng Dị giới khẽ cười, suy đoán theo bản năng của một người có kinh nghiệm bếp núc: ``Chắc là món người ta nhúng thịt bò vào giấm, rồi\dots\ ăn luôn?''

Suisei nghe những hậu bối giải thích mà nhăn mặt, nửa tò mò nửa rụt rè: ``Nghe có vẻ\dots\ hơi kinh dị nha.''

Kaien bật cười, lắc đầu tỏ vẻ cảm thông: ``Aki-chan đoán gần đúng rồi đó. Đây là món lẩu bò, nhưng phần nước dùng sẽ có vị chua thanh nhờ nấu từ giấm và một ít hoa quả, chứ không giống nước lẩu thường.''

Rushia, Noel, Subaru và Aki đồng loạt ``ồ'' lên, gật gù như vẻ hiểu. Nhưng Suisei thì vẫn chưa yên tâm, cô nàng nhìn chai giấm đặt sẵn trên bàn rồi khẽ rùng mình, thậm chí cô còn quặn môi vào bên trong, như thể cảm nhận được vị chua của gia vị này.

Người phụ nữ nội trợ tinh ý nhận ra, liền nghiêng đầu hỏi: ``Sao trông cháu có vẻ run thế, Hoshimachi-chan?''

Nàng Ca sĩ ngập ngừng, tay gãi gãi má, ánh mắt có chút lo lắng: ``Nó\dots\ có kinh dị thật không cô? Cháu thấy nghe thôi đã hơi run rồi\dots''

Kaien cười hiền, bước lại gần, nhẹ nhàng xoa xoa vai cô nàng: ``Không đâu, cô đảm bảo món này ngon tuyệt cú mèo luôn! Ăn một lần là nhớ mãi đó.''

Cô nàng nhìn bà, vẫn có chút ngờ vực, liếc sang những người còn lại rồi gặng hỏi lần nữa: ``Thật không cô? Thật sự không có gì\dots\ đáng sợ chứ?''

Bà gật đầu chắc nịch, nụ cười vẫn giữ nguyên: ``Tất nhiên rồi. Cứ làm thử đi, rồi cháu sẽ thấy.''

Suisei nghe thế, cuối cùng cũng chịu đứng vào hàng, ánh mắt vẫn còn thấp thỏm nhưng cũng đã tạm tin tưởng, nhoẻn miệng cười ngại ngùng.

\ct

Kaien bắt đầu giảng giải sơ qua, tay vừa sắp xếp các nguyên liệu ra bàn, vừa nói: ``Trước tiên, ta sẽ chuẩn bị một ít tỏi, hành khô, mấy củ hành tây, vài quả dứa, mấy cọng sả, và cuối cùng là\dots\''

Aki lập tức chen vào, nở nụ cười tinh nghịch: ``Dấm, đúng không ạ?''

``Chuẩn luôn! Bò nhúng dấm thì không thể thiếu giấm rồi,'' bà reo lên, tỏ vẻ phấn khích. ``Đầu tiên, ta phải băm nhỏ tỏi và hành tím trước.''

Subaru tò mò hỏi ngay: ``Ủa, sao phải băm trước hả cô? Lẩu thì cứ bỏ nguyên vào chứ?''

Kaien cười khẽ, ánh mắt lóe lên nét ranh mãnh: ``Cứ băm đi, lát nữa các cháu sẽ hiểu. Cô cần phi thơm lên trước khi cho vào nước lẩu.''

Noel ngơ ngác: ``Ơ, nhưng mà lẩu thì thường có nước luôn chứ, phi làm gì ta?''

Người nội trợ kiêm chủ kho cười bí hiểm, lắc đầu nhẹ: ``Cứ nhìn cô làm đi rồi sẽ rõ. Đây là bí quyết của nhà Hikawa đấy.''

Sau đó bà đảo mắt nhìn quanh, hỏi lớn: ``Trong nhóm này, ai biết cầm dao băm rồi nào?''

Không chần chừ, Aki và Suisei đồng loạt giơ tay lên. Kaien gật đầu hài lòng, rồi lấy từ trên chạn xuống hai con dao băm sáng loáng, cẩn thận đưa cho hai cô gái.

Nàng Ca sĩ vẫn không quên hỏi kỹ: ``Cô, tầm khoảng bao nhiêu tỏi và hành thì đủ ạ?''

Kaien chống nạnh, cười cởi mở: ``Cứ căn theo lượng ăn thôi, càng đông người thì càng phải băm nhiều.''

Aki suy nghĩ một lúc rồi ước lượng: ``Chắc phải tầm một củ rưỡi tỏi với bảy, tám củ hành khô nhỉ?''

Subaru nghe xong liền gật đầu lia lịa: ``Ừ, chắc cỡ chừng đó thật, vì hơn ba mươi người ăn mà\dots\ số lượng không thể ít được đâu.''

Và thế là, dao băm lách cách vang lên, nhóm chính thức bắt tay vào bước đầu tiên của món lẩu, trong lúc bầu không khí vẫn tràn ngập tiếng trò chuyện và tiếng cười.

Suisei và Aki đã nhanh chóng xắn tay áo, chuẩn bị vào việc.

Cũng giống như Choco và Mel ở tầng trên, hai cô gái này ra dáng rất chuyên nghiệp. Nàng Dị giới nhẹ nhàng cầm dao, điều chỉnh góc cắt tỉ mỉ, còn nàng Ca sĩ thì nhanh nhẹn giữ tỏi, băm đều tay, từng nhát dao gọn gàng và dứt khoát.

Kaien đứng bên cạnh, vừa quan sát vừa không giấu được vẻ ngạc nhiên và thán phục. ``Cô không nghĩ hai cháu lại cắt nhanh đến thế đấy,'' bà vừa cười vừa nói, giọng pha chút bất ngờ.

Chỉ trong chưa đầy ba phút, toàn bộ chỗ hành và tỏi cho hơn chục người ăn đã được băm nhuyễn, không thừa một cọng, không sót một lát. Công đoạn đầu tiên hoàn thành một cách gọn gàng và trơn tru đến mức Rushia và Subaru đứng bên cũng phải gật gù, thầm khen ngợi.

Người đàn bà của gia đình tiếp tục cầm lên mấy cọng sả, giảng giải tiếp: ``Tiếp theo, với sả thì ta bỏ phần gốc trắng ở dưới, rồi đập dập cho thơm, sau đó cắt khúc ngắn, dễ cho vào nồi.''

Nghe đến chữ ``đập,'' Noel - cô nàng Hiệp sĩ luôn hăng máu với những việc cần sức mạnh - lập tức sáng bừng mắt. Chẳng biết từ đâu, cô đã cầm sẵn một cái chùy nhỏ, giơ cao lên đầy phấn khích: ``Có ai bảo `đập' sao!? Phần này để em lo!''

Nàng Tomboy vội xua tay, nhắc nhở ngay lập tức: ``Không phải đập kiểu đó đâu, Noel-chan! Nhẹ tay thôi!''

Kaien khẽ cười, rồi đưa tám cọng sả cho Rushia trước. Bà nhẹ nhàng dặn dò: ``Đầu tiên bỏ gốc trắng, rồi hãy chuyển cho Noel-chan đập dập.''

Rushia gật đầu, cẩn thận cắt bỏ phần gốc sả một cách gọn gàng, sau đó đưa từng cọng sả và con dao cho Noel, người lúc này đang đứng bên cạnh, mặt háo hức như một đứa trẻ được phát quà.

``Đây, Noel- Đừng có dùng cây chùy đó!!'' nàng Pháp sư gắt lên khi cái vũ khí của đồng đội đang quay quay, suýt chút nữa thì trúng vào đầu.

Subaru vẫn chưa kịp nói hết câu cảnh báo: ``Đập nhẹ thôi nha, Noel-cha-''

Nhưng chưa để cô nói xong, Noel đã hô lớn một tiếng: ``\textbf{Hây da\~{}!}'' và hoàn toàn bỏ ngoài tai những lời khuyên ngăn từ mọi người.

*RẦM!*

Cái chùy nện xuống cọng sả, hết cọng này đến cọng khác, lực tay mạnh đến mức mặt bàn rung lên từng nhịp. Cứ thế, chỉ sau vài lần ``hây da'' dứt khoát, mấy cọng sả vốn dĩ thơm lừng giờ đã bị đập dẹp lép, bẹp nát như cái bánh giày.

Nhưng không chỉ có sả chịu trận, chiếc thớt gỗ đáng thương cũng ``hy sinh'' ngay tại chỗ, nứt đôi một cách không thương tiếc.

Tệ hơn, một mảng sứ trên bàn nấu ăn cũng bị nứt nẻ tới mức toác ra, rơi lả tả rơi xuống đất.

Cả phòng lặng đi một giây, rồi đồng loạt há hốc mồm.

Noel sững người, mặt tái mét: ``Chết rồi\dots\ em đập mạnh tay quá\dots''

Aki và Suisei nhìn nhau, nhún vai cười khổ: ``Ái chà\dots\ mạnh tay thật luôn kìa.''

Subaru ôm trán, lắc đầu thở dài: ``Mạnh tay cỡ này thì khỏi cần dao cũng thái được luôn ấy, mẹ trẻ! Đồ phá hoại!''

Rushia thì lặng lẽ nhặt lên nửa mảnh thớt gãy, lắc đầu cảm thán: ``Đi luôn cái thớt người ta rồi\dots\ \hl{May mà đó không phải là mình.}''

Kaien cũng chỉ biết cười trừ, tiến tới vỗ vai cô nàng Hiệp sĩ an ủi: ``Thôi, không sao đâu. Cái thớt gỗ đó rẻ mà, mai cô kiếm cái mới là được.''

Bà ấy không dám nói rằng đó là cái thớt của nhà Hijaki vì sợ rằng càng làm cho Noel cảm thấy tội lỗi hơn. Nói đoạn, bà xoay người bước ra khỏi bếp, đi thẳng lên tầng sáu. Một lúc sau, bà quay lại với một chiếc thớt gỗ khác -  chiếc mà Ryujin vừa rửa sạch, sau khi năm cô gái ở tầng trên đã hoàn thành việc thái thịt.

Mọi người lại nhanh chóng xốc lại tinh thần, tiếp tục công việc, chuẩn bị cho món lẩu bò nhúng dấm - trong khi tiếng cười khúc khích vẫn vang vọng khắp phòng, xua tan đi sự cố nhỏ vừa rồi.

Noel, dù vừa mới gây ra ``thảm họa'' với cái thớt gỗ đáng thương, vẫn chưa có dấu hiệu giảm nhiệt. Cô nàng vừa thấy Kaien quay lại với thớt mới, đã hào hứng giơ tay xin tiếp: ``Để em tiếp tục cắt nha, lần này em cẩn thận hơn, hứa luôn!''

Subaru lập tức khoanh tay, nhíu mày nhìn cô nàng Hiệp sĩ với ánh mắt nghiêm túc hiếm thấy: ``Thôi thôi! Noel-chan, em cứ đứng yên đấy giùm chị. Cái thớt thứ hai không chịu nổi em đâu!''

Rushia cũng chen vào, vừa đẩy nhẹ Noel sang một bên vừa trêu: ``Ừ đấy, nếu lại `hây da' nữa thì chắc lần này thớt bể, bàn bể, có khi sàn nhà cũng thủng luôn-nanodesu!''

Noel đành bĩu môi, ôm chùy đứng nhìn, trông có chút hụt hẫng nhưng vẫn ngoan ngoãn nghe lời.

Người phụ nữ nội trợ khẽ bật cười, lắc đầu nhẹ, rồi ánh mắt vô tình liếc lại đống thịt mà năm cô gái đã mua trước đó - một ``núi thịt'' chạm trần, so với phần thịt bò mà bà đã mua sẵn cho bữa ăn gia đình ban đầu.

Bà nhấc thử tảng thịt trong tay, rồi liếc sang nhìn núi thịt của nhóm \hll, thở dài: ``Không biết chừng này có đủ cho cả nhà không đây\dots\ Cô cứ cảm giác như phần thịt cô mua giống muối bỏ biển so với đống các cháu vác về ấy.''

Subaru cười ha hả, xua tay giải thích: ``Cháu nghĩ chắc không thiếu đâu cô. Bọn cháu còn có cả thịt nướng nữa mà!''

Aki gật gù, tiếp lời bằng giọng điềm đạm: ``Bọn cháu cũng tính kỹ rồi, chắc thế này là đủ cho ba mươi người no căng bụng luôn ạ.''

Kaien nhún vai, mỉm cười: ``Ừ, hy vọng là thế. Cô cứ tưởng là thiếu thịt cơ\dots\ Thôi, tiếp tục nào.''

Bà vừa nói vừa mở túi lấy ra mấy quả dứa chín vàng mà các cô gái đã mua trước đó từ trung tâm thương mại, cùng với vài củ hành tây tròn căng và những bắp ngô vàng tươm mà bà đã mua sẵn ở chợ, và đặt xuống bàn cạnh tảng thịt bò tươi rói.

``Với dứa, mấy đứa gọt vỏ, bỏ mắt rồi cắt lát mỏng ra. Ngô thì cắt ra từng lõi, sau này dùng để chế nước. Hành tây thì cắt làm đôi, một nửa băm nhuyễn để trộn vào thịt ướp, phần còn lại thái khoanh, lát nữa thả vào nồi lẩu. Thịt bò thì cắt miếng vừa ăn, cho vào bát to chuẩn bị ướp,'' bà từ tốn giải thích cho công dụng từng thành phần.

Suisei cầm củ hành tây lên, nghiêng đầu thắc mắc: ``Cô định ướp thịt bằng hành tây luôn ạ? Nghe có gì đó\dots\ là lạ.''

Người nội trợ cười hiền, gật đầu đáp: ``Ừ, đúng thế. Hành tây băm sẽ làm mềm thịt, vị có hơi khác một chút, nhưng đảm bảo ngon đến mức mấy đứa ăn xong sẽ muốn đổi vị liền.''

Nàng Ca sĩ nghe bà giải thích càng tỏ vẻ bán tín bán nghi hơn, nhưng rồi cuối cùng cũng chỉ có thể gật đầu đồng tình. Cô cầm dao bắt tay vào cắt hành tây cùng với Aki và Rushia.

Ba cô gái luân phiên nhau làm: Aki tỉ mỉ cắt dứa thành những lát mỏng tròn đều; Rushia băm hành tây, đôi mắt cay xè đến nỗi phải chớp lia lịa; còn Suisei vừa cắt hành vừa liên tục xì mũi, nước mắt chảy dài vì cay, miệng không quên than thở: ``Cháu không nghĩ hành tây lại\dots\ tàn nhẫn đến thế đấy.''

Subaru đứng cạnh nhắc vui: ``Chịu khó đi, senpai! Đổi lại chị sẽ được ăn thịt ngập mặt!''

Còn Noel thì đứng một bên lẩm bẩm, hai tay cầm chùy ôm chặt, mắt nhìn thèm thuồng mấy miếng thịt bò cắt ra: ``Ước gì có bài tập nào\dots\ dùng chùy để cắt thịt nhỉ\dots''

``Tôi nghĩ là bà nên im lặng mà đứng đó nhìn đi-nanodesu,'' nàng Pháp sư cau mày, mắt vẫn tập trung vào những lát hành tây. ``\textbf{Cay quá!!}''

Không lâu sau, phần hành tây băm nhuyễn được gom lại đầy một bát lớn. Kaien gật gù, tiếp tục cầm tay chỉ việc: ``Phần hành tây này cho vào chỗ thịt bò vừa cắt xong nhé.''

Subaru và Noel nhanh tay dùng muỗng gạt phần hành tây băm vào tô thịt. Hai cô còn đeo bao tay vào, chuẩn bị trộn đều, vừa làm vừa trò chuyện líu ríu.

Nhưng trong lúc trộn, Noel vì quen tay ``dùng lực'' mà vô tình trộn mạnh đến mức miếng thịt văng khỏi bát, bay thẳng vào mặt nàng Vịt.

*TOẸT!*

``\textbf{Ááá! Noel-chaaan!!}'' Subaru hét toáng lên, đứng hình mất mấy giây, trong khi cả nhóm cười lăn cười bò.

Nàng Hiệp sĩ giật mình, luống cuống gỡ miếng thịt bò khỏi mặt đồng đội, vội vàng xin lỗi: ``Chết rồi! Lần này em trộn nhẹ thật mà, không cố ý đâu!''

Rushia vừa lau nước mắt vì cười, vừa lắc đầu: ``Không ai tin đâu, mẹ trẻ!''

Kaien chỉ biết lắc đầu cười trừ, nhưng ánh mắt vẫn dịu dàng, không hề tỏ vẻ giận dữ: ``Không sao, mấy đứa cứ trộn cho đều vào là được.''

Cả nhóm tiếp tục hoàn thiện nốt phần ướp thịt, tiếng cười nói lại rộn ràng khắp căn bếp tầng bốn, khiến cho bầu không khí của buổi chuẩn bị sinh nhật bí mật này càng lúc càng trở nên sống động hơn.

\dots Chỉ ngoại trừ nàng Tomboy. Cô nàng cứ liếc mãi Noel bằng ánh mắt hình viên đạn khi cô đã bị một miếng thịt ``tấn công'' nhờ công của cô ấy.

\ct

Sau khoảng ba mươi phút chờ thịt ngấm gia vị, cả sáu người cuối cùng cũng hoàn thành xong công đoạn ướp. Mùi thơm của tiêu, sả, riềng và hành khô quyện đều vào từng thớ thịt, khiến ai nấy đều có cảm giác bụng bắt đầu réo lên dù chưa nấu nướng gì nhiều.

Noel, vừa tháo bao tay, vừa tò mò quay sang hỏi: ``Còn nguyên liệu nào nữa không cô ơi?''

Kaien nhẹ nhàng lau tay, rồi gật đầu, nói với nụ cười hiền hậu: ``Chuẩn bị thêm đậu phụ và\dots\ trứng vịt lộn nữa là ta có thể ăn được rồi!''

Chỉ nghe đến hai từ ``trứng lộn,'' bốn cô gái thần tượng-kiêm-hề lập tức đồng loạt khựng lại như bị đóng băng, ánh mắt hoảng hốt như thể vừa nghe tin sét đánh.

Noel trố mắt, nuốt nước bọt đánh ực: ``Tr-Trứng lộn\dots\ thật sao!?''

Suisei đứng đờ người, tay siết chặt vạt áo, miệng lắp bắp không nói thành câu, ánh mắt sợ hãi nhưng vẫn tò mò.

Aki rùng mình, mặt méo xệch: ``Cái món kinh dị ấy\dots\ cũng cho vào lẩu sao!?''

Subaru giơ tay ôm đầu, lắc lắc, giọng yếu ớt: ``T-Thật dã man\dots\ Còn tệ hơn vịt lẩu đầu năm của em luôn ấy!''

Trái lại, chỉ có Rushia là tròn mắt ngạc nhiên, thậm chí có phần thích thú: ``Có cả trứng lộn nữa ạ!? Nghe hay đấy chứ-nanodesu!''

Người phụ nữ gia đình Kuon bật cười, giải thích: ``Ở đây có một nguyên tắc rất dễ thương cho món lẩu: bất cứ thứ gì hợp khẩu vị, hay chỉ đơn giản là ăn được, thì cứ bỏ vào, không cần câu nệ gì cả.''

Ngay lập tức, Noel và Aki như bắt trúng cơ hội, lục trong túi lấy ra\dots\ một túi bánh quy và một củ cà rốt, giơ lên cùng đồng thanh: ``Vậy bánh quy với cà rốt thì có bỏ vào được không ạ!?''

Nàng Vịt vội vàng chen vào, giọng chán nản: `` `Bất cứ thứ gì' là mấy món hợp khẩu vị thôi, hai người đừng có bịa đại!''

Nàng Ca sĩ khoanh tay, nhíu mày phân tích: ``Bánh quy thì dẹp đi nha, đồ ngọt mà bỏ vào lẩu thì thành canh tráng miệng luôn rồi. Nhưng cà rốt thì hợp đó, Danchou.''

Nàng Pháp sư cười khúc khích, hí hửng bổ sung: ``Tất nhiên là cả rau củ cũng tính luôn nhé-nanodesu\~{}''

Suisei nghe đến ``rau'' thì mặt xịu xuống ngay, khoanh tay quay mặt đi: ``Đừng có nói chuyện rau với chị, ăn thịt thôi đủ rồi!''

Kaien nghe cả nhóm tranh luận nhau mà cũng phải bật cười: ``Chính xác, rau thì càng phải có cho món lẩu! Mấy đứa cứ nhớ: nguyên liệu nào hợp với món lẩu thì cho vào, còn bánh quy hay kẹo ngậm thì nên để dành tráng miệng. Nhớ chưa?''

Tất cả đồng loạt hô to: ``\textbf{Dạ, hiểu rồi ạ!}''

Không khí trong phòng lại trở nên sôi động. Cả nhóm bắt đầu lấy ra từng món nguyên liệu phù hợp với sở thích của mình, từ rau củ, nấm, đậu phụ cho tới các loại topping lẩu quen thuộc.

Sau đó, bước cuối cùng của nhóm - nấu nước dùng - cũng được bắt đầu. Dưới sự hướng dẫn của người phụ nữ gia đình Hikawa, năm cô gái cùng nhau chia nhau chuẩn bị tám cái nồi to: bốn chiếc lấy từ văn phòng \hll\ mang theo, bốn cái còn lại thì lấy từ trên gác và trong nhà bếp này.

Ai nấy đều khẩn trương nhưng vui vẻ, từng nồi nước được đun sôi, mùi dấm chua thanh thoang thoảng lan tỏa khắp phòng bếp tầng bốn, khiến không khí càng lúc càng rộn ràng, hứa hẹn một buổi tiệc đặc biệt.

Mười phút sau, tất cả tám nồi nước dùng đều đã sẵn sàng, bốc khói nghi ngút, mùi thơm đặc trưng của món bò nhúng dấm khiến ai cũng phải hít hà.

Trong lúc chuẩn bị bắc nồi xuống tầng hai, cả nhóm còn tranh thủ chuyện phiếm, mỗi người một câu, rôm rả chẳng khác gì một buổi dã ngoại.

Aki cười, quay sang Suisei trêu: ``Lát nữa chị có dám ăn trứng lộn không đấy?''

Nàng ``Sao chổi'' đỏ mặt, lườm nhẹ về đàn em của mình, lí nhí: ``Chị mà không ăn thì ai dám ăn hả trời\dots\ phải ăn trước mặt mọi người cho đỡ mất mặt chứ!''

Noel thì vừa khuân nồi vừa hào hứng lẩm bẩm: ``Không biết lát nữa Kuon-senpai sẽ ngạc nhiên thế nào nhỉ\dots\ Chắc vui lắm đây!''

Subaru chống nạnh, nói lớn: ``Anh ấy mà không rớt hàm vì ngạc nhiên thì em thề đổi nghề luôn.''

Rushia cười khúc khích, tay bê nồi, vừa đi vừa bảo: ``Chỉ mong lẩu không hết trước khi Kuon-senpai bước vào thôi!''

Rồi cả nhóm cùng nhau bê từng nồi nước dùng nóng hổi xuống tầng hai, chuẩn bị cho một bữa tiệc sinh nhật thấm đẫm tình cảm - mà một lần nữa, nhân vật chính vẫn chưa hay biết điều gì.

\subsubsection*{\phantomsection\label{P4.1c}}
\addcontentsline{toc}{subsubsection}{Nhóm 3 - Nhóm làm bánh sinh nhật}

\begin{center}
  {\color{T2} Nhóm 3 - Ryuuhou và các thành viên từ nhóm 1, nhóm 2\\Nhóm làm bánh sinh nhật}
\end{center}

``Một bữa tiệc sinh nhật mà không có bánh kem thì quả đúng là điều thiếu sót.'' Hoặc chí ít, đó là những gì mà Aki nêu ra.

Công đoạn nấu ăn cuối cùng mà họ sẽ phải trải qua, thậm chí có thể gọi là ``trùm cuối'' - đó là làm món bánh kem sinh nhật.

Và đó không phải là một chiếc bánh bình thường, mà mọi người dự định sẽ làm một chiếc bánh khổng lồ nhiều tầng, đủ cho hơn ba mươi người ăn.

Nhưng rồi, có một vấn đề rất nhanh chóng bày ra trước mắt: cả căn bếp tầng sáu tại phòng Kuon lẫn tầng bốn của phòng Hijaki đều không đủ rộng để nhóm đông người cùng lúc làm bánh. Nồi lẩu, thịt nướng, nguyên liệu các nhóm khác đã chiếm gần hết không gian. Dù có cố gắng xoay xở kiểu nào thì cũng không thể chen thêm một nhóm lớn vào, chứ chưa nói đến việc phải nhồi nhét cả nguyên liệu vào nữa.

Giữa lúc cả nhóm còn đang ngồi vòng tròn trên sàn tầng ba bàn bạc về kế hoạch tiếp theo, Shion bỗng nhiên chống cằm, nheo mắt nhìn trần nhà rồi chỉ tay lên phía trên: ``Này, tầng bảy có một căn phòng trống mà, sao không lên đó thử nhỉ?''

Aqua cũng phụ họa, hồn nhiên chen vào: ``Đúng rồi đó! Tôi nhớ tầng hai hình như cũng có phòng bỏ trống thì phải? Nếu không đủ chỗ ở trên thì xuống đó luôn!''

Nghe đến đây, Ryuuhou - người từ nãy tới giờ vẫn giữ vẻ bình tĩnh, vừa mới thở phào sau khi sóng gió mang tên ``Onii-san'' vừa qua - lập tức giật nảy mình, tim đập loạn xạ. Cái tên tầng bảy và tầng hai vừa vang lên, như một luồng điện lạnh chạy dọc sống lưng cô bé.

Căn phòng tầng bảy vốn là của Kiryu Coco, người đang ra ngoài có việc, và vì thế căn phòng được cô nàng khóa chặt cẩn thận.

Trong khi đó, tầng hai, nơi Mano Aloe đang sống và cũng đang không có nhà, nơi từng là cửa hàng viễn thông cũ đã đóng sau đại dịch. Dù biết các chị sẽ không hiểu, nhưng bản năng của một người ``biết nhiều hơn mức cần thiết'' khiến Ryuuhou không thể để họ tuỳ tiện vào đó.

Bối rối, cô bé nhanh chóng nghĩ ra một cái cớ: ``Cái\dots\ cái phòng tầng hai đó là\dots\ phòng riêng của em đấy! Nhưng mà\dots\ em chưa dọn đâu, bừa bộn lắm, các chị không vào được đâu!''

``Vậy còn căn phòng ở tầng bảy thì sao?'' Mel hỏi.

``Căn phòng đó thì\dots\ dành cho khách. Nhưng mà nó bị khóa rồi ạ. Gia đình em không có ai cầm chìa khóa, chắc là bác Hijaki-san cầm đấy ạ\dots'' Ryuuhou bối rối trả lời, cố gắng tìm cách biện minh.

``Bác Hijaki?'' tất cả mọi người cùng nhau đồng thanh, nghiêng đầu khó hiểu.

``À thì là\dots! Bác ấy sống chung với bọn em đấy mà! Nhưng mà bác ấy\dots\ đi ra ngoài có việc rồi!''

``Vậy à\dots'' Choco thở dài, tỏ vẻ hơi thất vọng. ``Giá như bác ấy có ở đây thì có thể hỏi bác ấy nhỉ\dots''

``E-Em nghĩ là không đâu!''

Vừa dứt lời, Noel lập tức xích lại gần, ánh mắt long lanh của một cô nàng Hiệp sĩ pha chút tò mò: ``Nhà đông người thế, mà phòng riêng cũng không cho bọn chị vào tham quan sao\~{}? Có gì đâu, chị với Shion-senpai dọn tí là xong mà!''

Luna thì\dots\ có lẽ khỏi phải nói. Cô nàng đưa đôi mắt long lanh về phía cô em gái, hai tay chắp trước ngực nũng nịu: ``X-Xin chụy đó, Ryuuhousa-chan! Làm bánh ở đây chật quá nanora\~{}! Em chịu hông nỗi đâu na\~{}!''

Trước ánh mắt cún con của nàng Công chúa xứ Kẹo và những lời thúc giục không ngừng của nàng Hiệp sĩ, Ryuuhou chỉ còn biết ngậm ngùi buông tay đầu hàng, miễn cưỡng gật đầu, giọng yếu ớt: ``\dots Được rồi\dots\ nhưng\dots\ các chị cứ chuẩn bị nguyên liệu đi trước đi, để em xuống trước dọn một chút đã!''

Vừa nói dứt câu, cô bé đã quay ngoắt người, rảo bước như chạy xuống tầng hai, lòng thầm cầu mong: ``\textit{Mong là họ không lục tung mọi thứ lên\dots\ Không thể để căn phòng của Aloe-chan bị phát hiện!}''

Còn phía sau, các cô gái định bước theo phụ giúp, nhưng cô em gái của Kuon đã chặn lại ngay trước bậc thang, nhoẻn miệng cười gượng: ``M-Mình em làm được mà! Các chị cứ chuẩn bị nguyên liệu mang xuống đi, đừng lo cho em!''

Và thế là, cả nhóm đành gật đầu đồng ý, chia nhau gom nguyên liệu, trong khi cô em gái của Kuon thì nhanh chóng trượt xuống cầu thang, lao thẳng về phía tầng hai - nơi cất giấu những bí mật mà cô bé buộc phải bảo vệ, đúng theo thỉnh cầu của YAGOO tới gia đình Hikawa.

\ct

Ryuuhou chạy vội xuống tầng dưới, gần như trượt chân xuống từng bậc, tay bám vịn đầy lỏng lẻo, tim đập thình thịch vì lo lắng. Vừa tới nơi, cô bé liền rút chìa khóa mở cửa phòng Aloe - căn phòng mà cả gia đình đều tuyệt đối giữ bí mật vì lời dặn nhẹ nhàng mà cương quyết của YAGOO: ``Nếu có thể\dots\ xin đừng để ai biết họ đang sống ở đây cho tới thời điểm thích hợp.''

Cánh cửa mở ra.

Ngay lập tức, mùi dịu nhẹ của tinh dầu hoa oải hương và hơi thở ấm áp quen thuộc của căn phòng tràn ra, khiến Ryuuhou càng thêm cuống cuồng.

Căn phòng của Aloe giờ đây trông rộng rãi và tiện nghi hơn bao giờ hết. Khác hẳn với đợt mà cô nàng đã dọn tới từ năm tháng trước, lần này cô cũng đã sơn căn phòng thành màu hồng nhạt và có chút áng xanh lá, cùng với việc có treo thêm vài món đồ lỉnh kỉnh khác.

Một bên được bố trí như một studio thu nhỏ, đầy đủ loa cỡ lớn, mixer, bộ máy tính cây, micro và tai nghe treo gọn gàng - chẳng khác nào phòng làm việc của một nghệ sĩ thực thụ. Bên còn lại là khu vực sinh hoạt cá nhân: chiếc giường đơn kiểu Nhật được xếp sát tường, một tủ quần áo nhỏ, kệ sách chứa đầy truyện tranh và sách hướng dẫn thanh nhạc. Những dải đèn LED viền góc phòng toát lên thứ ánh sáng dịu nhẹ màu tím hồng, bao phủ không gian bằng đúng cái không khí mang dáng dấp của nàng Succubus.

Cô em gái của Kuon nuốt nước bọt. Cô bé chạy nhanh vào phòng, kéo màn rèm lại để ánh sáng dịu bớt, rồi túm lấy một tấm khăn trải giường màu xám vắt sau ghế, phủ lên toàn bộ khu vực studio như thể đó chỉ là một ``khu làm việc bị bỏ hoang.'' Micro được tháo xuống và nhét tạm vào ngăn kéo tủ quần áo. Tai nghe bị giấu dưới gối, còn những poster nhỏ dán quanh tường cũng được cô gỡ xuống đầy cẩn thận nhưng nhanh nhẹn và xếp tạm vào hộp bìa các-tông. Chiếc hộp sau đó được cất tạm xuống kho ở bên rìa cầu thang.

Bên phía giường ngủ, Ryuuhou chỉ kịp gấp chăn sơ sài, rồi nhét hết gối vào hộc tủ. Cô kéo ngăn sách phía dưới ra, tống đại mấy cuốn manga về Succubus vào bên trong rồi đóng *cách* lại - mồ hôi ướt cả trán, cả người cô như đang bước vào kỳ thi vào cấp ba - nơi mà cô đã thất bại lúc đó.

Chỉ vài phút sau, căn phòng đã trông ``vô hồn'' hơn hẳn - ít nhất là với người ngoài. Cô thở hắt ra một hơi, chưa kịp lau mồ hôi thì nghe tiếng chân rộn ràng từ cầu thang vọng xuống.

``Ryuuhou-chan\~{}! Bọn chị mang xuống đồ xuống rồi đây\~{}'' giọng của Noel vang lên.

Theo sau đó là giọng nói nũng nịu đầy phấn khích của Luna: ``Có cả kem nữa nì-nora\~{} Cả bộc và xi-rô nữa-na\~{}!!''

Cô bé giật mình bật dậy, vừa kịp đóng cửa tủ thì cánh cửa phòng bật mở. Và rồi\dots\ một đoàn rước ``đại lễ nguyên liệu'' tiến vào như diễu binh. Những thứ mà họ mang theo - cả về số lượng lẫn chất lượng - khiến cho cô chết lặng.

Okayu và Korone dẫn đầu, mỗi người ôm một đống lọ đủ màu sắc: si-rô, hương liệu, tinh dầu thơm đủ vị từ dâu, socola, bạc hà đến\dots\ matcha. Theo sau là ba hộp bơ to gần bằng thùng mì tôm và ba hộp sữa lớn trông như dùng cho nhà hàng.

Nàng Khuyển hí hửng nói: ``Tụi em gom sạch kệ hàng khu nguyên liệu bánh á! Mỗi thứ ba phần cho chắc!''

Phía sau, Ayame bước vào như một tay buôn bột, trên tay là ba, bốn túi bột mì và bột nở, mỗi túi to như túi gạo mười ký. Kế đó là Shion tay xách hai túi đường và\dots\ một gói muối khổng lồ. Aqua lạch bạch theo sau với hai tá trứng ôm đầy tay, đi tới đâu là trứng lăn lóc tới đó. Cuối cùng, Suisei và Rushia cẩn thận bê một khay nhỏ chứa tinh dầu vani và kẹp nến trang trí.

Ryuuhou đứng ở cửa, miệng há ra vì sốc: ``\dots M-Mấy chị định mở tiệm bánh hay là làm bánh sinh nhật vậy\dots?''

Những ánh mắt long lanh quay sang nhìn cô, ai cũng cười toe toét đầy phấn khởi như vừa khám phá ra thiên đường nguyên liệu. Nàng Kẹo hớn hở đáp: ``Yên tâm đi, Ryuuhousa-chan, chỉ cần đủ đồ thì bánh từng nào cũng làm được hết-nora\~{}!''

``\dots Chị thực sự ngạc nhiên khi thấy mấy đứa mua nhiều nguyên liệu đến vậy đấy\dots'' Mel tỏ vẻ sửng sốt, trong khi hai tay vẫn đang xách đống lọ hương liệu.

``Thì đúng là chúng ta làm bánh cho tất cả mọi người ăn mà!'' nàng Miêu hồ hởi nói.

Cứ thế, hàng đống nguyên liệu cứ dần dà được mang xuống, tới mức lấp đầy nửa căn phòng. Ryuuhou chỉ biết đứng nhìn, lòng thầm nghĩ: ``\textit{Không biết làm xong bánh thì căn phòng này còn chỗ để thở không nữa\dots}''

Luna, trong lúc đặt hộp bơ xuống bàn, quay sang hỏi đầy háo hức: ``Ryuuhousa-chan, em có biết làm bánh không-nora?''

Cô bé giật mình, gãi đầu cười gượng: ``E-Em biết chút chút thôi\dots\ nhưng mà làm bánh lớn thế này thì chưa bao giờ thử.''

Aqua hí hửng chen vào: ``Không sao đâu! Cứ để chị chỉ cho! Chị có mua một cuốn sách hướng dẫn đây rồi!'' và lấy ra một cuốn sách mà chính tay cô mua khi vừa tách đoàn.

Shion nghe vậy liền bật cười, trêu chọc: ``Cần gì mua sách khi ta có thể tra mạng?''

``V-Vô duyên!''

Cả phòng bật cười rộn rã, không khí trở nên nhẹ nhàng hơn. Ryuuhou cũng khẽ mỉm cười, cảm giác lo lắng và hồi hộp ban nãy dần tan biến.

Cô bé hít một hơi thật sâu, rồi nói: ``Vậy thì bắt đầu thôi, các chị!''

``\textbf{Được rồi\~{}!}'' cả nhóm đồng thanh reo lên, bắt tay vào công việc với tinh thần đầy phấn khởi.

\ct

Ngay sau khi mọi người đặt những đống nguyên liệu như ``tiệm bán sỉ'' xuống sàn phòng, Ayame liếc nhìn quanh một lượt rồi thở dài một tiếng: ``Chà\dots\ chỗ này nhìn vui đó, yo. Nhưng mà mình với Korone-chan phải quay lại tầng ba thôi. Bọn mình còn chưa gắn hết dây đèn với sơn tường xong đâu.''

Korone đặt nhẹ lọ tinh dầu xuống bàn, quay sang vỗ vai Ryuuhou: ``Để bọn này lo phần trang trí! Còn bánh thì cứ để mọi người ở dưới đây lo nhen\~{}''

``V-Vâng ạ,'' Ryuuhou gật đầu, vẫn còn chưa hoàn hồn sau màn vận chuyển nguyên liệu siêu cấp. Hai người rời khỏi phòng, để lại một đám hỗn binh trong tay cô bé.

Ngay khi bầu không khí vừa mới lặng xuống đôi chút, hai cái đầu khác cùng ló lên - và cả hai đều tự nhận mình\dots\ có quyền chỉ huy.

Không ai khác, đó là nàng Hầu gái và nàng Phù thủy nhỏ - là cặp bạn thân.

``Để chị hướng dẫn mọi người nha! Chị có mua sách dạy làm bánh hồi sáng đây!'' Aqua lên tiếng trước, giơ cuốn sách nhỏ in hình bánh kem lòe loẹt, cười rạng rỡ như vừa tìm được kho báu.

``Ê này! Tôi cũng có đây!'' Shion lập tức rút điện thoại, vẫy vẫy trước mặt mọi người. ``Tôi thì không tin được công thức trên sách lắm đâu. Phải tra kỹ mới đúng chuẩn được, còn tỷ lệ nước, trứng, bơ nữa.''

Nàng Hầu gái nhíu mày, không chịu lép vế ``Nhưng sách tôi là bản tiếng Nhật nguyên gốc luôn đó! Có cả hình minh hoạ từng bước nữa!''

``Vậy bà định lật từng trang cho cả đám sao?'' Shion chọc giận cô bạn.

``Còn bà định đọc bài viết dài hai mét trên mạng chắc?'' Aqua không chịu kém cạnh.

Bầu không khí bắt đầu nảy lửa. Những người xung quanh chưa ai kịp làm gì đã thấy hai ``người hướng dẫn'' tranh nhau quyền lực như trong một game show nấu ăn hỗn loạn.

Ryuuhou ở chính giữa, mắt đảo như chong chóng giữa hai ``senpai'' hơn tuổi nhưng\dots\ chẳng lớn hơn mình là bao, thậm chí cô có thể xưng họ là ``-chan'', nhưng rốt cuộc vẫn phải hạ mình xuống.

Khác với ông anh trai, người tiếp xúc hầu như thường xuyên với các cô gái và từ đó biết cách kiểm soát tình hình, thì với cô em gái của cậu, đây mới là lần thứ hai cô gặp họ, và vì thế cô không biết phải xử lý tình huống như thế nào.

Sau vài giây quan sát tình hình, cô bé hắng giọng, khẽ nói: ``V-Vậy thì\dots\ hai chị làm chung luôn đi ạ. Một người xem sách, một người tra mạng. Có gì khó thì hỏi em ạ.

Aqua và Shion cùng quay sang, định phản đối, nhưng lại nhìn nhau một chút rồi\dots\ gật đầu gần như đồng thời:

\begin{dialogue}
  \speak{Aqua} Ừ\dots\ Cũng được.
  \speak{Shion} Ờm\dots\ vậy cũng ổn.
\end{dialogue}

Cuộc đấu tạm thời lắng dịu, làm cho cô em gái của Kuon thở phào nhẹ nhõm. Choco nhìn hai cô nàng vừa cãi nhau đó, nét mặt không vui: ``Hai đứa đang làm cho Ryuuhou-sama cảm thấy khó xử đấy. Em ấy không giống như Kuon-sama đâu.''

``Với cả, chị nghĩ rằng\dots\ cái gì nhỉ, ba cây chụm lại nên hòn núi cao chứ?'' Aki lên tiếng thắc mắc. ``Chẳng phải nhiều người làm cùng lúc sẽ hiệu quả hơn sao?''

Nàng Hầu gái và nàng Phù thủy nhìn nhau lần nữa, mặt đỏ hỏn không biết phải trả lời hai người lớn như thế nào.

\ct

Trong lúc hai ``hướng dẫn viên'' bắt đầu tranh luận xem nên dùng vani dạng bột hay dạng lỏng, thì Ryuuhou nhanh chóng tranh thủ mở ngăn kéo tủ, lôi ra các vật dụng làm bánh chuyên dụng mà nhà cô bé và Aloe từng để lại: máy đánh trứng cỡ vừa, khuôn tròn ba cỡ khác nhau, giấy nến, dao cắt bột, cán lăn bột và cả một bộ đui bắt kem ba mươi chi tiết.

``\vie{\textit{Ít ra mình còn kiểm soát được phần dụng cụ\dots}}'' - cô bé thì thầm, rồi tiếp tục kiểm tra độ sạch của từng món, sắp xếp lên mặt bàn. ``\vie{\textit{Không biết các senpai bàn bạc đến đâu rồi\dots}}''

Phía sau lưng, Aqua và Shion đã bắt đầu\dots\ cãi nhau vụn về việc trứng nên để nhiệt độ phòng hay lấy từ tủ lạnh là ngon hơn.

Subaru nghe cuộc tranh luận của hai người mà thở dài ngán ngẩm: ``Trứng thì thế nào mà chả được! Mau bắt đầu nhanh lẹ lẹ đi!''

``Liệu để hai người họ điều khiển này có ổn  không đây-nanodesu\dots'' Rushia đứng bên cạnh cô, mặt cũng tỏ thái độ bực bội.

Và thế là, hành trình làm chiếc bánh kem nhiều tầng bắt đầu một cách không hề suôn sẻ - với hai người chỉ đường, một người điều phối, và hàng tá các thần tượng phụ việc đang chuẩn bị nhúng tay vào ``cuộc hỗn chiến bột đường'' độc nhất vô nhị.

``Trước tiên thì chúng ta phải chia nhóm trước nhỉ?'' cô em gái của Kuon đứng chống hông hỏi, cố gắng mang dáng dấp của một người điều phối trưởng, dù trong lòng thì đang thầm tự nhủ: ``\textit{\vie{Mình còn chưa biết ai đánh trứng giỏi hay ai giỏi trộn bột cơ mà\dots}}''

Cả nhóm chợt im lặng vài giây, rồi bắt đầu rì rầm, tranh nhau làm công việc mà họ cho là làm giỏi nhất.

Aqua nhảy dựng lên đầu tiên, tay giơ cao như học sinh giành phát biểu: ``Chị muốn làm nhóm đánh trứng! Nhìn trứng quay quay vui lắm luôn đó!''

Shion bĩu môi: ``Vậy thì để chị trộn bột! Đừng có ai tranh phần này nha\~{}''

Okayu thì nhỏ nhẹ theo sau: ``Vậy em cũng làm trộn bột cho dễ phối hợp\dots''

Noel gật gù: ``Cho em làm với Okayu-senpai, nhóm này có vẻ lành tính, hợp với em á.''

Subaru chống nạnh, nói chắc nịch: ``Chị thấy đánh trứng có vẻ hợp! Dễ điều khiển máy móc hơn!''

Luna thì giơ hai tay lên, lon ton bước tới bên Subaru: ``Em cũng mún làm chứng\~{} trứng đánh zui lắm-nora\~{}!''

Rushia thì chỉ đứng ngó nghiêng nhìn các thành viên đang nhốn nháo, không rõ mình nên làm gì.

Mel, Choco, Aki và Suisei thì bắt đầu bàn nhau xem ai nên đứng lò, ai nên lo khuôn bánh, còn Ryuuhou thì mặt đầy lo lắng, tay ôm sổ ghi chú, mắt lia qua lia lại như đang đọc bảng điểm môn thể dục.

Cô bé giơ tay lên cắt ngang: ``Khoan khoan khoan! Nếu ai cũng chọn theo ý thích thì nhóm này sẽ lệch hẳn đó ạ! Trộn bột và nướng bánh đều cần người có sức nữa!''

``Thế sao không chia theo chiều cao nhỉ?'' Noel đề xuất, tự hào vỗ ngực. ``Chị khá là tự tin với khoản này đó!''

``Khum công bằn!'' Luna phản đối ngay. ``Luna đây thì nùn tịt à-nora!!''

``Thế chia theo\dots\ nhóm máu'' Suisei buột miệng. Ngay sau đó, Aqua reo lên: ``Chia theo cung hoàng đạo?''

Subaru thậm chí còn bông đùa hơn: ``Chị nghĩ là chúng ta nên chia theo\dots độ hỗn?'', chỉ ngay vào Luna.

``Ê ê! Khum chịu đâu-nanora\~{}!'' nàng Kẹo phồng má lên, giương đôi mắt bất cần. ``Nếu như vầy thì Noel-senpai xứng đáng hơn á-nora!''

Những đề xuất cứ thế bay loạn như pháo hoa, không ai chịu ai. Mãi cho đến khi Shion búng tay: ``Thôi oẳn tù tì đi cho rồi! Bốc thăm là công bằng nhất!'' thì lúc này, cả phòng đồng mới chịu đồng ý.

\begin{dialogue}
  \speak{Choco} Cuối cùng cách này vẫn ổn nhất nhỉ.
  \speak{Rushia} Biết thế làm ngay từ đầu có phải tốt hơn không.
  \speak{Mel} Được rồi, mọi người chuẩn bị chưa nào?
\end{dialogue}

Và thế là, giữa một rừng nguyên liệu, giữa tiếng cười nói rôm rả và vài cú cãi nhau lặt vặt, một vòng oẳn tù tì sôi nổi bùng nổ - có người thắng hô to ``YES!'' như đoạt giải Idol quốc dân, có người thua thì\dots\ lăn ra giả vờ xỉu như thể vừa mất suất debut.

Ryuuhou nhìn khung cảnh lộn xộn trước mặt, chỉ có thể thở dài ngao ngán: ``Bước đầu đã mệt thế này rồi\dots\ Không biết liệu mọi chuyện có êm đềm không đây\dots''

Cuối cùng, bằng đủ mọi thể loại ``phương pháp ngẫu nhiên bán nghiêm túc,'' đội hình làm bánh kem đã chính thức được phân chia:

Nhóm đánh trứng gồm Aqua, Subaru, Luna và Rushia.

Nhóm trộn bột gồm Okayu, Shion, Noel.

Nhóm chuẩn bị khuôn và lò nướng gồm Ryuuhou, Aki, Mel, Choco, Suisei.

Cuộc vật lộn với chiếc bánh kem sinh nhật dành cho Kuon của các thành viên \hll\ và cô em gái của cậu chính thức bắt đầu.

\ct

Bộ tứ của nhóm đánh trứng bắt đầu vào trận.

Aqua - người giành được quyền chỉ huy của nhóm này bằng cách\dots\ hô to nhất - lập tức lấy bốn chiếc âu to từ tủ bếp, đập trứng vào từng cái rồi ra lệnh, mắt liếc vào cuốn sách: ``\textbf{Subaru-chan! Rushia-chan! Luna-chan!} Mỗi người đánh một âu kem rồi cho một loại hương liệu nha!''

Nàng Pháp sư thì vẫn còn đang lúi húi đếm trứng trong những hộp vỉ họ mua về: ``Ủa\dots\ là trứng gà hay trứng vịt vậy? Sao có quả trắng, có quả hơi xám xám thế này\dots?''

Nàng Tomboy gãi đầu, bốc bừa mấy quả trứng vỏ nâu lên: ``Chắc là do ánh sáng thôi\dots\ Em đập ra là biết liền.''

Nàng Công chúa xứ Kẹo vẫn đang hì hụi lôi máy đánh trứng ra, chớp chớp mắt: ``Em muốn xài cái này cơ\~{}!''

``Dùng cái đó cũng được thôi,'' Ryuuhou vẫn đang tìm khay trong nhà bếp của Aloe, ``nhưng mà chị có biết dùng không vậy?''

Nàng Hầu gái thấy vậy đáp lại, cố lật đi lật lại trong cuốn sách tìm hướng dẫn sử dụng chiếc máy. Sau một hồi tra cứu tìm kiếm nhưng không có kết quả, cô nghiêm giọng nói đại: ``Phải để máy chạy ít nhất năm phút, tốc độ tăng dần, đừng để nó văng tung tóe nha!''

Thế rồi, ngay khi cô vừa dứt lời, một cái *BỘP!* vang lên, khiến tất cả mọi người khựng lại.

Một quả trứng gà trong tay của Luna rơi thẳng xuống nền nhà. Lòng trắng lòng đỏ loang ra thành một vũng lớn, chảy sát đến chân của Luna, làm cho cô hét toáng lên: ``\textbf{Á á á!! Trứng nổ nanora!!!}''

Do quá hoảng hốt, đôi cao gót của cô trượt phải mớ bầy nhầy đó làm cho Luna bật ngửa ra sau.

``\textbf{Luna-chan!!}''

Choco và Ryuuhou vội lao tới bám lấy cánh tay của cô, cứu thoát cô khỏi một cú ngã đáp đất bằng mông.

``Em/Chị không sao chứ!?'' cả hai người hỏi han, kéo nàng Kẹo đứng dậy.

Ở phía xa, nàng Vịt ngán ngẩm, tới mức độ đổ mồ hôi hột: ``Chưa bắt đầu mà đã hỏng rồi sao\dots?''

Trong mớ hỗn loạn đó, nàng Hầu gái vẫn đang loay hoay với mớ dây điện, tìm cách cắm chiếc máy đánh trứng.

Sau một hồi vật lộn, cuối cùng cô cũng thở phào nhẹ nhõm. ``Được rồi, xem nào\dots'' Cô nhấn nút khởi động chiếc máy.

*CẠCH!*

Và\dots\ không có gì xảy ra cả. ``Ủa? Cái máy này không chạy à?''

Ryuuhou xách những chiếc khuôn bánh ở gần đó, nhìn vào ổ cắm điện, thở dài đáp: ``Aqua-onee-chan chưa đóng điện kìa.''

Nàng Hầu gái và nàng Pháp sư ngoảnh đầu về phía ổ, nhìn thấy đèn công tắc không sáng. Aqua thấy vậy, chỉ nở một nụ cười ngượng ngùng: ``À, ra vậy.''

``Đúng là hầu gái hậu đậu-nanodesu,'' Rushia đáp lời với giọng mỉa mai.

``T-Tại chị quên chứ bộ!''

*CẠCH!*

Công tắc ổ cắm được bật lên, và ngay sau đó\dots

*RỪM RỪM! RỪM!*

Chiếc máy đánh trứng đột ngột khởi động với mức tối đa, không kịp để cho ai trở tay.

``Ể\dots? Ể!? N-Này! C-Chờ--!''

Những lát kem trong bát bắn tung tóe lên cả tóc Aqua, lên áo Subaru, và cả mặt Luna, thậm chí còn vương thẳng lên trần nhà, dưới sàn và tung vào những dụng cụ khác trong nhà bếp. Rushia và Ryuuhou thì may mắn né kịp, chỉ bị vương chút hương.

``\textbf{Ối trời ơi! Cái máy chạy kìa!!}'' nàng Vịt hét lớn, lấy chiếc khay gần đó làm bia đỡ, đồng thời luôn miệng: ``\textbf{Tắt đi, Aqua! Tắt nó đi!}''

Nàng Hầu gái luống cuống tìm công tắc chỉnh tốc độ xuống mức thấp nhất. Tới lúc này, cô mới thở phào nhẹ nhõm, nhưng khi nhìn lại, căn bếp nhà cô nàng Succubus - mà họ không hề hay biết - đã trở thành một chiến trường đầy mùi kem và hương liệu.

``\vie{\textit{Dọn chỗ này sẽ ốm người đấy\dots}}'' cô em gái của Kuon thở dài chán chường.

Luna đưa tay quẹt kem trứng trên má, liếm thử, rồi nở một nụ cười ngây thơ: ``Ngon quá na! Kem này chất lượng đó-nanora! Có hương vị quýt nữa!''

Subaru đập trán ngán ngẩm: ``Tất nhiên nó phải ngon rồi! Mà này, em ăn thử luôn đấy à!?''

Cô liếc quanh, nhìn đám kem loang lổ khắp sàn rồi thở phào: ``Ít ra vẫn chưa ai bị thương\dots''

*XOẢNG!*

Ngay khi nàng Vịt dứt lời, một tiếng bát vỡ vọng lên khắp căn phòng. Mọi ánh nhìn quay sang nàng Công chúa xứ Kẹo. Dưới chân cô là chiếc bát trứng thứ hai đã vỡ tanh bành, với phần kem đã nhuốm bẩn đôi giày của cô.

``Luna-chan\dots?''

Lập tức, khuôn mặt của nàng Hầu gái và nàng Vịt tối sầm lại, khuôn mặt cau có. Họ cùng đồng thanh: ``Nếu không giúp được gì thì đứng yên một chỗ đi!''

Luna thấy thế, toàn thân run bần bật: ``C-Cho em xin lỗi-nanora!!!''

Ryuuhou chạy lại, căn ngăn giữa hai bên: ``Thôi, thôi được rồi. Dù gì chị ấy cũng chỉ tò mò thôi mà! Hai người đừng giận!''

Rushia ở bên ngoài, lẩm bẩm: ``Giờ mình hiểu vì sao Kuon-senpai ghét em bé rồi-nanodesu\dots''

``Thôi được rồi, tập trung vào làm đi, mọi người!'' cô em gái nhanh chóng lấy lại tinh thần. ``Làm nhanh lên để còn có thời gian trang trí!''

Cuối cùng, nhóm vẫn phải tiếp tục làm việc, hoàn toàn ngó lơ mớ lộn xộn mà chiếc máy đánh trứng và nàng Kẹo đã gây ra.

\ct

Sau một hồi vật lộn với những quả trứng, cuối cùng nhóm Aqua cũng đánh xong phần lòng trắng và lòng đỏ theo chỉ dẫn. Aqua lại nhìn vào tờ ghi chú kèm công thức mà cô viết tay trong cuốn sách mua sáng nay.

``Giờ thì chúng ta sẽ bắt đầu đánh kem lần nữa nha!'' cô nói lớn, theo sau là những cái gật đầu từ Subaru, Luna và Rushia.

Nàng Hầu gái tiếp tục giảng giải: ``Mỗi tầng bánh sẽ có một vị khác nhau, nên ta cần làm năm loại kem! Vị dâu, socola, vani, trà xanh, kem ngô và nhiều vị khác nữa!''

Nghe cô nói thế, cả ba cô gái đồng loạt nghiêng đầu khó hiểu.

``Sao có nhiều loại kem vậy\dots?'' nàng Pháp sư chớp mắt.

``Lại còn có kem vị ngô nữa\dots?'' nàng Vịt nhăn mặt. ``Kem gì lạ vậy?''

Riêng chỉ có nàng Công chúa là vẫn phấn khích, chưa buồn lau toàn bộ số kem trên mặt: ``Nghe có vẻ ngon đó-nora!''

Aqua chỉ về đống lọ hương liệu đã bày sẵn trên một cái bàn, đáp lại: ``Chị đâu có biết đâu, tại chị thấy có nguyên một rổ hương liệu nên nói vậy thôi chứ.''

Okayu giơ tay, khuôn mặt giãn ra như một con mèo, vô tư đáp: ``Nhân tiện thì tôi với Koro-san mua chúng đấy.''

Ryuuhou và Subaru đổ mồ hôi hột, đồng thanh hỏi: ``Sao mà em/bà mua nhiều loại thế?''

``Tại bọn tôi không biết nên chọn loại nào nên cứ hốt tất về thôi\~{}'' nàng Miêu trả lời, không hề mảy may gì.

Chưa kịp để bất cứ ai tra hỏi thêm điều gì, Aqua hắng giọng, lật giở cuốn sách và phân nhiệm vụ: ``Luna-chan đánh kem dâu! Subaru-chan sẽ làm kem sô-cô-la! Còn Rushia-chan sẽ lo vani! Em sẽ lo hai phần còn lại là trà xanh và bắp nha!''

``Chưa kể rằng còn có kem vị cam, kem vị nho với kem vị mâm xôi nữa\dots'' nàng Vịt thở dài, lấy trên chạn một cái bát và một cái khấy đánh trứng kiếm được trên giỏ.

Cứ thế, mỗi người chọn một bát, lấy nguyên liệu riêng và bắt đầu làm mẻ kem đầu tiên.

Nhưng\dots\ mọi chuyện lại diễn ra không hề suôn sẻ.

Ngay khi nàng Hầu gái - người trưởng nhóm tự xưng - hoàn thành xong mẻ kem trà xanh, cô cảm thấy có gì đó không ổn.

``À-rế\dots? Sao cái phần kem này không đặc lại nhỉ?'' cô nhìn bát kem của mình trông không khác một bát canh rau muống được pha loãng.

Ryuuhou - đang sắp xếp khuôn bánh phía bên kia - nghe thấy liền ngẩng đầu: ``Aqua-onee-san, chị cho bao nhiêu whipping cream vậy ạ?''

Aqua gãi đầu: ``Chị thấy trong này ghi là 300ml nên chắc là tầm\dots\ một chai rưỡi.''

Cô đưa hộp sữa kem đã vơi một nửa cho cô em gái của cậu Tổng. Ban đầu, Ryuuhou gật gù, nhưng khi đọc kỹ hơn trên bao bì, cô trố mắt ra mà ngã ngửa: ``Aqua-onee-san! Cái này là \emph{sữa tươi} chứ có phải \emph{kem tươi} đâu!?''

Nghe vậy, Shion, người đang trộn bột lại từ đống kem mà Luna đã đánh xong ban nãy, ngoảnh lại ngay, cười nhếch mép: ``Biết ngay mà! Dám chắc là bà cũng làm hỏng luôn cả kem ngô luôn sao?''

Aqua gục đầu xuống bàn, giọng yếu ớt: ``\dots Cái đó thì tôi xay nguyên bắp ngô với kem sữa rồi\dots''

Tất cả mọi người đều toát mồ hôi hột với lời thú tội này. ``Đúng là hậu đậu mà.''

``Đâu nào, đâu nào? Để em nếm thử xem!'' Noel lật đật mò tới, cầm thìa nếm thử lớp bột.

Nàng Hiệp si vừa múc một thìa kem từ Aqua, cho thẳng vào miệng\dots

*ỰC!*

Một giây. Hai giây. Ba giây.

Ngay sau đó, khuôn mặt cô biến sắc lại, lè lưỡi ra. ``N-Ngọt quá\dots\ Y như ăn hộp sữa trộn với cháo ngô vậy!''

\ct

Nhưng Aqua không phải là người duy nhất mắc lỗi. Cả Subaru và Luna cũng ``góp vui'' không kém.

Nàng Vịt - trong khi Aqua còn đang hí hoáy với bát kem trà xanh đang pha sai kia - thì phụ trách phần kem sô-cô-la. Trong một thoáng lơ đễnh, cô ngoái đầu hỏi nàng Kẹo: ``Luna-tan bỏ dâu tươi vào chưa?''

``Bỏ rùi đó-nora! Tận ba quả lận á!'' cô nàng vui tươi đáp lại, tay vẫn cầm chiếc khuấy quay đều.

``Ba quả thì có vẻ hơi ít nhỉ\dots''

Bỗng chốc, Luna tỏ vẻ rụt rè hơn hẳn. ``Cơ mà nè\dots''

``Sao?''

``Luna vẫn chưa bóc zỏ đâu nạ\dots''

Subaru chết cứng. Cô quay sang nhìn Luna đang cho nguyên quả dâu còn cuống và cả\dots\ túi ni-lông vào máy xay.

``Này, này, \textbf{này!!} Sao lại cho nguyên túi vào thế!?'' nàng Vịt đặt bát kem xuống, vội lao tới ngăn lại.

``Thì, chả phải nó sẽ làm cho kem ngon hơn xao-nora?''

``Em học cái logic đó ở đâu vậy!?''

Ryuuhou bên kia chỉ còn biết lắc đầu nhìn bốn ``chiến sĩ đánh kem'' xoay như chong chóng giữa hương liệu, máy đánh và những tiếng cãi lộn rất vụn vặt.

``Sau đó chúng ta sẽ lọc lại kem bằng rây đó! Nên là hai người không phải lo đâu!'' cô bé vọng ra, gỡ găng tay khỏi bát kem rồi quay về khuôn bánh.

\ct

Thậm chí cả Rushia cũng không thoát khỏi những tai nạn khó đỡ như vậy.

Nàng Pháp sư trong lúc làm kem vani đã vô tình đổ tinh dầu tận\dots\ bốn lần liều lượng so với công thức.

Ngay khi cô nàng bắt đầu trộn bát kem, một lượng lớn hương vani đậm đặc phả khắp cả căn phòng Aloe.

``Mùi gì nặng vậy\dots?'' Ryuuhou là người đầu tiên phát hiện có điều bất ổn.

``Ngửi như vani à\dots'' Shion đáp, lấy một tay bịt mũi lại.

``Giống như hít kẹo bông gòn pha nước hoa. vậy\dots'' Okayu tiếp lời.

Họ hướng về bát kem vani mà nàng Pháp sư đang trộn với vẻ mặt vui vẻ, hầu như không nhận điều gì đó khác biệt.

``Rushia-onee-san?'' cô em gái của Kuon lên tiếng.

``Sao vậy?''

``Chị cho bao nhiêu vani vào bát vậy?''

Cô đồng tóc xanh khẽ ngạc nhiên, tay vẫn đánh tăm tắp: ``Ể!? Ờm\dots\ bảy đến tám giọt gì đó?''

Tất cả mọi người nghe vậy liền đổ sạp xuống đất, cùng lúc đồng thanh: ``\textbf{Chả trách!}'' với những khuôn mặt cam chịu.

\ct

Sau 30 phút hì hục với những bát kem, khi tất cả nhìn lại năm bát kem mà bốn người đã trộn, họ thấy:

Một bát kem dâu hồng nhạt có cả cuống dâu và còn dáng dấp ni-lông.

Một bát sô-cô-la chưa đặc hẳn, bị vón lại thành nhiều cục lổn nhổn.

Một bát trà xanh loãng toẹt hơn cả bát cháo.

Một bát kem bắp trông như cháo nước pha sữa với vị quá ngọt.

Và một bát vani thơm lừng, đến mức ``chấn động tiềm thức.''

Aqua hít sâu một hơi, lắc đầu đầy thất vọng. Cpp giơ tay tuyên bố như một nữ tướng vừa quyết định tái lập đế chế: ``Chúng ta\dots\ cần làm lại đống mẻ này!''

Ngay sau câu nói đó là một loạt tiếng: ``\textbf{Ể!??}'' thấu vang cả trời xanh, tới mức họ cảm tưởng căn phòng này có rung chấn nhẹ.

``L-Làm lại á!?'' nàng Vịt là người sững sờ nhất. ``Tất cả chỗ này ư!?''

``E-Em đánh kem mỏi hết cả tay rồi-nanodesu\dots'' nàng Pháp sư thì bủn rủn đôi tay, mắt trố ra như không tin vào mắt mình.

``Vậy đống này sẽ xử lý thế nào đây?'' Mel - người đang hoàn thành công tác chuẩn bị lò - đưa ánh mắt lo lắng hỏi.

``Chắc phải đổ hết thôi chứ làm sao nữa\dots'' nàng Hầu gái đáp lại, giọng như cạn kiệt năng lượng.

Ngay lúc đó, Ryuuhou - tay vẫn đang lau khuôn bánh - liền tiến tới, đứng chống hông đúng kiểu ``em gái trụ cột gia đình'', giọng nghiêm túc bất ngờ: ``Không làm lại gì hết! Gia đình nhà Hikawa ở đây không chấp nhận đổ mứa đồ ăn!''

Lời tuyên bố của cô bé khiến cho tắt cả mọi người sững sờ. Aqua định lên tiếng phản biện: ``N-Nhưng như thế này thì sửa lại làm sao?''

Cô em gái của Kuon liền đáp, giọng vẫn quả quyết: ``Ta vẫn sẽ luôn có cách để sửa chúng!''

``Vói cả ấy,'' Shion từ bàn bên cạnh cũng bước lại, mái tóc tím bạc bay nhẹ theo gió máy quạt, ánh mắt sắc như giáo sư môn Nghệ thuật Sống sót trong Bếp: ``nếu bọn mình bắt đầu lại từ đầu, sẽ không kịp làm xong ít nhất sáu mẻ kem đâu.''

Tất cảm mọi người đêu điếng người trước thông tin như sét đánh ngang tai: ``\textbf{T-Tận sáu mẻ!?}''

Ryuuhou gật đầu, hướng về đống lọ hương liệu: ``Ít nhất theo những gì mà chúng ta có ạ. Cũng tại Okayu-onee-san mua nhiều quá mà\dots''

Nàng Mèo lười biếng giải thích: ``Kế hoạch của chúng ta là sẽ làm một chiếc bánh nhiều tầng với toàn bộ các hương liệu ở đây, mỗi phần là một hương vị mừ\~{}''

Rushia cầm những chiếc lọ hương liệu lên, vừa đọc hương liệu vừa lẩm bẩm: ``Có nho, có dứa, có mâm xôi, có caramel\dots'' để rồi chính cô trố mắt: ``\textbf{Ể!? Hết chỗ này ư!?}''

``Đúng rồi á, meo\~{}''

Cả nhóm đánh kem nghe xong liền đứng hình. Rushia ngồi bệt xuống sàn, xém làm rơi mấy lọ hương liệu nếu không nhờ Mel tới đỡ.

Luna suýt thì ngất thật, hai tay lủng lẳng xuống dưới sàn, miệng than vãn: ``Nắm thế-nanora\~{}!? Luna hơm nàm nổi nữ đâu-!!''

Subaru lẩm bẩm: ``Đây là làm sinh nhật hay mở chuỗi buffet bánh kem vậy\dots?''

Choco nhìn chiếc đồng hồ treo tường, thở dài thay cho nhóm: ``Với tiến độ hiện tại thì chắc phải làm tới tận hai tiếng đấy\dots''

Nàng Phù thủy nhỏ thở dài, đặt cái rây xuống cái bát lần đó, cười nhếch: ``Không còn cách nào khác nhỉ.''

Cô búng tay một cái, ngay sau đó, một vòng tròn ma thuật tím ánh lên giữa lòng bàn tay cô. Làn khói mỏng xoáy quanh từng bát kem bị đánh hỏng, làm chúng bắt đầu bốc hơi nhẹ và biến đổi dần thành hỗn hợp mịn mượt hơn.

Kem trà xanh trở nên sánh hơn, kem va-ni bớt nồng, kem ngô không còn loãng như cháo, và dâu thì\dots\ ít nhất không còn thấy cuống hay bất cứ mảnh ni-lông nào.

Thế nhưng, cô chỉ có thể làm cho nó trông ``ăn được'', còn phần kết cấu chung vẫn còn vẫn hơi lỏng, tới mức hoàn toàn chưa thể sử dụng ngay được.

Ryuuhou liền lục trong ngăn tủ nguyên liệu ra vài món đặc biệt hơn: ``Choco-sensei, cô lấy cho em ít gelatin với bột ngô trên chạn ạ!''

Nàng Y tá lập tức hiểu ý, lấy ra vài gói làm thạch và một túi bột ngô từ túi nguyên liệu họ đã mua ở đó. ``Em định làm những gì mà cô nghĩ đấy chứ, Ryuuhou-sama?''

``Tất nhiên là vậy ạ! Em học mẹo đó từ bố em đấy!'' cô bé trả lời, nhẹ nhàng rắc gelatin và một chút bột ngô pha loãng vào từng bát, rồi hướng dẫn các cô gái idol dùng kỹ thuật khuấy tròn đều tay cho đến khi đạt độ sệt mong muốn.

Với sự phối hợp giữa ma thuật của Shion và kinh nghiệm thủ công của cô bé Hikawa, toàn bộ năm bát kem từ mẻ đầu tiên được cứu sống ngoạn mục tới mức ảo diệu.

Aqua thở phào nhẹ nhõm khi họ không cần phải làm lại từ đầu ``Vậy là vẫn có cơ hội à\dots''

Subaru xoa trán, lau mấy giọt lấm tấm mồ hôi: ``Cứ tưởng chúng ta phải làm lại tất cả chứ\dots''

Luna thì hí hửng quay lại bát kem dâu của mình: ``Chúng ta lại đực đánh típ-nannora\~{}!''

Rushia nhăn mặt nhẹ, cất giọng trầm hẳn vì kém vui: ``Còn năm mẻ nữa-nanodesu\dots\ Đến bao giờ mới xong đây\dots''

Ryuuhou dõng dạc: ``Cố lên! Nếu xong trước giờ Onii-san về, cả nhóm sẽ được thưởng bánh tầng ba trét kem vô hạn!''

Một sự im lặng ngắn ngủi, rồi tất cả hô vang: ``\textbf{ĐÁNH TIẾP!!!}''

Rất may mắn là sau đó, rút kinh nghiệm từ mẻ đầu tiên, nhóm bốn người họ đã có thể hoàn thành công việc đánh kem mà không gặp phải bất trắc gì nữa.

\ct

Bộ tam trộn bột - Okayu, Shion và Noel - cũng gây ra những rắc rối chẳng kém cạnh gì nhóm đánh kem.

Ngay khi nhóm đánh kem vừa tạm thời ổn định với mẻ đầu tiên được cứu sống thần kỳ, nhóm trộn bột của Okayu, Shion và Noel cũng bắt tay vào phần việc của mình.

Nàng Miêu kéo một cái bàn trống ra giữa phòng, trải tấm khăn lớn để tránh bột văng lung tung. Trong khi đó, nàng Phù thủy bê một thau bột mì to đùng mà Ayame đã mua, đặt phịch lên bàn, còn nàng Hiệp sĩ thì hí hửng lôi theo hai túi bột nở và gói đường lớn.

``Được rồi\~{}! Bột mì, đường, bột nở, bát kem\dots\ có đủ cả!'' Noel chắp tay lại, khơi dậy tinh thần của cả nhóm.

Shion chống nạnh, nhíu mày hỏi: ``Giò thì ai làm phần nào? Có cần chị lấy điện thoại ra xem công thức không?''

Nàng Hiệp sĩ hớn hở đáp: ``Không cần đâu chị! Em nhớ tỷ lệ này\~{} Hình như là cứ một phần bột mì, nửa phần đường, một ít bột nở, rồi trộn chung với kem đánh sẵn!''

``Nghe có mùi gì đó không ổn rồi đó nha\~{}'' nàng Miêu đáp lại với giọng lười biếng, tông giọng trông không có vẻ gì lo lắng cả.

Nghe câu trả lời từ hai người, nàng Phù thủy đảo mắt, tặc lưỡi nói: ``Thôi, không trông mong gì từ hai người rồi. Để đó chị đọc công thức vậy.''

Cô lấy điện thoại rồi, bắt đầu đọc công thức được ghi sẵn trên máy: ``Xem nào\dots\ ta cần 200 gam bột mì, trộn với 100 gam đường, và\dots''

Cứ thế, hai người con gái còn lại đều chỉ gật đầu lia lịa, tỏ vẻ hiểu được sơ qua cách trộn bột cho chiếc bánh.

Chỉ có điều, họ lại hiểu theo cách theo một cách không hề bình thường. Ryuuhou bằng một cách nào đó lại là người đầu tiên cảm nhận được sự bất thường từ họ. Cô bé chỉ có thể thở dài: ``\vie{\textit{Lại chuẩn bị báo nhà báo cửa nữa rồi\dots}}''

\ct

Okayu - tay cầm cái muôi to gấp đôi bát trộn mà không biết cô lấy nó bằng cách nào - xúc bột mì thẳng từ thau to ra\dots\ bằng cảm giác - hay nói cách bỗ bã của người châu Á, là lấy bằng cách ``tổ tiên mách bảo''. Không hề cân, không hề đong, mắt nhắm tịt lại.

Noel cũng chẳng kém cạnh gì, khi cô ``tiên phong'' đổ đường, với suy nghĩ rằng: ``Đường càng nhiều bánh càng ngọt càng ngon!'' tới độ đã đổ tới nửa túi.

Ryuuhou đang kiểm tra khuôn bánh gần đó, mặt nhăn lại, lầm bầm: ``\vie{\textit{Mình biết ngay mà.}}'' rồi gọi lớn: ``Đánh bột không phải cho lượng như vậy đâu các chị ơi!''

Nghe thấy tiếng gọi, Shion đang cắm cúi vào chiếc điện thoại tìm cách đánh bột thì tá hỏa nhìn thấy bột đã chất đống như núi Phú Sĩ.

``\textbf{Sao lại cho lắm thế!? Định nướng cho cả núi à!?}'' cô thét lên, làm cho mọi người giật mình.

Choco khẽ trố mắt ngạc nhiên: ``Cô biết chắc chắn đó không phải là liều lượng đúng cho một bát kem đâu. Ít nhất cũng phải đong chứ.''

Okayu gãi đầu, giả ngơ: ``Nhưng mà bọn em không có nên mới phải dùng cảm giác chứ\~{}''

Noel vỗ vai: ``Chả phải ta phải nướng cả núi bánh sao? Ta sẽ làm anh ấy bất ngờ mà!''

``Bất ngờ tới mức nhập viện vì tiểu đường đó\dots'' nàng Phù thủy kháy đểu, đảo mắt bất lực.

Sau cùng, cả nàng Miêu và nàng Hiệp sĩ phải san bớt số đường và bột mì đó đều cho năm bát, trong khi số bị thừa phải để trong một cái bát khác.

``Chỉ việc đong đếm thôi mà đã thế này rồi\dots'' cô em gái của Kuon thở dài, ``không biết liệu các bước sau có ổn áp không đây\dots''

\ct

Theo đúng kế hoạch ban đầu, Okayu nhận nhiệm vụ rây bột. Công việc chỉ đơn giản: rắc bột lên trên rây, rồi lắc nhẹ chiếc rây để bột rơi đều vào bát kem.

Nhưng, nàng Miêu lại không nghĩ vậy. Cô nàng cầm cái rây như đang rắc tiêu lên mì ramen, thi thoảng lại lắc một cách cực đại, tới mức có thể nghe thấy những tiếng lạch cạch làm cho mọi người chú ý.

``Thật là chán ngắt quá nha\~{}''

Rây được một chút thì chán, cô chuyển qua việc đổ thẳng cả túi bột vào bát kem luôn mà không cần tốn công gì.

Đã thế, Noel còn cao hứng cầm cái phới lồng\footnote{Dụng cụ đánh kem hoặc đánh trứng cầm tay. Từ này vốn dĩ là một từ Pháp: fouet và trông như cái lồng nên mới có cách gọi như vậy.} mà trộn bột với tràn đầy sự khí thế, không để ý rằng lượng bột đang bị quá nhiều.

Mel nhìn cảnh tượng đó, liền lên tiếng: ``T-Trộn bột chậm thôi! Đừng có làm bột xẹp--''

Nhưng rồi, quá đã muộn.

Làn bụi bột mỏng manh bay lên mù mịt như sương mù, khiến mọi người ho sặc sụa. Cả phòng chìm trong một biển bột trắng, vấy bẩn lên khắp mọi nơi, lan ra hẳn cả khu bếp.

Aqua từ nhóm đánh kem bên kia thò đầu sang, vừa lau mặt vừa hét: ``\textbf{Bên này còn đang đánh kem đó! Mồ! Đừng có biến phòng thành động tuyệt chứ!!}''

\ct

Sau màn trộn như vũ bão, bụi bột dần vơi đi, nhưng những trò nghịch phá của bộ hai vẫn không dừng lại. Noel còn cười hí hửng: ``Để em đảm bảo bột mịn nha, Shion-senpai!''

Nói rồi, cô nàng nhấc bát bột lên lắc mạnh.

``Thế nào, thế nào? Như vậy chẳng phải sẽ làm bột mịn và đều hơn sao?'' nàng Hiệp sĩ tươi cười nói, xóc chiếc bát kem như xóc lô tô.

``Không phải xóc kiểu thế!'' cả Choco và Shion đồng thanh. Nhưng nàng ``Não nho'' vẫn mặc kệ, xóc chiếc bát đầy nhiệt tình.

Và rồi\dots

*KÉT KÉT KÉT* *BỤP!*

Bát bột trượt khỏi tay Noel, bay thẳng vào người Shion.

Nàng Phù thủy xui xẻo trở thành một món ``kem sống'', nhem nhuốc kem từ đầu đến chân, chỉ trừ lại đúng đôi mắt vàng nâu và đỏ.

Tất cả mọi người nhìn cảnh tượng đó, người thì trố mắt, người thì bụm miệng cười, người thì thở dài ngán ngẩm.

``C-Chuyện quái gì thế?'' Ryuuhou thốt lên, suýt đánh rơi cả khuôn bánh.

``Cái mặt bà\dots\ *khì khì* Cái mặt bà kìa!'' Aqua thì cố nén tiếng cười thành tiếng của mình.

Mel thì chỉ có thể khẽ lắc đầu: ``Mấy đứa chỉ phá là giỏi thôi\dots''

Shion đứng im năm giây, rồi cất giọng rền rĩ: ``Tất cả là tại hai người đấy\dots\ Bà nhớ mặt hai đứa!''

Vậy là sau màn kết hợp bất ổn, nhóm trộn bột vẫn tạm thời gom lại được ba âu bột tương đối dùng được, chuẩn bị sẵn sàng cho bước tiếp theo.

Còn bản thân nàng Phù thủy thì phải đi gột bột dính khắp áo trước khi có thể tiếp tục đồng hành.

\ct

Lần này, tự tay Shion sẽ phải làm công việc này. Cô cẩn thận trộn bột mì, bột nở và một ít muối với nhau, theo đúng công thức phù hợp.

``Xem nào\dots\ Từng bày muối, từng này bột nở\dots'' cô nàng lẩm bẩm, mắt liếc qua liếc lại giữa cái bát kem và cái điện thoại của mình.

Noel ngó mặt vào, mặt rạng rỡ: ``Chị trông có vẻ cẩn thận nhỉ, Shion-senpai.''

``Tất nhiên. Trộn nguyên liệu khô chuẩn chỉnh là bước quyết định để bánh nở đẹp đó. Lỡ sai tỉ lệ một chút thôi là bánh xẹp lép luôn,'' Vừa nói, Shion vừa đong từng thìa bột nở và muối vào cái bát lớn, gõ nhẹ tay thành bát để nguyên liệu rơi đều xuống.

Okayu cũng ghé lại xem, mắt lim dim như đang mơ màng giữa tiết học gia chánh: ``Ồ\~{} Cứ như thể bà đang chế thuốc phép thuật vậy\~{}''

``Chuẩn. Pha mà sai công thức, cái bánh sẽ thành một quả bom đúng nghĩa trong lò đấy.''

Nói rồi, cô nàng cầm rây lên, nhẹ nhàng cho bột vào và lắc đều tay.

Một làn mưa bột trắng mịn từ từ rơi xuống, tạo nên những lớp bột mềm tơi xốp, tách rời hoàn hảo.

Dù đây là lần đầu cô làm việc này và có phần hơi vụng, nhưng so với những gì mà nàng Hiệp sĩ và nàng Miêu đã làm trước đó, cô nàng đã làm tốt hơn rất nhiều, trước mắt là không gây ra thêm bất kỳ thảm họa nào nữa.

``Tại sao lại phải rây bột vậy, Shion-senpai?'' Noel nghiêng đầu hỏi.

Shion vừa lắc rây vừa giải thích: ``Bột mì hay bị vón cục. Nếu không rây, lúc trộn với nguyên liệu ướt sẽ tạo thành mấy cục lợn cợn. Bánh nướng lên vừa xấu vừa dai.''

Nàng Miêu nhìn dòng bột mịn rơi xuống, cất giọng đùa: ``Mình mà không rây là bánh thành gạch lát sân nhà Kuon-senpai nhỉ?''

Tất cả mọi người - kể cả từ nhóm đánh kem và nhóm chuẩn bị lò - đều cùng bật cười với câu đùa vô thưởng vô phạt đó.

``Đúng là thế thật đó Okayu!'' Aqua đáp, giọng cao hứng. ``Thậm chí là có thể làm vật liệu cho Noel-chan tập đấm nữa chứ!''

``Ý của chị là sao chứ!'' Noel phản ứng lại, phồng má lên phụng phịu. Cô nàng dường như không thích bị trêu như vậy.

Như thể muốn thể hiện như là một tay làm bếp kinh nghiệp, Shion còn thêm một chi tiết tuy nhỏ nhưng quan trọng nữa: cô dùng thìa gỗ đảo đều hỗn hợp bột mì, muối, bột nở sau khi rây, đảm bảo mọi thứ phân bố đều trước khi mang đi trộn với nguyên liệu ướt.

``Kỹ tính thật đấy\dots'' cả Noel và Okayu nhìn thành phẩm mà nàng Phù thủy đã tự tay hoàn thành.

``Bánh ngon là phải làm được những điều nhỏ nhất đó,'' cô kết luận, ``Không có chỗ cho đường tắt đâu!''

Vừa nói, cô vừa đổ phần bột đã rây và trộn kỹ sang chiếc âu trộn lớn đặt giữa bàn, chuẩn bị cho bước tiếp theo: hòa quyện với phần nguyên liệu ướt.

Mọi thứ được chuẩn bị gọn gàng và cẩn trọng\dots\ ít nhất là cho đến khi Okayu và Noel nhúng tay vào.

``Bây giờ đến lượt hai đứa đấy,'' cô đưa ra trước mặt họ hai âu kem khác, một âu kem matcha, và một âu kem dâu. Cô nhướng mày ra hiệu: ``Trộn bột vào hai âu này, từng phần một. Nhẹ tay thôi đấy!''

Okayu hí hửng cầm cái phới lồng lên, trong khi Noel thè lưỡi liếm mép, trông như chuẩn bị xông pha trận địa.

``Okayu-senpai, chị trộn cái màu xanh, cho nó matcha!'' Noel sung sức đề nghị.

``Được thôi\~{} Em sẽ làm kem dâu,'' Okayu đáp, giọng lười biếng nhưng tay đã cầm cái muôi bột chực sẵn.

Từ phía xa, những con mắt quan ngại dần nổi lên khi đến lượt bộ đôi phá hoại thực hành.

Choco đang đứng gần đó kiểm tra lò nướng nhìn qua, cười nhẹ đầy lo âu: ``Liệu lần này có làm nổ tung bếp không đây?''

Mel thì ôm má, giọng nửa lo lắng pha nửa đùa: ``Đừng để biến thành thảm họa bột bay như nãy nữa nha\dots''

Ryuuhou, tay ôm bản ghi chú ghi ghi chép chép, cũng không quên nhắc nhở: ``Nhớ là rây với trộn nhẹ chứ không được quậy ào ào đâu. Nếu không là bột nó xẹp đó nha.''

``Rồi\~{}'' cả hai cô nàng đồng thanh đáp, dù thực sự giọng của họ trông không tỏ vẻ tin tưởng gì.

Nàng Mèo bắt đầu đổ từng muôi bột vào âu kem dâu, rồi nhẹ nhàng gập từ dưới đáy lên.

Động tác ban đầu khá chuẩn, nhưng chỉ sau ba lượt rây và trộn bột, cô nàng dường như có dấu hiệu\dots\ buồn chán.

``Giá mà Koro-san cũng ở dưới này làm với mình nhỉ\dots\ Chán quá nha\~{}''

Cô nhìn ngó xung quanh, dường như không ai để ý - mặc dù thực tế là vẫn đang quan tâm tới cô - rồi nhoẻn miệng cười gian: ``Để nhanh hơn thì quậy mạnh một tí cũng được nhỉ\~{}?''

``\hl{Biết ngay là kiểu gì cũng có chuyện mà\dots}'' nàng Dơi, nàng Y tá và cô em gái của cậu Tổng đồng loạt lẩm bẩm ngay khi nhìn thấy tia lóe trong con mắt của nàng Miêu.

Nói là làm, Okayu xoay tròn cái phới như đang đánh súp miso.

Bột bị trộn quá mạnh, các bọt khí trong kem bắt đầu vỡ tan, hỗn hợp trở nên đặc quánh lại, trông chẳng khác nào như cháo nếp.

Shion quay đầu lại, nhìn thấy khoảng bột mà nàng Miêu trộn dần trở nên cứng ngắc, liền to tiếng: ``\textbf{Trời ơi, Okayu! Bà đang làm gì vậy!?}''

``Trộn nhanh hơn một chút thôi mừ\~{}''

``Thế này thì bánh phẳng lì mất!!'' Mel vội chạy tới ngăn cô nàng Miêu đang hăng máu vô tư kia lại.

Choco thở dài, bĩu môi: ``Xem chừng em ấy không học được gì từ Shion-sama rồi.''

Ở phía bên kia chiến tuyến, Noel cũng chẳng khá khẩm hơn là bao, thậm chí có phần thảm họa hơn.

Ban đầu, cô nàng Hiệp sĩ trộn rất cần mẫn và cẩn thận, đúng như hướng dẫn của bậc tiền bối.

Nhưng mới chỉ được ba phút, cô bắt đầu sốt ruột khi thấy bát kem không đặc lại: ``Sao nó lâu đặc vậy?''

Và rồi, một ý tưởng điên rồ khác nảy ra trong đầu. ``Chắc có lẽ mình nên thêm bột cho đặc hơn nhỉ?''

Không cần chờ câu trả lời từ bất cứ ai, Noel lập tức xúc thêm một muôi bột mì thả vào âu kem trà xanh, rồi khuấy thêm một hồi.

``Có vẻ như nó đặc thêm rồi đấy!'' cô nàng hồ hởi khi thấy kem có dấu hiệu đông lại. ``Nhưng mà thế vẫn chưa đủ!''

Nói là làm, nàng Hiệp sĩ liền khuấy ào ào bằng những động tác nhanh và mạnh, như thể các công nhân đang trộn xi măng.

*CỘP CỘP!* *RẮC!*

Tiếng bột nở nứt vỡ dưới đáy âu, báo hiệu một điều chẳng lành. Ryuuhou thấy vậy, vội đặt bảng ghi chép xuống, lao tới giành cái bát kem: ``\textbf{Noel-onee-san!!} Chị trộn như thế là thành cháo luôn đó!!''

Shion cũng chạy tới giật lấy cái phới từ tay cô nàng: ``\textbf{Ta có làm bê tông đâu, đồ não nho này!!}''

Noel nhìn hai người, mặt ngơ ngác nhìn hai người họ: ``Ể? Em tưởng là chúng ta đổ bột vào rồi khuấy như thế chứ?''

``\textbf{Không phải là trộn kiểu thế!!!}'' cả hai cô bé cùng nhau đồng thanh, hét thẳng mặt nàng Hiệp sĩ.

Okayu lấm lét nhìn qua, cười khổ: ``Chị thề là chị làm còn nhẹ tay hơn Noel-chan rồi á\~{}''

\ct

Cả nhóm ngán ngẩm nhìn hai âu bột: một âu kem dâu xẹp lép như bánh mochi, còn một âu matcha thì đặc quánh như hồ xây nhà.

Nàng Phù thủy nghiêm nghị nhìn hai âu kem thảm họa, khuôn mặt tỏ vẻ rất không hài lòng: ``Tôi đã bảo rồi, không thể nào vội được đâu! Cứ quýnh quáng như thế thì chả làm được gì cả!''

Và rồi những lời vàng tiếng ngọc được thốt ra từ cô nàng xui xẻo nhất trong cả đội. Tiếng hai cô nàng Miêu và nàng Hiệp sĩ cùng lúc thốt lên: ``\textbf{Bọn tôi xin lỗi mà\~{}!!!}''

Chứng kiến cảnh vừa thương vừa buồn cười đó, Ryuuhou chỉ thở dài ngao ngán nói: ``Lần sau đừng bao giờ mời hai chị ấy vào bếp nữa.''

``Thôi nào Ryuuhou-chan,'' Mel lên tiếng an ủi, ``Chúng ta vẫn có thể chữa được mà, đúng chứ?''

``Vấn đề là em không biết cứu pha này thế nào nữa, Mel-onee-san ạ\dots''

``Để cô giải quyết cho.'' Dứt lời, Choco đi tới, vỗ vai Shion đang nóng tính kia: ``Rồi, rồi, Shion-sama. Ta vẫn có thể làm lại được mà, đúng không?''

Tới cuối cùng, Shion - một lần nữa - trở thành anh hùng bất đắc dĩ khi dùng ma thuật để chữa lại hai âu kem bị hỏng kia.

``Lần này để bọn chị trộn cùng nha,'' nàng Ma cà rồng nhỏ nhẹ nói, ``Không lại biến căn phòng của Ryuuhou-chan thành công trường thi công mất.''

``Lúc đó mà dọn dẹp thì ốm người đó-nanodesu,'' Rushia bổ sung.

Dưới sự chỉ dẫn kỹ càng hơn từ Shion, Choco và Mel, cả nhóm bắt đầu trộn lại mẻ bột mới. lần này, với tinh thần trân trọng từng hạt bột.

Mẻ bánh thứ hai bắt đầu khởi động trong tiếng cười khúc khích và những lời động viên ấm áp và nhẹ nhàng. Rất may, mẻ thứ hai, và các mẻ tiếp theo diễn ra đầy suôn sẻ và thận trọng, với một vài tai nạn nhỏ nhưng được giải quyết nhanh gọn.

\ct

Ngay sau khi nhóm đã được chia xong, nhóm chuẩn bị khuôn và lòng nướng - dưới sự chỉ dẫn của Shion và một vài sự góp tay của Mel và Choco - nhóm cuối cùng chuẩn bị khuôn và lò nướng - dẫn đầu bởi con gái nhà Hikawa cùng với Aki và Suisei cùng nàng Y tá và nàng Dơi - cũng khẩn trương bắt tay vào công việc.

Trên chiếc bàn lớn, đủ loại khuôn được bày ra: từ khuôn tròn, khuôn vuông, khuôn tim, khuôn sao truyền thống, cho đến cả khuôn hình mèo (mà Okayu đã mang đi từ nhà) và khuôn hình chó (Korone cũng ăn theo), khuôn hình quỷ (của Ayame).

Ryuuhou cố ngẫm lại những công đoạn mà cô cùng với gia đình làm bánh kem trước đó, rồi lên tiếng thận trọng: ``Trước tiên, chúng ta cần chọn khuôn bánh. Bởi vì chúng ta làm bánh nhiều tầng nên tầng dưới cùng phải lớn nhất, rồi giảm dần xuống.''

Mel cầm một cái khuôn tròn cỡ trung, tần ngần hỏi: ``Vậy mình dùng cái này cho tầng dưới cùng nha?''

Suisei lập tức lắc đầu quầy quậy: ``Thế thì không đủ đâu! Aqua-chan bảo tầng 1 phải to bằng một sải tay cơ!''

Cả nhóm đồng loạt ngoảnh lại nhìn Suisei như thể muốn nói: ``Sui-chan/Suisei-onee-san nói đùa à?''

Rồi, họ liếc về phía Aqua, người đang hướng dẫn cho ba tên ``đầu đất'' còn lại cách đánh kem, khuôn mặt không thể nào tin được.

Cô em gái của Kuon giơ hai tay sang ngang đo thử: ``Ý của chị ấy là một sải tay cỡ người lớn ạ? Nếu thế thì nó sẽ ngót nghét một mét đó, Suisei-onee-san.''

Thế rồi, cô bé thở dài: ``Em không nghĩ chúng ta có thể nhét vừa vào cái lò đâu ạ. Có khi phải nướng ở ngoài sân cơ.''

Aki phì cười: ``Q-Quả thật cái bánh ta định làm đúng là hơi quá sức nhỉ.''

Choco thở dài, cầm cái khuôn to nhất trong đống ra so thử: ``Kể cả dùng cái này, cũng chỉ bằng cỡ nửa sải tay thôi. Mà cái lò bọn mình đang có thì bé tí, chỉ đủ đút khuôn cỡ trung.''

``Khó ghê nhỉ\dots'' Ryuuhou thở dài, tỏ vẻ khó nhằn ngay từ công đoạn bước đầu.

Tưởng rằng họ sẽ phải bỏ cuộc ngay từ công đoạn đầu tiên, thì đúng lúc đó từ đằng sau, một luồng ánh sáng tím nhạt lóe lên.

``Cứ để đó cho tôi!'' Shion giơ tay lên cao, hô lớn như một vị anh hùng bất đắc dĩ.

Trước con mắt tròn xoe của mọi người, nàng Phù thủy vẽ một vòng phép rực sáng giữa không trung. Một cơn gió nhẹ quét qua phòng, những tia sáng ma thuật xoáy quanh cái khuôn bánh, rồi *phụt* - cái khuôn cỡ trung kia phình to ra gấp đôi, gấp ba, đến mức lấp gần nửa căn phòng.

``E-Em làm cái gì vậy, Shion-sama!?'' nàng Y tá hoảng hồn khi nhìn chiếc khuôn được hóa lớn.

``Thế này là chật quá rồi đó!'' Ryuuhou cũng tá hỏa, vội trèo lên một chiếc ghế xoay gần đó.

``Có gì đâu,'' nàng Phù thủy nhún vai, ``tại tôi thấy mọi người than vãn là khuôn nhỏ nên mới hóa lớn thôi.''

Chưa dừng lại ở đó, cô nàng còn không quên quay ra phía lò nướng, vẽ thêm một vòng tròn ma thuật khác.

Tức thì, cái lò nướng vốn chỉ to bằng cái tủ nhỏ, giờ đã phình ra chạm tới tận trần nhà, trông như một cỗ máy thời tiền sử trong phim viễn tưởng.

*CỘC!*

Và có lẽ bởi vì nó hóa hơi to, một vết nứt trên trần nhà xuất hiện, đúng lúc mà chiếc lò nướng ngừng lại.

Trông chiếc lò nướng giờ đã chiếm một nửa căn bếp, Suisei tái mét, suýt làm rơi khuôn bánh khỏi tay: ``Ể!?? Lò nướng thành lò sưởi kìa!!''

Aki cũng hoảng loạn khi nhìn thấy những vết nứt đang loan lổ, cùng với đó là một ít bụi rơi ra: ``Chết rồi! Phòng của Ryuuhou-chan!''

``Cái lò nướng đó\dots\ chắc phải đủ cho cả khu phố nhà Hikawa ăn đấy\dots'' Choco ngẩn ra, hơi ngạc nhiên.

``Quả đúng là cô nàng Phù thủy nhỏ Gen 2 có khác\dots'' Mel run rẩy. ``Cơ mà liệu như vậy có hơi quá không?''

Nhưng, Ryuuhou lại là người lo lắng nhất. Cô lẩm bẩm: ``\hl{Aloe-chan mà thấy được cảnh tượng này thể nào cũng la om sòm cho xem\dots}''

Cô thở dài một hơi, ôm đầu bất lực: ``Với cả, Shion-onee-san à, thực ra chị không cần phải làm thế đâu\dots''

``Ủa, gì nữa? Chẳng phải mọi người muốn làm bánh siêu to khổng lồ ư?'' Shion nghiên đầu thắc mắc.

Ryuuhou cố trấn an lại tinh thần, rồi từ tốn giải thích: ``Em nghĩ là\dots\ ta cứ nướng bánh như bình thường với kích thước vừa phải, rồi khi nướng xong ấy, nếu muốn làm bánh khổng lồ thì dùng phép hóa lớn thôi ạ.''

Cả căn phòng im lặng một hồi lâu, các nhóm khác cũng dừng tay lại, ngoái nhìn về phía cô em gái của cậu Tổng.

``Ừ nhỉ! Sao ta không nghĩ ra sóm!'' nàng Phù thủy búng tay, gãi đầu cười ngượng. ``Tự dưng tốn ma lực chỉ vì mấy chuyện như thế này\dots''

``Ryuuhou-chan đúng là thông minh thật đó!'' Suisei xoa đầu cô em gái, đáp lại cũng chỉ là một tiếng thở dài.

Aki gật gù tán thành: ``Lại còn tiết kiệm nguyên liệu nữa. Bánh vừa đủ ăn, khỏi lo dư mứa.''

``Cơ mà,'' Subaru lên tiếng, chỉ về phía đống nguyên liệu họ mang xuống, ``chỗ này ta sẽ giải quyết như thế nào?''

``Chị cũng nhỡ tay mua hơi lố rồi\dots'' Okayu lười biếng nói, khuôn mặt vẫn tỏ vẻ điềm tĩnh.

``Em nghĩ là chúng ta cứ mang về dùng dần thôi ạ,'' Ryuuhou nói, ``đằng nào các chị cũng chưa định bóc tất cả ra chứ?''

Tất cả các cô gái đều gật gù với ý tưởng đơn giản này.

``Đúng là ý tưởng như một cô bé tiết kiệm có khác\~{}'' Choco vỗ tay. ``Chị tán dương tinh thần này của em đấy.''

``Đâu cần nhất thiết phải lãng phí như thế đâu nhỉ?'' Mel cười mỉm, khẽ trêu Shion. Thấy thế, cô nàng Phù thủy nhỏ phản ứng lại, cố gắng ra vẻ điềm đạm bằng cách vuốt tóc: ``À thì\dots\ Tại em quá hăng hái thôi mà! N-Nhưng mà mọi nếu mọi người muốn như vậy thì được thôi!''

Nói rồi, cô bé vung nhẹ tay, hóa giải phép thuật. Cái lò và cái khuôn khổng lồ từ từ xẹp xuống trở về kích cỡ bình thường, chỉ để lại trong không khí một lớp bụi phép lấp lánh như sương mù mỏng.

\dots Và cả những gì còn sót lại: hai vết nứt lớn trên trần nhà và những vết xước của chiếc khuôn bánh khổng lồ ban nãy.

Cô em gái của Kuon cũng chỉ có thể cười đầy chua chát, thầm nghĩ: ``\vie{\textit{Thực sự thì mình sẽ phải nhờ Onii-san và bố mẹ dọn chỗ này trước khi chị ấy về đấy\dots}}''

\ct

``À mà này\dots''

``Sao thế, Aki-onee-san?''

``Căn bếp phòng em chỉ có mỗi một cái lò thôi, mà bây giờ ta phải làm nhiều tầng thì\dots''

``Liệu có kịp không?'' Suisei đặt vấn đề. ``Chị biết là ta có thể làm như bình thường, tức là nướng từng mẻ một, nhưng\dots''

``Nếu ta có thêm một cái lò nướng nữa thì tốt nhỉ?'' Mel chốt lại.

Ryuuhou nhìn cái lò bánh của phòng Aloe. Cô bé xem xét sơ lược một hồi, ước tính chiếc bánh họ định làm, rồi đáp: ``Em thấy cái lò đó đủ to mà.''

Thế rồi, cô nói tiếp: ``Nhưng mà nếu các chị muốn, nhà em còn một cái lò nướng nữa ở trên căn bếp nhà em ạ.''

Trước thông tin này, các thành viên trong nhóm nướng cảm thấy an tâm phần nào. ``May quá. Vậy là ta có hai lò chạy song song rồi,'' Mel hồ hởi nói.

``Nhưng mà nó cũng hơi nặng đấy ạ\dots''

Ngay khi cô bé dứt câu, một giọng nói từ đằng sau cất lên: ``Không sao! Việc này cứ để Danchou lo!''

Noel - người đang chờ được làm bột từ nhóm đánh kem, hí hửng đi tới. ``Để chị lo việc khuôn vác!''

``Ồ, cảm ơn chị, Noel-onee-san.''

Nàng Hiệp sĩ sau đó đã một mình vác chiếc lò nướng từ căn bếp của nhà Hikawa xuống mà không hề gặp phải bất trắc gì.

\ct

Tạm giải quyết xong vụ lò, nhóm lại gặp phải vấn đề khác cũng quan trọng không kém: cách chia các phần bánh kem.

Theo gợi ý của Aqua - mà thực chất là lấy từ ý tưởng điên rồ của bộ đôi Chó-Mèo - mỗi tầng sẽ được chia làm nhiều hương vị, ứng với mỗi mẻ mà họ sẽ dự định làm.

Câu hỏi lớn nhất là: làm thế nào để họ làm được điều đó?

Choco chỉ vào các túi nguyên liệu xếp la liệt: ``Có cách nào để ta có thể chia cho gọn mà vẫn đều các vị nhỉ?''

Mel nhìn quanh, tay chống cằm suy nghĩ: ``Chia khuôn nhỏ?''

Suisei thấy thế, liền lắc đầu: ``Chắc chắn là không đủ đâu. Nghe đồn bảo cái bánh mà Aqua-chan phải làm ít nhất chục vị cơ\dots''

Aki sáng mắt lên, đề xuất: ``Hay là ta dùng mấy cái vật dụng kim loại mình có trong bếp như khung inox, vòng bánh, hay mấy cái bát to nhỏ khác nhau?''

``Cụ thể là thế nào?'' nàng Ca sĩ và nàng Dị giới nghiêng đầu.

Ryuuhou lấy ra những vật dụng nhỏ từ trong chạn tủ, rồi giải thích: ``Em nghĩ là ta có thể dùng những thứ dĩa hay thìa để làm vách ngăn tạm, rồi khi ta có đổ bột vào khuôn to thì ta cứ căn chia tỷ lệ vào mấy cái phần nhỏ hơn vào thôi ạ.''

Đề xuất của cô liền được cả nhóm đồng tình. Nàng Ca sĩ thậm chí còn hóm hỉnh: ``Em gái của Kuon-senpai có khác, đúng là tinh ý thật!''

\ct

Sau khi đã chọn được căn chỉnh tỷ lệ vừa vặn trong các khuôn bánh, Suisei và Aki - một người hầu như chưa vào bếp và một người có chút kỹ năng nấu nướng - được giao nhiệm vụ chuẩn bị khuôn chống dính.

``Đơn giản thôi ấy mà!'' nàng Ca sĩ hào hứng tuyên bố, tay cầm cây cọ phết bơ, lẩy phẩy trên không khí.

Nàng Dị giới thì thủ sẵn một cuộn giấy nến to vật vã, mắt sáng lấp lánh: ``Vậy em cắt giấy lót đáy nha!''

Hai cô nàng bắt đầu thực hiện công việc được coi là ``dễ xơi'' nhất trong công đoạn làm bánh này. Người cầm cọ bơ, người cầm tấm giấy bạc, họ ngân nga những bài hát quen thuộc, như thể đang ca bài ca lao động.

Mới đầu, mọi chuyện dường như diễn ra xuôi chèo mát mái, không hề xảy ra một sự kiện khó đỡ nào. Mọi thứ vẫn diễn ra một cách nhẹ nhàng nhưng khá khẩn trương.

Để rồi, chính sự khẩn trương đó mà nàng ``Sao chổi'' đã có phần nóng vội. Cô nhiệt tình đẫm chiếc cọ xuống lát bơ được làm mềm, rồi lau cọ chiếc khuôn như không hề có gì xảy ra.

``Phết bơ nhiều để cho bánh dễ lấy hơn\~{}'' cô nàng vừa ngân nga vừa tiện tay\dots\ cho cả thìa bơ vào khuôn, rồi lấy cọ phết một lớp dày cộp như đang trát vữa.

``Sui-chan phết nhiều bơ ghê\~{}'' Aki cười tươi.

``Càng nhiều bơ thì bánh sẽ dễ lấy hơn mà, đúng chứ?''

Ryuuhou từ phía xa, nhìn chiếc khuôn đẫm màu vàng nhớp nháp đó, chỉ có thể cười chát: ``Không cần phải phết dày đến thế đâu Suisei-onee-san\dots''

Nàng Tóc bay cũng không kém cạnh hơn: cô ước lượng giấy bạc để nướng lại quá ngẫu hứng: miếng thì bé tí lọt thỏm trong khuôn, miếng thì to tới mức tràn ra ngoài.

``Như thế là đủ chưa nhỉ?'' Aki cầm xấp giấy giờ đã mỏng còn một phần ba, thắc mắc.

``Ừm, chăc là vậy rồi đấy,'' Suisei đáp, nhìn thành phẩm mà họ vừa làm.

Choco từ xa liếc mắt qua, nhíu mày: ``Ủa cái khuôn đó\dots\ Hai đứa định làm bánh hay bọc gói quà vậy?''

Cả hai cô gái nghe tới đây cũng chỉ cười trừ, cố giương đôi mắt cảm thông: ``Biết làm sao được, Choco-sensei\dots\ Đây là lần đầu bọn em làm bánh mà.''

Mel bước tới, chỉnh lại cách đúng: ``Quét bơ mỏng thôi nha Sui-chan\~{} Còn Aki-chan, giấy thì cắt đúng đáy khuôn, viền vừa tầm, không thì bánh nở bị méo á.''

Sau một hồi vật lộn, khuôn cũng dần được chuẩn bị đâu ra đó, dù khuôn nào khuôn nấy cũng lem nhem một ít dấu vết ``nỗ lực hết mình.''

\ct

Ryuuhou, Choco và Mel làm công đoạn chuẩn bị lò nướng. Cô em gái sẽ lo phần lò phụ do quen tay hơn, trong khi hai người còn lại sẽ lo lò chính.

Với cô bé nhà Hikawa, cô không có vấn đề gì với chiếc lò nhỏ của gia đình. ``\vie{Xem nào\dots{} Hình như là đặt 160\degree C trong 25 phút\dots}''

Ở bên kia, hai cô nàng thạo bếp cũng bắt tay vào công việc. Choco lướt qua những núm chính, vừa xoay vừa căn chỉnh theo trí nhớ của cô: ``Cô nghĩ là chúng ta nên để\dots\ 170\degree C trong 60 phút nhé?''

Mel cũng gật đầu: ``Cho chắc! Em cũng toàn để một tiếng lận!''

Aqua ngoảnh đầu lại, chen vào: ``Một tiếng liệu có đủ không vậy, Choco-sensei?''

Nàng Y tá chống nạnh, tự tin đáp lời: ``Đương nhiên rồi! Cô dùng lò nướng nhiều rồi nên biết chứ!''

Trong lúc hai cô nàng thạo bếp vẫn đang kiểm tra những bước cuối cùng trước khi có thể rảnh tay hỗ trợ các nhóm khác, Ryuuhou - lúc này đang tính toán chia phần bột - nghĩ bụng: ``\vie{\textit{Nghe đồn là hai chị ấy vào bếp cũng nhiều\dots\ Nên chắc là cũng không sao đâu nhỉ?}}''

Chẳng ngờ, đó lại là một suy nghĩ cực sai lầm.

Một giờ sau, khi bột và kem cơ bản đã hoàn thành các mẻ đầu tiên và chuẩn bị được đưa vào lò nướng, cả Mel và Choco đều cùng hí hửng mở chiếc lò nóng sẵn ra.

``Để xem chiếc lò này được chưa nào\~{}'' Choco phấn khích, mong chờ cho bước tiếp theo.

Và ngay khi chiếc lò được mở\dots

*PHỪNG!*

Một luồng hơi nóng hừng hực phả ra.

Lập tức, cả nàng Dơi và nàng Y tá đều vội lùi ra sau, rụt tay lại như thể bị bỏng dù mới chỉ đưa tay gần vào bên trong lò.

``\textbf{N-Nóng quá!!!}''

Suisei đang ngồi gọt dâu để trang trí chiếc bánh kem, nghe tiếng hét của hai người liền giật mình, suýt cắt vào tay: ``C-Chuyện gì vậy!?''

Okayu núp đằng sau bàn, bông đùa: ``Không lẽ lò của nhà anh ấy lại là lò siêu cấp\~{}?''

Aki khịt mũi, rồi bịt miệng lại: ``Chị còn nghỉ thấy mùi khen khét kìa\dots''

Tất cả mọi người đều toát mồ hôi hột với chiếc lò nóng hổi kia, không một ai dám tiến lại gần.

Ngoại trừ cô bé nhà Hikawa. Cô vội vã chạy lại, nhìn vào nhiệt kế của chiếc lò nướng chính.

Hơn 230\degree C.

Ryuuhou thỞ dài ngán ngẩm, rồi lo lắng hỏi hai người: ``Choco-sensei, Mel-onee-san\dots\ Hai chị quen dùng lò 100V đúng không ạ?''

``Ý em là sao?'' nàng Y tá chớp mắt, ngạc nhiên với câu hỏi bất ngờ đó.

``Chị tưởng là ở đâu mà chả dùng 100V?'' nàng Ma cà rồng càng sững sờ hơn với lời nói tưởng chừng như hiển nhiên đó.

Cô em gái của cậu chỉ vào nhiệt kế, đáp với giọng lạnh tanh: ``Trời ạ. Chẳng trách. Ở Etsunan bọn em, dòng điện 220V mới là điện dân dụng ở đây.''

Cả nhóm ngỡ ngàng với lời giải thích của cô. Nhưng chưa kịp để họ nói bất cứ điều gì, cô tiếp tục: ``Thế nên các chị mà để một giờ thì kiểu gì cũng cháy bánh!''

Choco nhìn lại chiếc lò đang hầm hập kia, thở phào nhẹ nhõm: ``May mà chúng ta chưa cho bánh vào đấy\dots''

``Cũng phải. Nếu như vậy thì có khi chiếc bánh đã thành than mất rồi,'' Subaru toát mồ hôi hột.

``Có lẽ bọn chị nên hỏi kỹ trước khi dùng lò ở đây quá\dots'' Mel cười ngượng.

Sau đó, họ đã phải để lò nguội bớt một chút, trước khi có thể cho mẻ bánh đầu tiên. Cả nhóm lắc đầu cười xòa, rồi cùng nhau chuẩn bị cho công đoạn tiếp theo trong việc làm bánh.

\ct

Aki lấy ra những chiếc khuôn bánh đã được phết bơ và bọc giấy sẵn. Cô nhìn đống âu bột từ hai mẻ đầu tiên, giơ sáu ngón tay lên: ``Ta sẽ chia từng này bột thành sáu âu nha!''

Choco gật đầu, chuẩn bị sẵn cái cân điện tử mà Ryuuhou mang ra: ``Cô hiểu rồi. Tầng dưới cùng sẽ là tầng lớn nhất, càng lên cao càng nhỏ dần.''

``Như thể là\dots\ tháp Tokyo sao?''

``Ừm. Có thể hiểu là như vậy, Aki-sama.''

Hai cô gái bắt đầu rót từng âu bột vào cốc đong, canh lượng chính xác một cách tỉ mỉ, rồi sau đó đưa cho Suisei - người cũng đang tò mò - để đổ khuôn và căn chỉnh.

``\textit{\vie{Mọi người chỉ cần dùng cảm giác để định lượng thôi mà\dots}}'' Ryuuhou lẩm bẩm, nhưng vẫn tiếp tục nhìn mọi người.

Choco - vốn là một cô nàng Y tá trường kiêm cả giáo viên bộ môn - tỏ ra rất cẩn thận, thậm chí đo chuẩn từng gam một.

``Tầng đầu tiên cần từng này gam bột\dots\ sau đó từng này gam kem\dots'' cô đưa cho hai nhóm còn lại, tính toán một cách kỹ lưỡng.

Nhưng, Aki lại tính toán một cách cảm tính hơn - đúng theo tinh thần của người châu Á, chỉ khác là theo một cách rất ẩu đoảng.

Cô nhìn cái khuôn chỉ bằng mắt, rồi không hề do dự gì mà\dots\ đổ nguyên một âu bột vào bát để cân đong, thậm chí còn không thèm nhìn vào đồng hồ điện tử.

``Trông có vẻ ổn rồi đấy!'' Aki hí hửng nhìn bát đong, với phần bột bị tràn lên đến tận ngọn. Cô em gái của Kuon nhìn thành phẩm của cô nàng, liền toát mồ hôi hột: ``\vie{\textit{Không phải định lượng kiểu đó\dots}}''

Hiển nhiên Choco cau mày khi nhìn đồng đội của cô đong bột một cách bất cẩn. Cô gắt nhẹ: ``Ít nhất cũng phải đong cho nghiêm túc chứ, Aki-sama!''

Aki lè lưỡi, tỏ vẻ như bị bắt tại trận: ``Bị phát hiện rồi ha\~{}''

\ct

Sau khi bột được chia đủ sáu phần, Mel và Suisei nhanh chóng bắt tay đổ vào khuôn.

Nàng ``Sao chổi'' cẩm thận cầm bát đổ từ từ, nhưng tay cô lại đang\dots\ run lẩy bẩy vì hồi hộp. Chẳng thể trách được cô nàng khi đây là lần đầu cô làm bánh.

``Eek\dots! T-Tay chị run quá!'' cô rên lên khe khẽ, cố gắng nhắm vào tâm khuôn.

``Thư giãn nào, Sui-chan!'' Mel lên tiếng trấn an - ``Đừng có căng quá như thế, thả lỏng người ra nào!''

``N-Nhưng mà chị sợ bánh méo lắm!'' hai tay của Suisei càng lúc càng rung mạnh hơn, như thể cô đang vung tay lên xuống chứ không đơn thuần là đổ bột nữa.

Thế là, phần bột vani bị lệch hẳn sang bên trái khuôn, thậm chí còn lan hẳn sang các phần vốn sẽ đổ các phần bột hương vị khác.

``Ái chà chà\dots'' Luna cười chát với thành quả mà nàng Ca sĩ vừa tạo. ``Hơm xao đâu, Suisei-senpai! Lần đù ai cũng có xai xót mà-nora!''

``Hy vọng là Mel-senpai sẽ tốt hơn\dots'' Subaru liếc về nàng Ma cà rồng, người đang chuẩn bị đổ bột vào các khuôn nhỏ hơn.

Chẳng ai ngờ, ngay khi nàng Vịt vừa dứt câu, một thảm họa ẩm thực khác lại vừa xảy ra: Mel đổ bột như rót nước vào xô do quá tay.

``Nghe đồn bánh có lõi đang là trend đó, Mel-senpai\~{}'' Okayu đang mài thanh sô-cô-la, nói góp vào bằng giọng nói bình thản quen thuộc.

``Danchou cũng muốn thử món đó!'' Noel hớn hở xen ngang. ``Xin chị đấy, xin chị đấy, xin chị đấy-''

``Thôi đi, Noel,'' Rushia lên tiếng chặn đà hưng phấn của đồng đội, ``ta đang làm cho Kuon-senpai, không phải cho bà đâu-nanodesu.''

Ryuuhou nhìn ``chiến lợi phẩm'' của hai cô nàng mà chỉ biết ôm đầu thất vọng: ``\vie{Anh ấy mà biết được chắc là xỉu lên xỉu xuống luôn đấy.}''

Sau vài lần chỉnh sửa, cùng với một chút hướng dẫn từ các thành viên khác và tài nguyên nguyên cứu - mà lấy từ sách của Aqua và bài hướng dẫn trên mạng - cuối cùng cả sáu khuôn bánh cũng được lấp đầy bởi toàn bộ loại bột mà họ đã làm trước đó.

Một bảng mày bột sống trông đẹp mắt và rực rỡ, nhưng đi kèm với đó cũng là một sự rối rắm và có phần khá\dots\ màu mè.

\ct

Cô em gái của Kuon - người chịu trách nhiệm canh cả hai lò sau khi Choco và Mel để nóng lò quá mức trước đó - cho hai tầng lớn nhất của chiếc bánh vào trong lò nướng.

``Được rồi, 140\degree C đặt trong 30 phút\dots'' cô xoay núm của chiếc lò nhỏ một cách đầy dứt khoát, thao tác nhanh nhẹn.

Tiếp đến, cô đặt bốn tầng còn lại vào chiếc lò lớn và căn chỉnh thời gian cũng như nhiệt độ tương tự.

``Giờ ta chỉ còn việc chờ bánh thôi!'' Ryuuhou tháo bao tay chống nhiệt ra, tiến về chiếc bàn bày sẵn hoa quả và đồ ngọt. ``Giờ chỉ còn việc trang trí bánh nữa là xong!''

Aqua nhìn cô bé với một ánh mắt ngưỡng mộ, tò mò hỏi: ``Ryuuhou-chan có vẻ rành việc này lắm ha! Em làm nhiều lần rồi à?''

Tới đây, con gái nhà Hikawa bỗng khựng lại vì giật mình. Cô mím môi, gãi gãi đầu rồi cười khì đáp: ``Thật ra\dots\ đây mới là lần thứ ba em làm thôi ạ\dots''

Lời thú nhận của cô khiến cho cả căn phòng dường như nín lại.

Một giây, hai giây\dots

``\textbf{Ể!?/Ý em là sao chứ!?}'' tất cả mọi gười từ ba nhóm cùng đồng loạt thốt lên đầy ngỡ ngàng.

``À thì\dots'' cô nàng ngụng nghịu nói, ``lần gần nhất em làm việc này từ tận chín tháng trước, hồi sinh nhật bố chúng em cơ. Sau này vì gia đình em bận quá, cũng không có dịp trổ tài nên\dots''

``Vậy tại sao em điều phối tự tin vậy!?'' Suisei thắc mắc, vẫn còn hơi choáng.

``Nói ra thì các chị có lẽ sẽ không tin, nhưng,'' Ryuuhou giơ hai ngón tay lên: ``em chỉ nhớ được đúng hai thứ thôi ạ.''

``Và chúng là\dots?'' Subaru nhíu mày.

``Thứ nhất là nhiệt độ lò phải luôn ổn định, và thứ hai - cái này quan trọng hơn - là luôn làm nóng lò trước khi nướng ạ,'' cô bé đáp với giọng bình tĩnh đến lạ.

Bầu không khí trong căn bếp lại im bặt một hồi, không biết nên xử trí như thế nào với câu trả lời của cô.

``Ít nhất thì Ryuuhousa-chan cũng nhớ đực căn bản rùi còn gì-nora,'' mặt của cô tái mét lại, không biết vì áp lực từ bầu không khí hay vì sợ sệt trước sự bình thản của cô em gái của cậu Tổng.

Suisei giơ ngón tay cái, chữa thẹn: ``Vậy là hơn các chị rồi còn gì. Nhớ như vậy là tốt rồi!''

``Và em còn biết\dots\ không nướng 1 tiếng nữa cơ!'' Shion châm chọc, nhìn đểu hai cô nàng Y tá và Dơi.

Mọi người phá lên cười, còn Ryuuhou thì đỏ mặt xua tay: ``Thôi mà, đừng trêu em nữa! Đ-Để em canh lò cho mọi người!''

Sau khoảng gần nửa tiếng, tiếng *TINH* của chiếc lò vang lên, cũng là lúc cả sáu chiếc bánh đã được nướng chín vàng rộm. Những mẻ bánh đầy màu sắc ngon lành đã được bưng ra, phả một mùi hương ngào ngạt của bơ, sữa và các phẩm phụ.

Chiếc bánh sinh nhật cũng đã dần thành hình.

\ct

Và giờ, công đoạn cuối cùng cho chiếc bánh sinh nhật dần thành hình này: ``khoác áo'' cho chiếc bánh này bằng hoa quả và kem.

Shion, Rushia và Noel - tập trung trước chiếc bàn dài nơi các tầng bánh đã nguội bớt, xếp gọn gàng theo thứ tự kích thước.

Nàng Phù thủy đẩy gọng kính (tưởng tượng) của cô lên sống mũi, tiếp tục thuyết trình: ``Giờ là bước quan trọng nhất. Hai đứa cẩn thận đấy, trượt cái bánh là đi tong ngay.''

``Xem chừng có vẻ rất cẩn thận đấy-nanodesu\dots'' Rushia nhăn mặt, toát mồ hôi hột.

Tầng đầu tiên - cũng là tầng nền gồm các kem từ mẻ thứ nhất - được đặt một cách dễ dàng trên một tấm đĩa xoay.

Mặc dù nó khá lớn và rộng (chưa hóa lớn mà đã rộng tới mười lăm xăng-ti), nhưng với sự phối hợp khéo léo của cả ba cô gái, tầng lớn nhất đã được đặt ngay ngắn mà không có miếng nào bị vỡ.

``Thành công bước đầu!'' Noel giơ tay chữ V, hớn hở reo lên.

``Chúng ta vẫn còn năm tầng lận nữa đó, Noel\dots'' Rushia nhăn mặt. ``Các tầng sau có khi còn căng thẳng hơn đấy.''

Tầng thứ hai của chiếc bánh cũng bắt đầu triển khai. Tầng bánh này nhỏ hơn một chút so với tầng nền, được chia thành năm phần với năm hương vị khác nhau. Nó cũng nhẹ hơn đôi chút so với tầng một, nhưng không có nghĩa là sẽ dễ thở hơn.

Lần này, Rushia và Noel bê chiếc bánh lên, thận trọng đặt lên trên chính của tầng nền.

``Cẩn thận nha Noel\dots''

``Mình biết rồi!''

Có lẽ bởi vì quá cao hứng, nàng Hiệp sĩ đã hơi quá tay mà ấn chiếc bánh lệch hẳn sang một bên.

``N-Nó đang nghiêng kìa-nanodesu!!'' nàng Pháp sư hốt hoảng. ``Đừng có ấn mạnh xuống như vậy chứ!!''

Noel giật mình, vội nhấc nhẹ chiếc bánh lên. Cô quýnh quáng nhấc bánh lên, vô tình làm cho chiếc bánh bị vụn lại một mảng.

Thấy thế, Shion nhanh tay chữa lại chiếc bánh. Tầng thứ hai được đưa về dạng nguyên mẫu, nhưng đổi lại là một con mát ngờ vực từ chính cô nàng Phù thủy nhỏ đó.

``Đúng là PON mà\dots'' cả Subaru và Ryuuhou đều thở dài chán chường.

Các tầng tiếp theo, tầng ba, tầng bốn và tầng năm, cả ba người họ cũng gặp phải những sự cố nhỏ khác.

Lúc thì kem bị tràn vì ấn quá mạnh\dots

``Noel!! Đã bảo là đừng có ấn quá mạnh mà-nanodesu!!'' Rushia rít lên khi nhìn phần kem mà Shion đã phết bị tràn ra ngoài.

\dots lúc thì bánh bị xô lệch vì đặt không đúng tâm\dots

``Cẩn thận! Đừng để lệch nữa!'' Shion nhắc nhở, tay nhanh chóng chỉnh lại vị trí tầng bánh.

``Tại em run tay quá mà!'' Noel phân bua, mặt đỏ bừng vì ngượng.

``Run tay thì để đó tôi làm!'' Rushia chen vào, nhưng vừa nhấc tầng bánh lên thì một mảng kem dính vào tay cô, khiến cô hét lên: ``\textbf{Á! Kem dính hết vào tay rồi!}''

Cả nhóm bật cười, nhưng vẫn cố gắng phối hợp để hoàn thành việc xếp bánh.

\dots để rồi cuối cùng, tầng thứ sáu - tầng nhỏ nhất có đường kính bằng một chiếc bánh một phần ba kích cỡ bánh tiêu chuẩn, được chia làm bốn hương vị khác hẳn với năm tầng còn lại, và cũng là tầng đỉnh - được đặt lên một cách cẩn thận.

``Xong rồi!'' Nàng Phù thủy nhỏ thở phào, nhìn chiếc bánh sáu tầng đã hoàn chỉnh. ``Cơ mà\dots''

Tất cả mọi người nhìn sơ qua chiếc bánh kem không được hoàn hảo: có tầng thì hơi lệch tâm, có tầng thì bị trồi kem đôi chút, có tần thì hơi trũng xuống, nhưng mặt bằng chung dưới con mắt của những người làm bánh lần đầu, thì\dots

``Đẹp quớ-nora!'' Luna reo lên, mắt sáng rực khi nhìn chiếc bánh đầy màu sắc.

``Ít nhất cũng tạm ổn\dots'' Ryuuhou gật gù, chẹp miệng tạm chấp nhận với kết quả này.

``Và giờ thì\dots'' cô giơ bàn tay của mình lên, lập tức xuất hiện một vòng tròn ma thuật màu vàng.

``\textbf{Phóng lớn!}''

Một luồng sáng vàng dịu nhẹ tỏa ra bao phủ lấy chiếc bánh kem nhiều tầng. Trong chớp mắt, nó lớn dần lên - mỗi tầng phình ra đúng tỷ lệ, đường kính tầng nền giờ bằng gần một sải tay, đúng như mong muốn ban đầu của mọi người. Phần đỉnh bánh cũng ngấp nghé chạm tới trần nhà.

Okayu thốt lên từ xa, mắt cô mở to lần đầu: ``Ối trời\dots\ nhìn như lâu đài đường ấy\~{}''

Subaru huýt sáo:

``Nó\dots\ to bằng cái bàn ăn phòng Ryuuhou-chan rồi á!'' Rushia và Noel đứng bên cạnh nhìn thành quả mà nước mắt rưng rưng.

``Giờ chỉ còn trang trí thêm nến và hoa quả nữa là hoàn hảo!'' Ryuuhou nói, tay cầm sẵn con dao và các loại trái cây đã chuẩn bị từ trước.

``Hãy cùng biến nó thành một tác phẩm nghệ thuật nào, mọi người!'' Shion giơ nắm đấm lên cao, khích lệ tinh thần của mọi người.

``\textbf{Đồng ý!!!}'' cả thảy mười hai cô gái cùng nhau đồng thanh, nhanh chóng bắt tay vào công đoạn cuối cùng, trang trí chiếc bánh với tất cả sự tỉ mỉ và sáng tạo, để tạo nên một chiếc bánh sinh nhật thật đặc biệt dành cho cậu Tổng của họ.

\ct

Một lần nữa, đội hình trang trí chiếc bánh kem lại chia thành ba nhóm nhỏ: nhóm phết kem gồm Suisei, Luna, Subaru và Choco; nhóm cắt và gắn trái cây cùng các nguyên liệu khác gồm Okayu, Aqua, Rushia và Aki; và nhóm viết chữ là Shion, Noel, Mel cùng Ryuuhou - người thân cận nhất với Kuon.

Nàng Y tá và nàng Công chúa xứ Kẹo đảm nhận ở ba tầng trên cùng, trong khi nàng Ca sĩ và nàng Vịt xử lý ba tầng trệt và phụ trách xoay chiếc đế - vốn dĩ cũng được phóng to để phục vụ cho việc xoay bánh.

Cả nhóm bốn người chuẩn bị sẵn các âu kem còn dư lại để tiến hành công việc này, với các âu kem cho các tầng trên được đặt trên một tấm ván kéo bằng ròng rọc. Subaru - cô nàng quản lý - bắt đầu lèo lái mọi người trong khi hai tay giữ lấy thang: ``Phết đều nha! Và nhớ là phải đúng loại kem đấy!''

Choco và Luna ở phía trên cùng nhau đồng thanh đáp: ``Được rồi(-nanora)!''

``Chị bắt đầu xoay nhé!'' Suisei từ phía dưới gọi vọng lên.

Chiếc đĩa xoay bắt đầu chuyển động chầm chậm theo hướng chiều kim đồng hồ. Cả hai người đứng vững trên thang gập và bắt đầu công việc của mình.

Trong khoảng thời gian năm phút đầu, mọi việc đều diễn ra thuận buồm xuôi gió. Từng lớp kem mà Choco phết ở tầng năm và tầng sáu một cách mịn màng, đều tăm tắp và cực kỳ tỉ mỉ, khiến nó trông như một chiếc bánh làm bởi một nghệ nhân đồ ngọt lâu năm.

Nhưng cũng chính vì sự tỉ mẩn đó mà Luna cùng Suisei bắt đầu sốt ruột.

``Làm nhanh lên nào, Choco-sensei!'' Subaru nhặng xị. ``Ta sắp hết thời gian rồi!''

``Cứ bình tĩnh nào, Subaru-sama! Đừng có giục cô chứ!''

Không chần chừ thêm một giây phút nào nữa, Luna cũng bắt đầu làm công việc của mình. ``Em cũm bắt đầu phải nàm thui-nora! Cứ để cô nàm như thế thì nâu nắm na!!''

Thay vì dùng thìa phết bột (spatula) như mọi người, nàng Kẹo lại\dots\ cầm chiếc thìa ăn sữa chua.

Và cô cũng không phết kem một cách đàng hoàng, mà thay vào đó, nàng ``Em bé'' lại đang\dots\ bôi phết như đang vẽ tranh mẫu giáo.

``Luna-chan làm nghiêm túc đi chứ!!'' Subaru gằn giọng.

``Thì Luna đang nàm nghim túc mà!'' nàng Kẹo đáp, tay vẫn đang vân ve âu kem ở bên cạnh. Nhưng cách mà cô nàng nghịch kem trong âu cho thấy nàng Kẹo không hề để tâm tới công việc một chút nào.

Để rồi, trong lúc đang lấy kem trong âu, Luna vô tình vấy một ít kem việt quất lên tóc của Ryuuhou.

Cô em gái của Kuon ngẩng đầu lên, chống hông hỏi: ``Em làm cái trò gì vậy, Luna-chan!?''

Nàng Kẹo hét lớn: ``\textbf{Em nhỡ tay-nanora!}'' đồng thời lè lưỡi tỏ vẻ vô tội.

Cô bé nhà Hikawa thở dài với lời biện minh đó, rồi tiếp tục chú tâm vào công việc chính.

Ở phía dưới, Suisei đang cố gắng phết kem tầng ba. Mặc dù đã rất cố gắng căn đều từng góc, nhưng do tay hơi run và đĩa bánh xoay hơi nhanh hơn dự tính, đường phết của cô trở nên xoắn xít như mô hình xoắn ốc của hàm số Fibonacci.

``Ối trời\dots\ sao nó lại vặn như bánh rán thế này!?'' nàng Ca sĩ rên rỉ, khẽ trố mắt với ``tác phẩm'' của mình.

Subaru bên cạnh quay mặt sang, cũng không giấu nổi sự ngạc nhiên: ``Gì vậy!? Chị đang vẽ\dots\ mê cung trên bánh à!?''

Suisei chống nạnh: ``Chị đang cố vẽ xoắn đều đó!''

Nàng Vịt chặc lưỡi: ``Thôi chị để em sửa lại\dots'' rồi giành lấy thìa phết từ tay Suisei.

Nhưng chính vào lúc đó, chiếc đế xoay bất ngờ bị trượt bánh răng một chút, khiến nó giật nhẹ. Subaru suýt mất thăng bằng, đâm thìa xuống bánh một đường như máy xúc đào đất.

Một đường vệt lộ ra ở trên hông của chiếc bánh khiến cả hai toát mồ hôi hột.

``Không được nói với ai nha?'' Suisei thì thầm.

``Để em bôi lại. Nhanh thôi,'' Subaru đáp, mồ hôi nhễ nhại.

Thế là, tầng ba có một vết lõm sâu đáng ngờ, nhưng được cả hai người che kín lại một cách nhanh gọn, mặc dù có phần hơi đoảng - và cũng từ đó, họ cẩn thận hơn nhiều.

``Mình không phải là pâtissier\dots\ nhưng chắc ít nhất\dots\ cũng không nên làm bánh thành hồ bơi.'' Suisei kết luận.

``Ừ, mình là idol mà.'' Subaru gật đầu.

``Idol làm bánh méo.''

``Idol méo làm bánh.''

``\dots''

``Được rồi đừng nói nữa.'' - cả hai cùng im lặng và tiếp tục công việc. ``Hết tầng ba rồi, còn hai tầng nữa đó!''

``\dots''

Dù loay hoay, méo mó, hỗn loạn đủ đường, nhưng chiếc bánh kem vẫn dần dần được hoàn thiện, dưới những tiếng cười, tiếng càu nhàu, và sự phối hợp bất chấp logic của bốn cô gái không chuyên.

``Nghệ thuật là ánh trăng lừa dối.'' nàng Ca sĩ thở dài, hai tay xoay chiếc đế bánh.

``Bánh này là ánh trăng méo rồi.'' Subaru chốt.

``Kệ thui, mĩn là Kuonsa-senpai ăn được là OK-la thui na\~{}'' Luna từ trên lầu hét vọng xuống.

``Cả đầu em dính kem rồi kìa, Luna-sama.'' Choco thở dài.

\ct

Khi những vệt kem cuối cùng được phết khẩn trương của nhóm trước, nhóm cắt gọt trái cây và những phụ phẩm liên quan cũng đã bắt đầu tiến hành công việc.

Okayu và Aqua vui vẻ cắt những trái cây trong khi tán gẫu nhau - một điều mà nàng Hành tím ao ước từ lâu.

Nàng Miêu hí hoáy cắt quả kiwi, trong khi nàng Hầu gái thì tỉa quả nho, như thể đây là tiết thủ công đầu bếp.

``Okayu-chan! Bà chọn nho xanh kiểu gì mà nhỏ vậy? Tôi tỉa tới nửa giờ rồi mới ra được cái hình trái tim này đó\dots'' Aqua than thở.

``Tôi đâu có phải là người mua chúng đâu\~{} Subaru-chan với Choco-sensei chọn chúng mừ\~{}'' Okayu đáp. Cô liếc nhìn về quả nho mà nàng Hầu gái vừa cắt gọt, phì cười: ``Sao trái tim của bà trông như quả phổi vậy?''

``M-Mồ! Đây là quả trái tim!'' nàng Hầu gái đỏ mặt, định lên tiếng đáp trả thì bỗng nhiên\dots

*XOẸT!*

Cô nàng cắt trượt quả nho khiến cho nó bay thẳng lên trời, rồi nhẹ nhàng đáp xuống phần kem ở tầng hai.

``Ồ\~{}'' Okayu và Suisei vỗ tay, trong khi Ryuuhou thì nhăn mặt: ``Tập trung vào cắt quả đi chứ, Aqua-onee-chan. Cắt vào tay đấy.''

Aki thì điềm nhiên như không có gì xảy ra, cầm quả dâu lên, bắt đầu cắt tỉ mỉ thành hình hoa hồng.

``Chị đúng là cao thủ bếp núc\dots'' Rushia lẩm bẩm, nhìn quả dâu đang được gọt. ``Mà chị có cắt được thành hoa không vậy?''

``Chị nghĩ là được đó,'' nàng Dị giới đáp, cười nhẹ. ``Chị có học được mẹo này từ một video trên mạng à.''

Về phía Rushia, phần việc của cô là dùng vỏ cam để tạo thành những dải xoắn cam vàng uốn lượn quanh viền bánh.

Cả bốn người tiếp tục bày biện các loại trái cây đã cắt được: gồm những quả dâu hồng tươi mọng phủ mặt bánh ở tầng ba, những quả nho và việt quất xếp thành hình vương miện ở tầng năm, những lát kiwi hình ngôi sao và lát cam mọng nước xếp thành hình xoắn ở tầng một.

Và dĩ nhiên, vẫn còn nhiều loại quả nữa mà họ chưa động tới, như đào, lê, táo hay mận (Bắc). Thế nhưng, với mọi người, chiếc bánh kem này đã được trang trí hoàn thiện.

Suisei đứng lùi dần sát đến tường đối diện, thốt lên trong sự kinh ngạc: ``Ái chà\dots\ Bánh mà đẹp thế này thì Kuon-senpai chắc sẽ không nỡ ăn mất.''

``Mặc dù chiếc bánh có vẻ như chưa được hoàn chỉnh lắm\dots'' Rushia nghiêng đầu, nhìn ra những lỗi vặt trong chiếc bánh, nhưng rồi cũng lắc đầu gạt phăng đi: ``Nhưng mà thế cũng ổn rồi-nanodesu.''

Okayu vẫy tay: ``Đừng lo, chắc chắn anh ấy sẽ thích thôi.''

Aqua gật đầu đồng tình: ``Đúng đó đúng đó! Đây là công sức hơn một tiếng rưỡi chúng ta làm mà!''

Cô cầm lấy một đĩa đào định xếp tiếp vào chiếc bánh, nhưng rồi Subaru lên tiếng: ``Akutan, thôi được rồi. Tôi nghĩ ta cứ để cho nhóm trang trí phòng ở trên hoàn thành nốt đi.''

``Đúng đó, Aqua. Tôi nghĩ chúng ta thế là quá đủ rồi,'' Shion nói.

Aki bổ sung thêm, xách giỏ đựng những nguyên liệu khác: ``Ta vẫn còn si-rô, kẹo lấp lánh, hạt rắc với bánh quy nữa cơ mà. Nhóm ở trên tầng làm nốt cho.''

``Zà cũng zảm tải áp lựt nữ-nanora!'' Luna kết luận, giơ ngón tay cái lên, trong khi miệng vẫn đang ngậm quả nho.

``Quả nho đó là để xếp mà\dots'' Ryuuhou đổ mồ hôi hột. ``Em ăn vụng đấy à?''

``Đâu có đâu-nora\~{} Tự dưng nó chui vào mịng em đấy chứ\~{}'' cô nàng ngây thơ đáp lại, nuốt quả nho như thể chưa có chuyện gì xảy ra.

Cả bốn người cùng ngồi bệt xuống sàn, ngắm nhìn tác phẩm bánh ngọt cao sáu tầng đang rực rỡ hơn bao giờ hết - dù vẫn còn một vài chỗ chưa được hoàn thiện.

``Chỉ còn phần dặm màu và trang trí viền nữa thôi,'' Choco nói, lau tay vào tạp dề.

``Để lại cho đội cuối cùng. Mình làm như vậy cũng nhiều rồi,'' Mel vươn vai.

Noel lau chiếc chùy như một thói quen, tươi cười: ``Làm bánh cũng như làm phép, cần kiên nhẫn, tinh tế và\dots\ không nên cắt vào tay.''

Aqua cười khúc khích, suýt làm rơi luôn con dao gọt. ``Xem đứa phá hoại kia nói gì kìa!''

Tất cả mọi người đều cùng bật cười với lời bình của nàng Hầu gái. Sau những giờ làm việc khá cực nhọc và vất vả, họ xứng đánh tự thưởng cho mình một khoảng thời gian nghỉ ngơi.

Nhưng rồi, trong một thoáng chốc nói chuyện, cô em gái của Kuon bất chợt nhăn trán lại, rồi cất giọng hỏi: ``Ơ mà\dots\ mọi người\dots\ có ai biết mình sẽ đưa cái bánh này ra ngoài kiểu gì không?''

Một sự im lặng bao trùm cả căn phòng. Họ hướng tới chiếc cửa, vốn chỉ rộng vừa đủ cho một người cao lớn cỡ Noel.

``\dots Ừ nhỉ.'' / ``Giờ chị mới để ý\dots''

Lại một sự im lặng khác. Dường như họ không tính đến trường hợp này.

Luna nhìn chiếc bánh khổng lồ kia, rồi ước lượng: ``Chíc bánh này bự gấp đui, gấp ba lần cửa luôn đó-nanora.''

Subaru đưa tay đo thử chiều ngang tầng một của bánh, rồi quay sang đo khung cửa. Sau một hồi ước lượng, cô Vịt chỉ có thể lắc đầu: ``Không ổn rồi. Ta không tài nào mang nguyên cái bánh này ra ngoài được.''

Choco thở dài, khẽ lắc đầu thất vọng: ``Giá như chúng ta nên làm từng tầng xong mang lên ghép lại thì tốt hơn á.''

``Hoặc ít nhất là lắp ghép xong rồi mang trên kia phóng to cũng ổn mà,'' Mel nói góp, tỏ vẻ e ngại.

Suisei tiếp lời: ``Bây giờ thì\dots\ làm sao ta có thể cho chiếc bánh ra ngoài được ta?'', chau mày suy nghĩ.

Tới đây, mọi ánh nhìn đổ dồn về phía Shion. Nàng Phù thủy dường như đã lường trước được tình huống này, nên cô nàng không còn tỏ vẻ ngạc nhiên nữa.

``Rồi rồi, tôi hiểu rồi. Thu nhỏ nó lại là ổn chứ gì?'' cô thở dài một tiếng. ``Mà, mọi người nói cũng đúng đấy. Đáng ra ta phải xây bánh xong rồi mang lên hóa to mới phải.''

``Làm tốn ma lực của chị rồi, Shion-onee-san\dots''

``Không sao, không sao,'' nàng Tỏi phẩy tay. ``Nếu là vì Kuon-senpai thì từng này còn không ăn nhằm gì.''

``Ghê nha, ghê nha\~{}'' nàng Miêu buông lời trêu chọc. ``Thế tại sao bà không dùng phép thuật của bà làm việc có ích vào đời sống đi?''

``Im đi nha.''

Cô giơ bàn tay của mình hướng về chiếc bánh kem khổng lồ, lẩm bẩm đọc câu thần chú, rồi hô lớn: ``\textbf{Thu nhỏ lại!}''

Một tia sáng màu bạc xám xoáy quanh chiếc bánh trong tích tắc, và rồi, chiếc bánh khổng lồ giờ chỉ còn cỡ bằng nồi cơm điện.

``Ể\~{} Trông nó dễ thương quá đi!'' Aqua nhìn chiếc bánh kem ở dưới sàn, tỏ vẻ thích thú.

Subaru tỏ vẻ dè chừng: ``Nhỡ như ai tiện tay ăn hết trước khi phóng lớn thì sao?'' đồng thời mắt liếc về kẻ khả nghi lớn nhất đang mút kem trên ngón tay.

``Ý chị là xao chứ, Subarucha-senpai!!'' Luna hét toáng lên, lắc đầu nguầy nguậy như đang chối bỏ chủ ý của mình.

``Nhìn mặt em là đoán được em đang định làm thế rồi\dots'' / ``Không hề giấu giếm luôn kìa\dots'' cả Choco và Mel nhăn mặt.

Sau một hồi gói ghém, cả nhóm gom thêm các món trang trí còn dư gồm vài bịch kẹo lấp lánh, vài lọ rắc trang trí, một gói bánh quy hình idol, những chai lọ si-rô cùng bạc hà và một tấm bảng mini - mà thực chất là thanh sô-cô-la ăn được - để ghi chữ.

Tất cả được cẩn thận cho vào khay, buộc dây và giao lại cho nhóm phòng tiệc ở tầng ba. Riêng Shion và Subaru thì đưa chiếc bánh kem lên trên tầng ba đem lên trước.

Mỗi người một túi, vừa xách đồ lên vừa trò chuyện với nhau đầy rôm rả về buổi làm bánh bát nháo nhưng vui vẻ vừa rồi.

\dots Ngoại trừ một người.

Ryuuhou - cô bé ra khỏi phòng sau cùng - quay đầu nhìn lại căn phòng mà họ vừa rời khỏi, vốn đó là phòng của Aloe với căn bếp họ mượn tạm.

Và những gì ngổn ngang trước mắt càng làm cho cô thở dài chán nản hơn: một mảng trần bị nứt từ cú hoắ lớn chiếc lò của Shion; những hạt bụi bột mì, bột nở và các phụ phẩm khác bay vất vưởng trong phòng, phủ đầy trên mặt sàn và trên cả rèm cửa.

Ở góc phòng, mọi chuyện lại càng tồi tệ hơn: vỏ trứng và vết lòng trắng lòng đỏ còn ở trên tường, trên sàn gốm và chất đầy tại nơi đã từng là căn bếp. Những vết kem mà Luna nghịch vẫn còn dính đầy ở nắm cửa và ở trên tường.

Những gói bao bì đường, bột nở, các loại giấy bạc bị xé bỏ bị vứt lộn xộn ở một nơi khác. Và ngớ ngẩn nhất, là một chiếc nồi bị để ngược đặt trên chiếc loa - mà có lẽ nàng Hiệp sĩ đã nghịch trong khi chờ được làm.

Cô bé thở ra thật dài, gần như tuyệt vọng: ``\vie{Aloe-onee-chan mà thấy được cảnh tượng này chắc là ngất xỉu mất\dots}''

Chiếc bánh kem sáu tầng với hàng chục hương vị khác nhau giờ đã được thu nhỏ an toàn, và nhóm làm bánh lục tục rời tầng hai. Họ để lại phía sau không gian hậu trường mà chỉ có thần linh hoặc một công nghệ hiện đại nào đó mới có thể dọn dẹp hết mớ này.

Nhưng cũng tại căn bếp đầy hỗn độn ấy, một điều ngọt ngào đã thành hình.

Một chiếc bánh sinh nhật khổng lồ, do mười hai cô gái làm từ con tim, bằng bột, kem, bằng sự kiên trì và cả phép thuật.

Họ mong chờ rằng chủ nhân của bữa tiệc này - Hikawa Kuon - sẽ thích chiếc bánh kem đặc biệt, thậm chí là đậm chất \textit{\hll} như thế này.

\subsubsection*{\phantomsection\label{P4.1c}}
\addcontentsline{toc}{subsubsection}{Nhóm 4 - Nhóm chuẩn bị phòng tiệc}

\begin{center}
  {\color{T2} Nhóm 4 - A-chan và những người còn lại\\Nhóm chuẩn bị phòng tiệc}\\
  \uline{Quay lại một khoảng thời gian ngay sau khi mọi người phân công nhiệm vụ.}
\end{center}

Trong lúc mười một cô gái đảm nhận phần bếp núc ở các tầng khác nhau, hai mươi người còn lại - dưới sự chỉ đạo chung của A-chan và Tokino Sora - bắt đầu tiến hành công cuộc trang trí căn phòng tầng ba, nơi sẽ diễn ra bữa tiệc sinh nhật tối nay.

Đây là một căn phòng lớn, vốn từng là phòng sinh hoạt chung của nhà Hikawa trước khi được cải tạo thành không gian mở, nối liền hai cầu thang (trái và phải) - một hướng dẫn lên phòng Kuon ở tầng năm, một hướng dẫn đến phòng Hijaki ở tầng bốn. Hai nhà vệ sinh đặt đối xứng nhau ở hai góc trong cùng, và diện tích phòng đủ rộng để chứa khoảng chừng năm mươi người một lúc mà không quá chật chội.

Để dễ chia việc, cả nhóm được phân làm hai đội nhỏ:

Bên trái (bao gồm nửa trái trần, tường trước mặt, và nửa trái bức tường sau lưng) là nhóm của A-chan, gồm: Botan, Flare, Nene, Aqua, Ayame, Matsuri, Marine, Korone, Fubuki và Pekora.

Bên phải (nửa còn lại của trần và hai tường) do Sora chỉ huy, với các thành viên: Roboco, Miko, Haato, Mio, Polka, Towa, Kanata và Watame, Polka.

Ngay từ khi bắt đầu, A-chan và Sora đều rất nỗ lực điều phối đội hình. Thế nhưng, thực tế chứng minh một điều: không phải ai trong số các thần tượng của \hll\ cũng phù hợp với công việc trang trí!

\begin{center}
  \uline{Bên trái bức tường.}
\end{center}

Mười một cô gái trong nhóm bên trái - dẫn đầu bởi A-chan - bắt đầu tụ lại gần bức tường trước mặt, nơi được giao nhiệm vụ chung ``trang trí ấn tượng nhất phòng''.

Aqua - dù sau này phải quay xuống tầng hai để làm bánh cùng nhóm với Ryuuhou - hiện tại vẫn đang có mặt cùng nhóm, đúng vào thời điểm trước khi Kuon bước xuống nhà.

Polka nhảy lò cò quanh cuộn dây trang trí rồi nhanh nhảu hỏi: ``Chúng ta sẽ làm gì với bức tường này vậy?''

``Chị nghĩ chúng ta phải làm gì đó để trang trí bức tường này,'' Mio khoanh tay đáp lại, rồi gục đầu xuống suy nghĩ: ``Nhưng vấn đề là ta nên làm gì nhỉ\dots''

A-chan chỉnh lại gọng kính, cất giọng nghiêm túc nhưng không quá căng thẳng: ``Ai có ý tưởng gì cho mảng tường này không? Gợi ý chút cũng được!''

Matsuri ngồi bệt xuống sàn, tay chống má, mặt lơ đãng nói: ``Tính cách của anh ấy như kiểu `gì cũng được' mà\dots\ nên em chẳng biết bắt đầu từ đâu luôn.''

Aqua ngó sang, nhíu mày đầy lo lắng: ``Xem chừng có vẻ khó nhỉ\dots''

Fubuki đưa tay xoa cằm, ánh mắt hơi nheo lại suy tư: ``Hmmm\dots\ Treo băng-rôn có được không nhỉ?''

Nàng Cổ động giật mình phản đối, mắt sáng lên đầy nhiệt huyết: ``Thế thì bình thường quá rồi! Nghĩ cách nào sáng tạo hơn chút đi!''

``Vậy thì ta treo băng-rôn rồi rắc bột lấp lánh bên trên!'' nàng Cáo trắng hớn hở đáp.

Flare nhướng một bên chân mày, búng nhẹ tay vào trán tiền bối: ``Có khác gì đâu trời!? Cùng là treo đồ lên tường thôi!''

Nene giơ hai tay lên như thể đang cầm đèn lồng vô hình: ``Sao mình không dùng đèn trang trí nhỉ? Treo kiểu vòng cung chẳng hạn!''

``Ý tưởng tuyệt đấy!'' Matsuri reo lên đầy phấn khích, nhưng rồi nhìn lại đống đồ họ mua, thắc mắc: ``Nhưng mà chúng ta có mua đèn đâu\dots?''

Ayame chớp mắt một cái, rồi cười khẩy như một con quỷ: ``Ta cứ lấy dây đèn ở quán cà phê bên cạnh là được ha-dazo\~{}''

``Liệu có được không vậy?'' Marine tò mò với ý tưởng có phần khác thường này. ``Chưa kể việc căn phòng này rất rộng nữa.''

``Thừa lúc không ai để ý, ta sẽ\dots\ nhẹ nhàng lấy một cuộn dây thôi! Không ai biết đâu!''

Cả A-chan và Flare ngã khuỵu ra phía sau. Họ cùng đồng loạt lên tiếng phản đối: ``Đó gọi là ăn cắp đấy Ayame-senpai/Nakiri-san!''

Pekora vừa ngồi xếp bánh quy và đồ vặt, vừa gật gù đồng tình: ``Đúng đó, peko. Người ta báo cảnh sát đấy.''

``Nhất là sau những gì chúng ta đã làm ở trên đường phố Kawachi đó\dots'' Korone gật gù.

``Nếu là cảnh sát thì chắc em cũng sẽ hạ gục họ thôi,'' Botan nhíu mày, đảo mắt không vui.

Flare thở dài, tay chống hông, đảo mắt: ``Mọi người đừng lạc đề nữa. Quan trọng bây giờ là\dots\ chủ đề. \emph{Chủ đề trang trí.} Cái mà Kuon-senpai thích.''

Không chút do dự, nàng Khuyển vỗ tay: ``A! Cái này chị biết! Làm việc tăng ca á!''

Cả đám lập tức quay sang nhìn A-chan. Cô nàng đứng hình đúng ba giây rồi thét gào trong lòng: ``\textbf{\textit{Tại sao lại lôi chị vào vụ này chứ!?}}''

Aqua toát mồ hôi hột: ``Tại sao bà ấy lại nghĩ đến chuyện đó nhỉ?'' Pekora cũng đồng tình: ``Đúng là anh ấy làm việc nhiều thật, nhưng đâu có nghĩa là anh ấy cuồng việc đâu, peko.''

Nàng Thuyền trường khoát tay: ``Vậy là tàu cướp biển! Chắc chắn ảnh sẽ thích phong cách phiêu lưu mạnh mẽ đó!''

Nàng Quỷ sừng chen vào: ``Hoặc là\dots\ một thanh kiếm Quỷ to bự giữa tường luôn, yo!''

Botan thì tỏ vẻ bình tĩnh hơn tất cả, cô khoanh tay ngồi dựa tường và đề xuất: ``Sao không treo một khẩu súng mô hình?''

Nene nhảy tưng tưng: ``Hay mình làm một chiếc\dots\ xe đẩy hàng phiên bản bento nhỉ?''

Fubuki thì đưa ra đề xuất một cách tỉnh bơ: ``Hoặc là treo luôn cả hang động tuyết thu nhỏ?''

Flare lúc này chỉ còn biết nhìn lên trần nhà, tay vò trán, miệng thở ra một tiếng dài: ``Trời ơi\dots\ mấy người đang nói cái gì vậy không biết\dots''

A-chan vội xua tay, điều chỉnh lại luồng suy nghĩ của nhóm: ``Ý của Shiranui-san là\dots\ một chủ đề mà ai cũng thích cơ. Điều gì mà mọi người nghĩ đến đầu tiên?''

Mọi người thoáng im lặng ba giây, người thì khoanh tay, người vò đầu, rồi tất cả đồng thanh thốt lên với những giọng pha trộn: ``Trò chơi điện tử ạ?''

Flare vỗ hai tay vào nhau, gật gù: ``Ví dụ như thế! Nhưng mà\dots\ chủ đề nào thì được?''

Korone, Aqua, Botan và Nene cùng đồng thanh, mắt sáng rỡ: ``Minecraft chăng?''

Nàng Thỏ nghe vậy liền nhíu mày, lắc đầu nhẹ: ``Tôi không nghĩ là anh ấy khoái cái trò lắm đâu. Mà nhắc mới nhớ, Kuon-senpai thích chơi trò gì được nhỉ, peko?''

Senchou chống cằm suy nghĩ, rồi bất chợt bật ra:``Em nhớ là hồi trước ảnh có chơi cái trò bắn chim gì cơ mà\dots\ con chim đỏ bắn lợn xanh á.''

``Hình như là A*g*y B***s, đúng không?'' nàng Cáo Trắng góp lời, giả vờ bắn đạn bằng chiếc súng cao su tưởng tượng.

Nàng Hầu gái nghe tới đó suýt ngã: ``Cái trò cổ lỗ sĩ đó á!? Năm 2021 rồi ai còn chơi cái đó nữa?''

``Thế mà có người vẫn chơi á\~{}'' nàng Chó vẫy đuôi. ``Đổ cổ vẫn là đồ tốt mà!''

Nàng ``Hải cẩu'' hồn nhiên tiếp lời: ``Nene để ý thấy anh ấy còn chơi cái trò bóng gì cơ mà! Cứ vuốt vuốt rồi né mấy cục gì đó ấy.''

Cô bạn cùng thế hệ nghe thế, nhướng mày khó hiểu: ``Trò gì nghe lạ vậy? Nhưng mà\dots\ nếu là chúng ta thì hay chơi gì?''

Pekora nhún vai rồi nói đầy tự tin: ``Nếu là em thì thích chơi mấy trò như MSG lắm nha, peko!''

Botan gật đầu: ``Em lại chơi mấy game FPS nặng đô.''

Marine nâng ngực, hất tóc: ``Trò gì chị cũng chơi được hết! Gái đa tài là thế!''

Ayame cười ngượng: ``Em thì lại giống Aqua-chan.''

Aqua lập tức nắm tay Ayame: ``Đúng đúng! Có đồng minh rồi!''

Matsuri vắt tay sau gáy: ``Khó rồi đây nhá\~{} Không biết Senpai thích gì thật.''

Flare thở ra một hơi. Thực sự cách tiếp cận này cũng không mấy hiệu quả. Cô bèn tính nước đi khác. ``Ý mình là tìm một thứ mà Kuon-senpai thích, cơ mà\dots\ thôi, hay là thế này: ta cứ trang trí theo cái mình thích đi?''

A-chan chớp mắt rồi mỉm cười đưa ra quan điểm: ``Chị thấy cách đó lại hợp lý hơn đấy. Mỗi người một phong cách, hợp với sự đa dạng của nhóm. Nhưng mà các em định trang trí như thế nào? Ta có một bức tường chính này và một nửa hai bên cạnh tường nữa.''

Fubuki bước tới, giơ hai tay ra căn chỉnh không gian như một đạo diễn sân khấu: ``Phải thêm cả hình ảnh Kuon-senpai vào nữa chứ, vì đây là bữa tiệc của ảnh mà!''

Cô đặt hai ngón cái và hai ngón trỏ căn khung như camera: ``Nhưng chỗ này\dots\ có vẻ hơi ít không?''

Aqua nghiêng đầu: ``Trông có vẻ đủ mà? Tôi thấy ta có thể vẽ một bức tranh thật lớn!''

Ayame chỉ tay, hớn hở nói: ``Vậy thì mười người chúng ta chia nhau ra vẽ nha! Em xí phần giữa này!''

Nói xong, nàng Quỷ sừng bật người chạy tới ôm chặt lấy khu vực giữa bức tường bên trái như thể những người hâm mộ đang tranh chỗ ngồi đẹp nhất trong buổi biểu diễn của các cô gái nhà Tam Giác Xanh.

Chỉ trong chớp mắt, tám người còn lại lập tức lao theo và tranh giành chỗ của mình.

``\textbf{Á! Chỗ đó là của chị mà!!}'' Matsuri hét lên, giơ tay túm áo nàng Quỷ sừng.

Fubuki trượt tới từ bên cạnh: ``Là của tôi chứ!''

``\textbf{Của em mà, peko!}'', Pekora chen ngang, đẩy Nene ra khỏi bức tường lớn.

``\textbf{Của bọn này chứ!!}'' Nene và Botan đồng thanh, rồi quay sang nhìn nhau đầy nghi ngờ: ``Khoan, sao bọn mình lại tranh nhau nhỉ?''

Cả nhóm lao vào thành một khối tranh giành, chen lấn, gào thét.

Matsuri lôi cả tay áo Fubuki, Pekora dùng chiêu nhảy hai chân cố đẩy Ayame ra, Marine thì giơ tay hô to ``\textbf{Chiếm quyền thống trị khu trung tâm!!}''

Ở góc phòng, Flare - người đã khởi xướng cả ý tưởng này - chỉ biết thở dài, tay chống hông, mắt lườm nhẹ: ``\dots Y như một đám con nít lần đầu được phát bút màu vậy.''

A-chan lúc này cũng không thể đứng ngoài nữa. Sau một cái chỉnh kính rất khẽ, cô bước tới giữa nhóm đang tranh cãi, vỗ tay hai cái thật mạnh: ``Rồi, rồi! Mấy đứa, bình tĩnh nào!''

Tiếng ồn ào giảm xuống một chút. Những cái đầu còn đang giành nhau phần giữa bức tường từ từ quay lại, có vài người còn lấm lét như học sinh bị bắt gặp đang nói chuyện riêng giữa giờ học.

Nàng Đeo kính dõng dạc, tiếp tục đưa ra ý tưởng: ``Mấy đứa cứ tranh nhau một chỗ thế này thì chẳng ai làm được cái gì cả đâu! Chị thấy là tốt nhất nên chia khu vực rõ ràng, giống như bên kia ấy!''

Flare nhướng mày khó hiểu: ``Bên kia?''

Cả hai người cùng quay về phía bên phải căn phòng - nơi đội của Sora đang làm việc.

Cảnh tượng bên ấy khác hẳn với hỗn độn tại đây: tuy không hoàn toàn yên bình, nhưng có thể thấy họ đã chia tường thành từng khu rõ ràng, mỗi người đang tập trung làm một phần riêng biệt: Miko dán chữ giấy bằng băng keo màu, Roboco đang xếp giấy bóng lên viền tường, Mio thì trải vải bàn, còn Towa đang vẽ hình các ác ma phiên bản hiện đại hóa.

``Ủa? Bên đó chia khu rõ ràng từ sớm luôn á hả?'' nàng Dị giới lẩm bẩm.

A-chan gật đầu: ``Chắc là do Sora-chan sắp xếp đấy. Mà chị thấy cũng hợp lý\dots''

``Vậy sao mình không `học lỏm' họ luôn nhỉ?''

``Ý hay đấy.''

Không chần chừ thêm, nàng Đeo kính quay sang cả nhóm, chỉ tay về phía tường trái: ``Giờ thế này. Chúng ta có bức tường chính này, và mỗi bên tường hông là một nửa. Chị sẽ chia làm ba phần: bên trái, giữa, và phải.''

Matsuri giơ tay phản đối: ``Khoan đã! Em vẫn muốn phần giữa đó nha!''

A-chan nghiêng đầu, giọng vẫn dịu nhưng cứng rắn: ``Muốn phần giữa thì phải oẳn tù tì, Natsuiro-san. Công bằng với mọi người.''

``G-Giờ mới oẳn tù tì á, peko!?'' Pekora lắp bắp.

``Biết thế làm vậy ngay từ đầu,'' Aqua thở dài, đoán định được mọi chuyện sẽ diễn ra như vậy.

Marine cười khẩy: ``Thế thì chị không nhường đâu nhé.''

Botan bình thản kéo tay áo lên như chuẩn bị vào trận, không nói một lời nào.

Nene thì đang vẽ sẵn sơ đồ chiến thuật như sắp ``xí đất''.

Ayame và Korone nhìn nhau đầy quyết tâm, rồi cùng gật đầu.

A-chan nhìn khung cảnh như sắp có một giải đấu cấp tốc, thì chỉ biết mỉm cười nhẹ: ``Rồi, ai muốn phần giữa thì oẳn tù tì. Còn không thì chị chia phần cho.''

Tất cả lập tức\dots\ bước vào tư thế.

``Oẳn tù tì, ra cái gì\dots\ ra cái này!''

\ct

Sau một trận oẳn tù tì căng thẳng và đầy mưu mẹo - kèm theo một vài pha ``ra lại'' đáng ngờ từ phía Marine và Nene - cuối cùng, kết quả cũng đã rõ ràng: phần giữa - vị trí trung tâm và được xem là tiêu điểm của cả bức tường thuộc về Marine, Ayame, Korone, Pekora và Nene; phần tường bên trái do Matsuri, Aqua và Botan hoàn thiện, còn phần phải sẽ do Fubuki, Flare và A-chan làm.

A-chan nhanh chóng vạch ranh giới tạm bằng băng keo màu để các nhóm không đụng chạm phần nhau, rồi lùi lại để cả đội bắt tay vào việc.

``Chị không cần các em phải làm giống nhau,'' A-chan dặn dò, ``chỉ cần mỗi người đều để một chút dấu ấn cá nhân vào đó là được. Nhưng nhớ là vẫn phải xoay quanh Hikawa-san nha.''

Cả chín người đồng thanh: ``\textbf{Rõ!}''

Ngay sau đó, không khí lao động trở nên sôi nổi hơn bao giờ hết. Mỗi nhóm triển khai phần mình với phong cách hoàn toàn khác nhau - không ai giống ai, nhưng vẫn hướng về một chủ đề chung.

Mặc dù phong cách hoàn toàn lệch pha, thế nhưng bằng cách kỳ diệu nào đó, chúng lại hòa quyện với nhau một cách hài hòa đến không ngờ.

A-chan lùi về một bước, ngắm toàn bộ mảng tường vừa hoàn thành với ánh mắt hài lòng. Cô gật đầu với Flare: ``Không cần phải giống nhau, nhưng vẫn kết nối với nhau - chị nghĩ Kuon sẽ thấy ấm lòng lắm đấy.''

Flare khẽ cười: ``Dạ. Có khi lại bất ngờ đến mức\dots\ cạn lời cũng nên.''

\ct

Ở phía bên kia chiến tuyến, không khí tuy có phần bình lặng hơn nhóm của A-chan, nhưng cũng không kém phần bối rối. Mười thành viên đứng thành hình vòng cung trước bức tường trắng, ai nấy đều ôm trán hoặc chống cằm, đôi mắt dõi lên khoảng không như mong ý tưởng sẽ tự rơi xuống.

Sora - người được giao nhiệm vụ chỉ huy - phá vỡ sự im lặng bằng giọng nhẹ nhàng, nhưng đủ để gọi sự chú ý: ``Vậy\dots\ mọi người định sẽ làm gì với bức tường này đây?''

Haato nhún vai, nghiêng đầu tỏ vẻ vô tư: ``Thì cứ trang trí theo kiểu mình thích thôi mà? Giống như đợt sinh nhật của YAGOO ấy!''

Mio chớp mắt ngạc nhiên, gật đầu nửa tin nửa ngờ: ``Ờm\dots\ Nhưng mà Kuon-senpai không như YAGOO đâu. Anh ấy nghiêm túc lắm á! Với lại đây là lần đầu ta tổ chức sinh nhật cho senpai, không biết anh ấy sẽ phản ứng như thế nào nữa\dots''

Miko đưa tay ra hiệu như đang đóng khung một bức tranh tưởng tượng: ``Chả phải Kuon-senpai dễ tính lắm sao ne? Kiểu `gì cũng được' ấy. Em nghĩ chắc ảnh không soi mói đâu.''

``Chính vì `gì cũng được' nên mới khó đấy,'' Kanata thở dài, xoa gáy, ánh mắt nhìn lên trần như đang cố tìm câu trả lời từ các vị thần tiền bối của mình.

``Xem chừng việc này cũng không hề đơn giản nhỉ,'' nàng Thần tượng cười gượng. Cũng giống như nhóm ở tường bên trái, dường như họ vẫn không tìm được cách giải quyết vấn đề ngay lập tức.

Cô đưa ra một giải pháp ban đầu: ``Vậy mọi người có ý tưởng gì chưa?''

Towa bỗng chống hông, ra vẻ nghiêm túc hiếm thấy: ``Thế\dots\ ta làm theo phong cách cyberpunk thử xem?''

Polka lập tức nảy số: ``Còn không thì chủ đề xiếc? Trang trí như rạp diễn luôn!'

Watame cười nhẹ, vỗ vai đàn em của mình: ``Chị tưởng em sẽ nói là làm hình Kuon-senpai đội mũ chóp với cái mũi đỏ đó chứ.''

Nàng Cáo sa mạc khoanh tay nghiền ngẫm ý tưởng này, rồi nghiêm túc đáp: ``Thật ra cũng là một lựa chọn không tệ đó!''

Roboco ngồi xổm xuống, vẽ vẽ gì đó lên sàn bằng tay: ``Còn chị thì thích làm không gian kiểu kỹ thuật số - như là trang cá nhân của anh ấy vậy á\~{}''

Haato lẻo mép chen vào: ``Em thì thích trang trí những ruy băng đỏ và hình trái tim-chama!''

Mio lắc đầu với đề bạt có phần hơi quá sức này: ``Nghe thì hấp dẫn đấy, nhưng\dots\ liệu nó có phức tạp quá không?''

``Chưa kể rằng nó còn khó phối nữa chứ\dots'' nàng Thiên sứ cũng không đồng tình với đề xuất này.

Mio khoanh tay, mặt nhăn lại đầy khổ sở và thở dài một tiếng: ``Mỗi người một kiểu thế này thì đúng là khó tìm được tiếng nói chung thật\dots\ Có ai có ý tưởng gì không?''

Watame, giọng nhẹ nhàng, nghiêng đầu như đang mơ mộng: ``Đồng cỏ nghe thế nào?''

``Bầu trời thì sao nhỉ?'' Kanata nêu ý kiến, bàn tay đặt lên ngực như thiên thần đang suy tư. ``Giống như thiên thần bọn tôi ấy?''

Towa hăng hái góp lời: ``Tôi lại thích ở trên vũ trụ hơn. Giống cái hồi mà tôi với Roboco-san, Luna-chan với Aqua-senpai cùng Kuon-senpai bay lên\footnote{Xem lại {\flaregothic ÉPISODE 103}, quyển số Tám.} á!''

Nàng Người máy bật cười khúc khích: ``Ý em là cái lần chị vô tình đáp xuống văn phòng \hll\ chi nhánh Ai Cập sao?''

``Đúng, đúng! Rồi tụi mình đi theo cái cỗ máy thời gian lạ lùng đó!'' nàng Ác ma hào hứng vung tay tả lại.

``Ể!? Thật ư!? Chị cũng muốn thử cảm giác đó, ne!!'' Miko phụng phịu, tỏ vẻ ghen tỵ với hai cô gái đó.

``Ê, mấy người đừng có buôn dưa lê nữa được không!?'' nàng Sói đen nhăn mặt, cắt lời.

Haato chen vào, bỏ ngoài tai Mio: ``Trang trí thêm một chút ruy-băng đỏ thì sao nè? Cho nó tình cảm!''

``Và thêm chút phong cách xiếc nữa!''  Polka reo lên, ánh mắt lấp lánh hình trái tim.

Mio chắp hai tay lên trán, gào lên trong tuyệt vọng: ``Thế này còn ra thể thống gì nữa!?''

Sora khẽ cười trừ, cố giữ không khí ôn hòa: ``Được rồi, chúng ta bắt đầu từ thứ cơ bản nhất đi. Kuon-senpai thích gì nào?''

Watame ngẫm nghĩ một lúc rồi lắc đầu: ``\dots Ngoài game ra thì em chẳng nghĩ được gì. Còn lại em thấy\dots\ 50-50 lắm.''

Towa gật gù: ``Game cũng có nhiều loại lắm. Ta nên chọn cái gì đây?''

``FPS chẳng hạn?'' Roboco đề xuất.

``Một cái gì đó dễ thương thì sao?'' Haato nhíu mày hỏi.

``Cho thêm một chút\dots\ hỗn loạn thì sao nào?'' Polka cười tinh quái.

``Không, không\dots\ Phải là một thứ gì đó nhẹ nhàng, hiền hòa một chút\dots'' nàng Cừu phản đối.

Mio toát mồ hôi hột, cố gắng đưa mọi thứ về quỹ đạo: ``Ý của Sora-senpai là cái trò mà anh ấy thích chơi cơ mà!''

Towa khoanh tay, vẫn bỏ ngoài tai nàng Sói đen đầy phũ phàng: ``Thôi thì, chủ đề vũ trụ vẫn hợp hơn!''

Roboco hào hứng nắm tay hậu bối của mình: ``Để chị giúp em nha\~{}''

``Không, đồng cỏ xanh vẫn hơn chứ!'' Watame giật tay áo Roboco, nằng nặc thuyết phục được cô.

``Xiếc hợp hơn!'' nàng Cáo sa mạc vẫn giữ vững lập trường.

``Phải thêm ruy-băng đỏ nữa chứ!'' nàng Song tính lên tiếng.

``\textbf{Bầu trời!}'' nàng Thiên sứ hét lên lần nữa.

Và cứ thế, năm người họ tiếp tục tranh cãi như thể đang thi đấu khẩu chiến trên sân khấu truyền hình. Mio - cố gắng đóng vai người điều phối - bị cuốn vào như tâm bão, tay liên tục xua xua, thậm chí hét đến lạc cả giọng nhưng không ai chịu nghe.

Cho đến khi nàng Thần tượng đứng giữa, khuôn mặt tối lại và im bặt không nói câu nào, họ mới bắt đầu nhận ra điều bất thường.

Một bầu không khí u ám bất giác bao trùm khắp cả căn phòng, nhiệt độ trong phòng như thể giảm đột ngột xuống mặc dù trời vẫn còn nắng oi.

Cả nhóm đồng loạt im lặng, chuyển ánh nhìn về phía cô, rồi đồng loạt run lên bần bật.

``S-Sora-chan?'' Miko ngập ngừng lên tiếng đầu tiên.

``C-Chị làm sao vậy?'' Mio hỏi, giọng nhỏ đi hẳn, khuôn mặt tái mét.

Sora vẫn không nói gì. Mặt cô vô cảm, ánh mắt nghiêm nghị.

``B-Bọn em làm gì sai à?'' Haato sợ hãi.

``M-Mặt chị\dots\ đáng sợ quá\dots'' Polka lùi nhẹ một bước.

``Bình tĩnh nào chị ơi\dots'' Kanata và Towa đồng thanh.

``Chúng em không cố ý tranh cãi nảy lửa tới mức này mà\dots'' Watame giọng run run.

Giữa bầu không khí căng như dây đàn, khuôn mặt Sora bất ngờ\dots\ giãn ra. Cô chớp mắt khó hiểu, rồi hỏi: ``Mấy đứa bị sao vậy?''

``C-K-Khuôn mặt chị\dots'' nàng Cáo xiếc nói khẽ.

``\dots Đáng sợ quá luôn.'' Towa rụt cổ lại.

``Gì cơ?'' nàng Thần tượng nghiêng đầu. ``Ý mấy đứa là\dots\ thế này hả?'' cô gồng mặt làm lại biểu cảm ban nãy.

``Aaaaa đừng mà!'' cả nhóm gần như hét lên ngay khi nhìn thấy khuôn mặt kinh dị đó.

``Cho chị xin lỗi nha\~{}'' Sora bật cười, ngượng nghịu. ``Chị\dots\ không cố ý vậy đâu\dots''

Roboco, đứng phía sau, khoanh tay: ``Bà ấy lúc nào cũng vậy đấy\dots''

``Như vậy là mình vẫn chưa chốt được chủ đề nào nhỉ\dots\'' Sora thở ra, nhìn cả nhóm một lượt.

Mọi người giờ đây đã lấy lại tinh thần, đồng tâm hiệp lực suy nghĩ, lần này thực sự nghiêm túc hơn.

Người chỉ huy của nhóm đứng yên lặng một lúc, mắt dõi quanh từng người. Cô có vẻ đang cân nhắc, rồi nhẹ nhàng giải đáp: ``Nếu mọi người mỗi người một phong cách, chị sợ là nó sẽ lộn xộn và không có điểm nhấn đâu.''

Cả nhóm ngừng bàn tán và đổ mắt về phía cô sau lời nhận xét đó.

``Hay là\dots\ mình chọn một chủ đề riêng cho mình?'' cô đề xuất một ý tưởng khác, ánh mắt dịu dàng nhưng kiên định.

``Chủ đề?'' nàng Cừu nghiêng đầu.

``Ừ.'' Sora gật đầu. ``Chủ đề nào đó liên quan đến Kuon-senpai, chẳng hạn.''

Kanata chống cằm, mắt sáng lên:

``Kiểu\dots\ Kuon-senpai phiên bản siêu anh hùng?''

Haato lẩm bẩm: ``Hoặc là Kuon-senpai với những sợi tơ hồng\dots''

Polka phấn khích: ``Kuon-senpai bước ra từ rạp xiếc!''

Towa nheo mắt: ``Hay Kuon-senpai bước ra từ cyberpunk!''

Miko giơ tay: ``Còn Kuon-senpai trong thế giới đền thờ thì sao?''

``Bất cứ điều gì cũng được! Miễn là liên quan tới Kuon-senpai là được thôi!'' Sora mỉm cười dịu dàng.

Ý tưởng bắt đầu đổ về, đa dạng như chính tính cách của các thành viên đang đứng ở đây.

Mio tiến lại gần, tay khoanh nhẹ, nghiêng đầu nói nhỏ: ``\hl{Thật ra em định hỏi nãy giờ rồi\dots{} Cái ý tưởng chia chủ đề con này chị nghĩ ra từ khi nào vậy?}''

Sora nháy mắt tinh nghịch: ``Thú thật thì\dots\ chị nghe lỏm từ bên kia đó.''

``\dots Dạ?''

``Lúc nãy ấy,'' nàng Thần tượng thì thầm, che miệng lại và ghé sát vào đôi tai sói của Mio, ``chị nghe A-chan với nhóm bên trái bàn nhau chia tường và chia mảng vẽ theo nhóm cá nhân. Thấy cách làm đó hợp lý quá nên\dots\ chị `mượn' luôn.''

Nàng Sói đen che miệng bật cười: ``Vậy là\dots\ Sora-senpai học lỏm ạ?''

Sora gật gù, mặt không chút ngại ngần: ``Ừ, học lỏm đấy. Nhưng học có chọn lọc nha\~{}''

Cả hai cùng cười khúc khích, rồi liếc mắt nhìn sang phía bên kia tường, nơi nhóm A-chan cũng vừa lúc hoàn tất việc phân khu và đang hí hoáy bắt đầu nét vẽ đầu tiên.

Cuối cùng, Mio nheo mắt đùa: ``Đến cuối cùng thì cả hai nhóm đều học hỏi của nhau nhỉ.Như thể thần giao cách cảm ấy!''

``Ừ.'' Sora mỉm cười, ánh mắt dịu dàng dõi lên bức tường trước mặt, nơi giấc mơ của cả nhóm đang dần hình thành.

Việc phối hợp hai mươi con người - mỗi người một cá tính, một phong cách, một thẩm mỹ khác nhau - dần dần rơi vào cảnh ``chín người mười ý'' hay ``lắm thầy nhiều ma'', đúng như câu thành ngữ.

Nhưng, chính trong cái mớ hỗn độn đó lại ánh lên một thứ gì đó rất\dots\ riêng \hll.

Khi Nene và Polka tranh nhau màu nền pastel, Ayame với Pekora giành nhau độ dày của viền chữ, hay Kanata và Towa còn đang bàn xem màu thiên thanh hay tím trầm hợp với ``tầng khí quyển'', thì Sora lại chỉ nhẹ nhàng xoa dịu:

``Mỗi người có gu riêng, nhưng điều tuyệt nhất là được cùng nhau tạo ra một bức tranh không ai khác có thể làm được ngoài các em.''

Câu nói ấy, như một sợi chỉ nối những đường nét lộn xộn lại với nhau, tạo thành bố cục của một bức tranh chưa hoàn thành nhưng đầy tiềm năng.

Mio cũng góp lời, giọng đậm chất ``chị cả'' của hội bạn đồng niên: ``Vẽ sao cũng được, tô kiểu gì cũng được. Chỉ cần đừng làm mất đi cái tinh thần nhóm chúng ta là được.''

A-chan - vừa từ bên kia phòng bước sang, mỉm cười nhìn cảnh tượng trước mặt, nhẹ giọng tiếp lời: ``Cứ coi như đây là một MV sống động đi. Cảnh quay mỗi người khác nhau, nhưng đạo diễn là cùng một trái tim.''

Flare thì khoanh tay, gật đầu chậm rãi, lời ít nhưng sắc bén: ``Chỉ cần nhớ - đây là sinh nhật của Kuon-senpai. Tức là không cần phải vẽ về anh ấy, mà vẽ cho anh ấy thấy tinh thần của chúng ta.''

Lời nói ấy, như một lời nhắc nhở đơn giản nhưng đầy trọng lượng, khiến không khí cả phòng lặng đi một thoáng - rồi bật lên như sóng vỡ. Người này gật đầu với người kia, kẻ bắt đầu phác thảo, người bày dụng cụ, từng tốp nhỏ lại tự chia nhau việc mà không ai bảo ai.

Thay vì cãi cọ, giờ đây họ hỏi nhau những câu nhỏ nhặt:

``Tớ mượn màu này chút nha?''

`'`Mình phối phần trái tim bên trái được không?''

``Vẽ hình idol version của Kuon-senpai ở giữa nhé?''

Giữa hỗn loạn, là điều chỉnh. Giữa khác biệt, là hòa âm.
Không ai giống ai, nhưng tất cả đều cùng vì một người.

Và thế là, những mảng màu đầu tiên bắt đầu được phác lên bức tường, như thể cuộc hỗn chiến thẩm mỹ này chưa từng là hỗn loạn.

\ct

Trong lúc nhóm 1 đang hoàn tất những chi tiết cuối cùng trong gian bếp - những mẻ thịt được ướp đầy đặn, những nồi lẩu cũng đã sôi - thì từ phía hành lang vang lên tiếng bước chân đều đặn. Ryujin và Kaien, vừa hoàn thành công việc dưới bếp, thong thả bước lên tầng ba để nghỉ tạm. Vừa rẽ qua hành lang dẫn vào phòng khách, họ sững người lại trước khung cảnh trước mắt.

Trên mảng tường dài chia căn phòng làm hai, các thành viên \hll\ đang chăm chú vẽ, dán giấy màu, xếp các mảng giấy origami, cắt dán đề-can và sắp đặt những nét đầu tiên trong bản hoạ lớn. Bầu không khí vừa im lặng tập trung, lại vừa tràn đầy năng lượng sáng tạo.

\begin{dialogue}
  \speak{Ryujin} Ồ\dots\ Cái này là\dots''
  \speak{Kaien}\dots Một bức tranh ghép khổng lồ? Hay là một bức họa trang trí tường?
\end{dialogue}

Sora, đang cùng Watame điều chỉnh lại một nét vẽ màu xanh trời, liền quay lại và mỉm cười dịu dàng đáp lễ: ``Ryujin-san, Kaien-san. Bọn cháu đang vẽ bức tường này cho sinh nhật của Kuon-senpai đấy ạ. Chúng cháu chia tường ra thành ba phần - trái, phải, và giữa - để mỗi nhóm có thể tự do thể hiện ý tưởng của mình.''

A-chan, đứng gần đó với bảng màu trên tay, tiếp lời: ``Bên cháu là nhóm trang trí tường giữa. Chủ đề thì chúng cháu vẫn đang dần định hình, nhưng tụi cháu thống nhất là mỗi người sẽ thể hiện một góc nhỏ bằng phong cách riêng, rồi cùng nhau tái hiện hình ảnh Hikawa-san qua nhiều khía cạnh ạ.''

Người nội trợ nhìn lên tường, nét mặt trầm ngâm rồi gật gù: ``Cũng như một bức tranh nghệ thuật, ghép từ từng mảnh tâm hồn riêng lẻ vậy.''

Thợ cơ khí khoanh tay trước ngực, ánh mắt hiền từ dõi qua từng gương mặt đang chăm chú làm việc. ``Các cháu đã dành nhiều tâm huyết lắm đấy. Chú không nghĩ một bữa tiệc sinh nhật mà làm đến mức như thế này.''

Sora nghiêng đầu mỉm cười, đáp lại: ``Bởi vì\dots\ chúng cháu rất yêu quý Kuon-senpai. Đây là lần đầu tiên bọn cháu được mừng sinh nhật cùng anh ấy - nên ai cũng muốn làm hết mình ạ.''

``Thực ra thì,'' A-chan ngượng ngùng lên tiếng, ``ý tưởng chia nhóm và trang trí tường theo phong cách cá nhân cũng không hoàn toàn là của cháu đâu. Cháu nghe loáng thoáng nhóm kia nhắc đến ý đó, rồi tụi cháu học hỏi luôn.''

``Có khi nào\dots\ họ cũng đang làm điều tương tự không nhỉ?'' nàng Thần tượng mỉm cười, liếc về phía tường của nhóm mình.

Kaien mỉm cười, liếc nhìn sang phía bên kia của căn phòng, nơi vẫn còn khuất sau vách tường ``Cuối cùng thì\dots\ là cả hai nhóm đều đang tương trợ lẫn nhau. Như thể thần giao cách cảm vậy.''

Ryujin gật đầu, một tiếng cười trầm ấm vang lên: ``Hay lắm, đúng là không khí của một đại gia đình. Hẳn là thằng bé sẽ cảm động lắm đây.''

Vừa lúc đó, từ phía bên kia của căn phòng, một giọng nói rộn ràng vang lên: ``Này, này! Bên tường giữa ngó qua đây làm gì vậy! Định `đạo văn' bên nhóm tụi này hả?''

Câu nói nửa đùa nửa thật khiến cả nhóm cười rộ. Một vài người trong nhóm 2 tò mò ngó đầu ra. Khi thấy Ryujin và Kaien đang đứng cạnh Sora và A-chan, họ bỗng khựng lại, như bị bắt quả tang khi đang vẽ bậy lên tường nhà người khác.

Haato, với tay vẫn còn dính đầy kim tuyến, lúng túng: ``A-A-nô\dots\ Chú Ryujin-san, cô Kaien-san\dots!?''

Towa thì chỉ gật đầu ngắn gọn, nhưng vẫn tiếp tục tô những đường viền màu tím đậm trên nền đen - mô phỏng một màn đêm trừu tượng mà cô khẳng định là ``hiện thân sự bí ẩn của Kuon-senpai.''

Kanata đứng bật dậy, nhanh chóng giải thích: ``Bọn cháu\dots\ à không, nhóm chúng cháu chọn vẽ lại những khoảnh khắc đáng nhớ của anh ấy, theo cảm nhận riêng của mỗi người. Có người vẽ anh ấy như một người thầy nghiêm khắc, có người lại vẽ như anh hùng, hoặc\dots\ người anh trai dễ dỗi ạ!''

Polka ngoe nguẩy cái đuôi cáo phía sau, đắc ý khoe với cặp mắt long lanh: ``Cháu thì khắc họa nguyên một mê cung! Ý tưởng là `bước vào tâm trí Kuon-senpai' á! Mỗi hành lang sẽ dẫn tới một phiên bản khác nhau của anh ấy!

Kaien bật cười: ``Vậy ra mỗi đứa lại thấy Kuon theo một cách rất khác nhau nhỉ?''

Ryujin gật đầu, ánh mắt ánh lên sự xúc động thầm lặng: ``Nhưng mà nghe cách các cháu nói\dots\ hình như ai cũng đều quý nó cả.''

Nàng Song tính lí nhí, tỏ vẻ rụt rè hẳn: ``V-Vâng ạ\dots\ Chúng cháu mà không quý anh ấy thì đâu vẽ được thế này.''

Nàng Thiên sứ tiếp lời, nhẹ nhàng: ``Thật ra, chính anh Kuon-senpai cũng là người gắn kết mọi người lại, dù cháu nghĩ anh ấy chẳng bao giờ nói ra ạ.''

A-chan nở một nụ cười dịu dàng: ``Đúng thật.Hikawa-san chẳng bao giờ nhận, nhưng ai cũng cảm thấy được.''

Kaien lùi một bước để nhìn toàn cảnh căn phòng, nơi hai bức tường dài đang được dệt nên bằng những nét vẽ chân thành, tự do nhưng vẫn chan chứa tình cảm. Với tinh thần bay bổng của một người yêu thích nghệ thuật, người nội trợ nói: ``\vie{Một bên là cảm nhận, một bên là lý tưởng. Cả hai cùng vẽ về nó, bằng hai ngôn ngữ khác nhau. Mà lại cùng hòa về một điểm.}''

Ryujin đứng cạnh người vợ, thì thầm như nói với chính mình: ``\vie{Cái đó gọi là đồng tâm ý hợp chăng?}''

``Cũng có thể lắm chứ.''

\ct

Ba giờ chiều, một tiếng kể từ khi các thành viên của nhà Tam Giác Xanh tới.

Công đoạn chuẩn bị đang được ráo riết và khẩn trương đẩy nhanh tiến độ trong khắp ngôi nhà. Các nhóm làm đồ nướng, nấu lẩu và trang trí đều chia nhau từng phần công việc một cách kín đáo, ăn ý, như thể đã tập dượt từ trước. Tiếng dao thớt, tiếng cười khúc khích xen lẫn mùi thơm của gia vị, nước lẩu và thịt nướng hoà quyện khắp các tầng, nhưng tuyệt nhiên không ai để lộ ra một dấu hiệu nào có thể đánh động tới cậu Tổng trong nhà.

Ở tầng năm, Kuon và cô em gái của cậu vẫn đang ngồi lặng lẽ trong phòng khách - mỗi người một góc, mỗi người một việc riêng. Cậu chăm chú gõ bàn phím, vừa làm báo cáo cho công ty, vừa thỉnh thoảng liếc sang em gái.

Cậu nhận ra Ryuuhou có vẻ không tập trung xem tivi như mọi khi, tay cứ bấm bấm điện thoại, đôi lúc còn khẽ cười, khẽ chau mày - nhưng cậu cũng chẳng mấy để tâm, chỉ cho rằng con bé đang tám chuyện với bạn học cùng lớp.

Nhưng, điều mà cậu Tổng không thể nào ngờ tới, là cô em gái nhỏ đang âm thầm tham gia vào một ``chiến dịch đặc biệt'' mà chính cậu là trung tâm. Trước đó, trong lúc Kuon đang bận làm việc, điện thoại Ryuuhou liên tục rung lên những tin nhắn từ một cái tên quen thuộc: Ookami Mio, người đã gọi điện cho cô hồi ban sáng.

Chính nàng Sói đen đã chủ động liên lạc với Ryuuhou, dịu dàng gợi ý rằng:

\txtbx{``Kuon-senpai sắp sinh nhật rồi, bọn chị đang bí mật tổ chức, nếu em có thể giúp bọn chị một tay, anh ấy chắc chắn sẽ rất bất ngờ và vui lắm đó!''}

Ban đầu, cô em gái có phần hơi lưỡng lự, nhưng khi được Sora - người mà cô vẫn còn giữ liên lạc sau dịp Tết năm ngoái - ngọt ngào nhắc tới việc nhóm đang thiếu người phụ làm bánh, cô bé đã đồng ý, đơn giản chỉ vì lý do:``\textit{Giúp một lần cũng có sao đâu, nếu làm bánh thì em biết chút chút, với lại\dots\ cũng muốn xem các chị ấy dạo này như thế nào nữa.}''

Và thế là, cô bé âm thầm nhập hội, trở thành một mảnh ghép bí mật trong kế hoạch chuẩn bị này - mà người anh trai đang ngồi ngay trước mặt vẫn chẳng hề hay biết.

Nhưng, một biến cố đã xảy ra.

Trong lúc mọi thứ tưởng chừng như yên ổn, Kuon khẽ đẩy bàn sát vào tường rồi đứng dậy, lấy chiếc quần thể thao mắc treo gần nhà tắm, chuẩn bị rời khỏi phòng.

Cô em gái của cậu vốn đang mải nhắn tin hẹn thời gian xuống giúp nhóm làm bánh bỗng giật bắn mình, hoảng loạn nhìn theo anh trai, tim như muốn nhảy ra khỏi lồng ngực.

Cô vội vàng giấu điện thoại sau lưng, rồi vội vàng đứng dậy, hỏi dồn dập: ``\vie{Ơ\dots\ Onii-san, anh đi đâu đấy!?}''

\pov{Hikawa Kuon}

Nghe thấy tiếng hoảng loạn của em gái, tôi dừng lại. ``\vie{Anh đi ra ngoài có việc,}'' tôi đáp.

Cơ mà, nhìn kỹ lại thì\dots\ có gì đó khang khác. Thái độ của cô bé rõ ràng có gì đó bất thường, thậm chí tôi còn thấy Ryuuhou lấm tấm mồ hôi, dù điều hòa vẫn còn đang chạy và căn phòng vẫn đang mát rượi.

Chuyện gì xảy ra với em ấy vậy?

``\vie{Anh lên trường, bàn bạc bài tập nhóm với bọn bạn. Kỳ tới có đồ án to mà lại dễ trượt, không chuẩn bị sớm thì toang là cái chắc,}'' tôi nhàn nhạt trả lời, kiểm tra lại xem có quên đồ không.

Ryuuhou lắp bắp gật đầu, nhưng rồi lại thở phào, gượng cười: ``\vie{À, ra là thế\dots\ Em cứ tưởng anh lên bếp tranh phần ăn trước chứ! Hồi nãy bố mẹ bảo có lẩu bò cơ mà.}''

Tôi nhướng mày, khẽ ngạc nhiên với câu nói của em gái mình. ``\vie{Em nói linh tinh gì thế? Anh đã ăn tranh của em đâu mà phải nhặng xị lên vậy? Hôm nay em bị sốt à?}''

``\vie{K-K-Không có đâu, Onii-san! C-Chỉ là\dots!}'' em gái tôi mấp máy môi, cố tìm ra lý do để giải thích. Còn đó là cái gì thì có giời mới biết. Dù gì có khi cũng chỉ là mấy chuyện vụn vặt mà tôi chẳng hề để tâm tới.

Tôi hướng ra phía cửa, tìm tôi dép lê đi tạm. Cô em gái tôi ra ngoài cửa định tiễn tôi đi.

\dots Quái lạ. Em ấy chưa bao giờ làm thế với tôi. Hôm nay mặt trời mọc ở đằng Tây à?

``\vie{Anh đi luôn bây giờ sao? Bao giờ về vậy?}'' Ryuuhou hỏi, giọng có chút hối thúc.

Tôi nhẹ giọng đáp, nhún vai: ``\vie{Chắc cũng muộn, sáu giờ hơn mới về. Thậm chí còn muộn hơn do tắc đường nữa.}''

Nói rồi, tôi xách chiếc máy tính lên, xuống cầu thang. Tôi ngước lên cô em gái lần nữa.

``\vie{Anh nhớ mua thêm bánh sữa chua về nha\~{}! Để em còn ăn ké luôn đó!}''

``\vie{Biết rồi,}'' tôi cười mỉm đáp.

Và thế là tôi sải bước xuống cầu thang, chẳng mảy may nghĩ ngợi gì nhiều. Dù trong lòng có một chút cảm giác mơ hồ rằng, hôm nay không khí trong căn nhà này có gì đó\dots\ khang khác.

Vừa bước xuống hành lang tầng năm, tôi khẽ hít một hơi, bất giác nhận ra một mùi thơm lạ len lỏi từ tầng dưới. Không giống mùi nước dùng hay thức ăn quen thuộc mẹ tôi hay nấu, mà lại có chút gì đó chua chua, thanh thanh, vừa xa lạ vừa kích thích vị giác.

Tôi ngẩng đầu nhìn xuống cầu thang, khẽ nghĩ bụng: ``\textit{Chắc là bác Hijaki-san nấu cơm cho bà nội đây mà.}''

Nghĩ thế, tôi lắc đầu, chẳng bận tâm thêm. Nhưng khi vừa định bước đi, một âm thanh khác lại vang lên - tiếng lộp cộp khe khẽ, nghe như ai đó đang di chuyển hay gõ vào vật gì đó bằng gỗ.

Tôi hơi dừng bước, quay sang ngó nghiêng một vòng. Tiếng đó nhanh chóng tan vào khoảng không, chỉ còn lại mấy tiếng quạ kêu vọng từ cửa sổ hành lang.

Tôi nhún vai, gạt phắt: ``\textit{Chắc lại mấy con quạ gõ vào cửa kính thôi, mùa này bọn nó hay vậy.}''

Tiếp tục bước xuống tầng ba, tôi bất ngờ bắt gặp bóng dáng quen thuộc của bố - ông đang đứng ngoài hành lang, một tay cầm điện thoại, chăm chú bấm bấm, trông có vẻ như đang trả lời tin nhắn công việc.

Tôi hơi ngạc nhiên, lên tiếng trước: ``\vie{Sao bố lại đứng đây thế?}''

Bố tôi ngẩng lên, cười rất tự nhiên, chẳng có gì bất thường, rồi đáp: ``\vie{À, chờ mẹ con hút bụi phòng xong để bố lắp lại cái điều hòa ở đây. Mấy hôm nay nóng bức, bà chịu không nổi.}''

Nghe cũng có lý, tôi cũng chẳng nghi ngờ gì thêm, chỉ gật đầu rồi đáp lại: ``\vie{Con tưởng bố đứng đây lén lút hút thuốc cơ.}''

Bố tôi cười khà khà, lắc đầu: ``\vie{Chật chội thế này thì hút làm sao được con. Đứng hóng gió thôi!}''

Tôi cười nhạt, tay với vịn cầu thang, nói thêm: ``\vie{Con lên trường đây, có chút việc quan trọng. Có thể sẽ về muộn, chắc phải qua sáu giờ tối ạ.}''

Ông chỉ ừ một tiếng gọn lỏn, vẫn giữ nụ cười điềm nhiên quen thuộc, rồi lại cúi đầu xuống tiếp tục bấm điện thoại như chưa từng có gì đặc biệt.

Tôi rời khỏi tầng ba, tiếp tục bước xuống từng bậc cầu thang, lòng vẫn lẩn quẩn với mấy suy nghĩ vụn vặt. Công việc trên công ty thì ngập đầu, bài tập lớn với đồ án thì đang chờ, mà lịch trình học kỳ tới đã sớm được cảnh báo là ``đánh rớt thẳng tay'' nếu không chuẩn bị kỹ.

Nghĩ tới đó, tôi bất giác thở dài. Một tiếng thở dài đầy mệt mỏi với những thứ đang bủa vây xung quanh tôi.

``\textit{Sinh nhật à\dots\ có lẽ đó sẽ là dịp duy nhất trong kỳ nghỉ này để xả hơi trước khi lại quay vào guồng chạy đua với deadline, ở cả trên công ty lẫn ở trên trường.}''

Vừa nghĩ, tôi vừa bước ra cửa chính, mở cổng, rảo bước về phía con đường quen thuộc dẫn tới trường - chẳng hay biết gì về những điều đang xảy ra ở ngôi nhà phía sau lưng.

\pov{-}

Cùng thời điểm mà Kuon bắt đầu xuống thang, đôi mắt của Ryuuhou nhanh chóng lia về phía điện thoại, ngón tay gõ liên tục dòng tin nhắn báo động tới các nhóm chuẩn bị bên dưới:

\txtbx{``Onii-san đang ra ngoài! Báo động, Onii-san đang ra ngoài! Mọi người cẩn thận, giấu hết mọi thứ đi, đừng để anh ấy thấy!''}

% Tại tầng sáu, nhóm làm nướng - Choco, Mel, Okayu, Shion và Luna - đang chuẩn bị dọn dẹp đĩa thịt vừa ướp xong thì đồng loạt nhận được tin nhắn. Không ai bảo ai, mỗi người nhanh như chớp giấu những bát thịt, túi gia vị và dụng cụ vào trong hộc bếp hoặc sau tủ lạnh, những chiếc bao tay còn vương mùi thịt cũng vội vo lại nhét vào túi nilon rồi ném gọn vào ngăn chứa rác kín.

\begin{center}
  \uline{Tầng 6 - phòng bếp nhà Kuon.}
\end{center}

Tiếng *TINH!* vang lên trong điện thoại của cả năm cô gái gần như cùng một lúc. Tất cả mọi người đều lấy điện thoại ra xem, và rồi bầu không khí yên ổn lúc trước bị xé toạc bởi một dòng chữ gọn lỏn.

Shion là người bật dậy đầu tiên, giọng cao vút đầy hốt hoảng: ``\textbf{Ế-Ế-Ế!? Xuống thật á!? Bây giờ luôn hả!?}'', nhưng chính cô phải tự bịt miệng lại để không bị Kuon nghe thấy. Okayu thì suýt làm rơi bát thịt đang cầm, vội ôm chặt lấy nó như bảo vệ bảo vật: ``Thôi toang rồi! Giấu đi, giấu nhanh lên\~{}!''

Mel thì cuống cuồng quơ lấy những bát thịt và túi gia vị, cất vội vào tủ bếp, tay run đến mức suýt làm rơi nắp hộp. Cô không dám ho he một tiếng, đôi mắt quay như chong chóng trước biến cố không ngờ tới này. Choco giật luôn cái túi bao tay từ tay Luna, vừa vo lại vừa thúc giục: ``Mọi người! Cất hết đi! Đừng để sót gì cả, mau lên! Luna-sama, cả em nữa đó!''

Luna thì ngồi thụp xuống, hai tay lóng ngóng ôm mấy cái túi ni-lông, đôi mắt long lanh ngước lên đầy hãi hùng và ngơ ngác: ``Kuonsa-senpai tới ư!? Zấu nó đi đâu được na!? Nhà em đâu có túi ma thụt đâu nanora\~{}!''

Cả năm người hốt hoảng chia nhau, người nhét đồ vào ngăn bếp, người chụp lấy khăn che lên bát thịt, người thì loay hoay giấu dụng cụ vào hộc tủ - mồ hôi túa ra như tắm, nhưng vẫn cố gắng giữ cho mọi thứ thật tự nhiên.

\begin{center}
  \uline{Tầng 4 - phòng bếp nhà Hijaki.}
\end{center}

% Tầng bốn, nhóm nấu lẩu - Rushia, Subaru, Noel, Aki và Suisei - cũng cuống cuồng không kém. Những nồi nước dùng vừa mới ninh xong lập tức bị đậy nắp kín mít, giấu gọn xuống dưới bàn hoặc xếp chồng lên nhau như thể chúng là nồi chưa rửa. Đống nguyên liệu còn lại thì tạm thời được dồn vào tủ lạnh của nhà Hijaki, trong khi Noel thì suýt chút nữa lại ``vung chùy'' dọn đồ cho nhanh nếu không bị Subaru chặn lại kịp.

Không khí ở tầng bốn và cầu thang dẫn xuống tầng ba cũng lộn xộn và náo loạn không kém ngay sau khi tin nhắn cảnh báo được gửi tới.

Ngay khi nhận được hung tin, Rushia đã đọc to tin nhắn, không hề kiêng nể gì: ``\textbf{Kuon-senpai đang xuống nhà! Mọi người, Kuon-senpai đang xuống nhà kìa!!}''

Suisei nghe thấy thế, giật mình đến mức làm rơi con dao làm bếp, mặt mày xanh mét: ``C-Cái gì!? Anh ấy đang xuống thật à!?'' rồi lập tức quẳng con dao vừa làm rơi đó vào nhà vệ sinh, hối thúc: ``N-Nhanh lên! Trốn đi!''

Subaru gần như phắt khỏi ghế, hai tay chống hông ra lệnh như một quản lý thể thao thực thụ: ``Rushia-chan, đậy nồi vào! Aki-senpai, giấu thịt tạm vào tủ lạnh đi! Che nắp lại hết cho chị rồi tìm chỗ giấu, nhanh lên, nhanh lên!!''

Noel đang bưng nồi nước dùng xuống tầng ba thì khựng lại, tay run bắn: ``Trứng lộn chưa giấu! Trứng lộn đâu rồi!?''

Aki vội gom đống trứng vào một túi nilon, nhét tọt vào ngăn tủ: ``Đây, đây này! Mau đóng tủ lại, còn mùi thì làm sao đây!?'' Nàng Ca sĩ ôm mặt, hoảng loạn: ``Làm sao mà che mùi được chứ! Sắp bị lộ rồi!!''

Nghe thấy những tiếng lách cách của nồi niêu xoong chảo và những tiếng tri hô của mọi người, Kaien bình tĩnh bước ra, cầm lấy nồi nước mà nàng Hiệp sĩ đang bê - suýt đổ - đồng thời bình tĩnh ra hiệu: ``Đóng nắp hết lại, giả vờ dọn dẹp! Cô sẽ xử lý phần còn lại!''

Nàng Pháp sư và nàng Vịt nhanh tay bê hết đống bát nguyên liệu xuống tầng dưới, đặt tạm vào các kệ tủ dưới đó, hoặc để dưới gầm bàn và lấy những chiếc khăn không dùng tới để che chúng đi.

Nàng Hiệp sĩ thì luống cuống lau sạch quầy bếp, mặt vẫn không giấu nổi vẻ hoảng loạn, liên tục lẩm bẩm: ``Không phải ngày thường đâu\dots\ đúng ngày sinh nhật lại có nhiệm vụ trốn tìm thế này nữa chứ\dots!''

Mọi thứ chỉ vừa kịp an bài thì tiếng bước chân Kuon từ tầng bốn bên kia cầu thang vọng xuống. Cả nhóm nín thở, tim như muốn nhảy ra khỏi lồng ngực, chỉ còn chờ xem ``bước tiếp theo'' cho cậu Tổng của họ sẽ ra sao.

% Ở tầng ba, nhóm trang trí đang trong quá trình nghĩ ý tưởng cũng vội vàng tháo gỡ tạm thời, cuộn gọn lại, giấu hết vào gầm bàn hoặc sau rèm cửa. Fubuki thì nhanh tay giấu những cuộn ruy băng vào ngăn tủ, trong khi Korone giả vờ cầm khăn lau bàn, mồ hôi vã đầy trán nhưng mặt vẫn cố giữ vẻ ``vô can''.

\begin{center}
  \uline{Tầng 3 - phòng trống.}
\end{center}

Ngay lúc cả nhóm vẫn còn đang mải mê với công đoạn trang trí, điện thoại của Fubuki *Tinh!* lên một tiếng khẽ.

Chỉ thoáng liếc qua màn hình, đôi tai cáo trắng của cô lập tức dựng đứng. Mặt tái mét, Fubuki hét lên hoảng loạn ``\textbf{C-c-c-cảnh báo! Kuon-senpai đang xuống nhà! Trạng thái đỏ toàn đội!!}''

Cả căn phòng lập tức ngừng bặt như thể bị bấm nút ``tạm dừng.''  Tất cả đông cứng tại chỗ trong tích tắc. Korone thì đang cầm chổi sơn, tay vẫn còn dính màu nước, mắt tròn xoe: ``Ể-Ể!? Anh ấy đang xuống sao!?''

Mio giật phắt cây cọ sơn khỏi tay nàng Khuyển, tống vội vào xô nước rồi gấp gáp hạ lệnh: ``Dừng ngay lại! Mức cảnh báo cao nhất! Nguy cơ bị phát hiện trực tiếp!!''

Ở một góc phòng, Lamy đang đứng ghế treo dây đèn lên tường liền giật bắn mình khi nghe tin sét đánh ngang tai. Suýt chút nữa cô trượt chân ngã, đành vội buông dây xuống, nhét vội sau bức rèm cửa như đang giấu tang vật. Watame thì gần như bản năng chui tọt xuống gầm bàn, ôm theo chiếc gối hình bông hoa, chỉ còn đôi sừng đen lò xo nhấp nhô nhè nhẹ trên sàn.

Miko vừa dán xong băng rôn thì mặt cắt không còn giọt máu: ``K-Kami-sama phù hộ cho bọn con\dots\ Đừng để anh ấy nhìn thấy cái này, ne\~{}!!''

Mọi người nhốn nháo trong cơn hoảng loạn. Người gỡ dây, người tháo bóng bay, người thì kéo ghế chắn lối vào như thể đang diễn tập cứu hỏa. Mà đúng là cảnh huống chẳng khác gì huấn luyện phòng cháy chữa cháy - chỉ có điều phiên bản này hỗn loạn, rối rắm và hoàn toàn\dots\ phi chính thống.

Korone luống cuống lôi hết mấy lọ sơn vào túi giấy, mắt đảo quanh tìm chỗ giấu, rồi trong khoảnh khắc bối rối tột độ, cô còn thì thào với chính mình ``Nếu bị phát hiện\dots\ mình uống hết mấy lọ này chắc không để lại dấu vết đâu nhỉ?''

May thay, Mio kịp thời chặn lại, giật tay nàng Chó ra khỏi ý tưởng tự sát: ``Đồ ngốc! Bà là chó, không phải quạ! Uống sơn chẳng giải quyết được cái hú gì đâu!!''

Ở góc khác, A-chan vẫn giữ được vẻ bình tĩnh hiếm hoi. Vừa chỉ huy tạm thời nhóm tẩu tán đồ đạc, cô vừa nhắc nhở bằng giọng khẩn trương mà vẫn chặt chẽ: ``Mọi người, chia ra trốn tạm vào hai phòng vệ sinh! Tôi với Kaien-san sẽ lo phần còn lại!''

Ngay lúc đó, trong bếp, Ryujin và Kaien - vốn đã được con gái báo tin từ trước - chỉ nhìn nhau một cái. Không cần nói lời nào, cả hai đã ngầm hiểu phải làm gì.

Ryujin nhét điện thoại vào túi áo, thong thả ra hành lang tầng ba. Ông tựa vào lan can, vờ như đang chăm chú nhìn gì đó trên màn hình điện thoại, tạo dáng vẻ chẳng khác nào một ông chú đang hóng gió buổi chiều mát mẻ, hoàn toàn vô can.

Trong khi đó, Kaien thì giả vờ như chưa biết chuyện gì đang xảy ra. Bà trở lại khu vực bếp, nhanh tay thu dọn một vài túi nguyên liệu đang ngổn ngang, sập cửa tủ lạnh *RẦM!* một cái nghe rất đời thường, rồi vừa quét sàn vừa huýt sáo khe khẽ. Tất cả nhằm tạo ra khung cảnh của một ``gia đình bình thường,'' như thể chẳng hề có hai chục thần tượng đang trốn chui trốn nhủi trong nhà, khiến bà không khác một bà mẹ xứ Etsunan anh hùng đang giấu bộ đội tránh địch.

Bên trong căn nhà, cả nhóm \hll\ như nín thở, mọi ánh mắt tai nghe dõi theo từng tiếng bước chân vọng xuống từ cầu thang, chờ đợi cuộc chạm trán giữa bố con nhà Hikawa diễn ra.

Bầu không khí căng như dây đàn.

Một vài người vẫn đang thu dọn nốt phần còn lại của chiến trường trang trí, nhưng mắt vẫn liếc về phía cửa phòng tầng ba. Những cái đầu lần lượt thò ra từ sau ghế, sau hộp quà, sau tủ. Có người ngồi xổm, có người bò sát đất, tai dựng lên như radar.

Korone níu lấy tay áo Fubuki, thì thầm, giọng run run: ``\hl{Nè! Nhỡ đâu anh ấy thấy hết mọi thứ thì sao? Kuon-senpai đa nghi lắm á, gâu gâu!}''

Fubuki nuốt khan, mắt vẫn dán chặt vào khe cửa, thì thào đáp lại: ``\hl{Bình tĩnh đi Korone, Ryujin-san đang che chắn cho rồi\dots{} Chỉ cần anh ấy không bước vào phòng là ổn thôi!}''

Sora nấp sau chồng hộp quà, hai tay bấu chặt vào mép hộp, mắt mở to đầy lo lắng nhưng vẫn thì thầm đầy niềm tin: ``Chú Ryujin-san sẽ xử lý ổn thôi\dots\ Kuon-senpai sẽ không nhận ra đâu, chắc chắn mà\dots''

Mio thì cầm chiếc ruy băng dở dang, lặng người, nín thở nghe ngóng, tai sói dựng đứng cả lên, mắt chỉ chờ tiếng chân cậu con trai bước lên hoặc quay đầu rời đi.

Phía xa hơn, Botan khoanh tay, đứng dựa lưng vào tường, miệng nhếch lên cười mỉm, nhỏ giọng chọc nhẹ: ``\hl{Nếu Kuon-senpai mà phát hiện\dots{} chắc phải giả vờ ngất hàng loạt thôi.}''

``\hl{Không có vui đâu, Shishiron! Im lặng đi!}'' nàng Yêu tinh tuyết gắt lên nho nhỏ, hạ giọng ngay lập tức vì sợ lộ, nhưng vẫn lườm Botan một cái sắc lẹm.

Noel và Subaru - những người còn chưa kịp thu dọn xong nồi nước dùng - thì áp sát vào cửa, dỏng tai lên nghe trộm.

Nàng Hiệp sĩ lẩm bẩm: ``\hl{Nghe thử xem anh ấy nói gì, nếu bảo `quay lên lấy đồ' thì\dots{} hết đường cứu luôn rồi\dots}''

Nàng Vịt thấy thế vội bịt miệng đối phương, thì thầm đáp lại: ``\hl{Yên lặng nào! Nghe kỹ, nghe kỹ kìa\dots{} Anh ấy đang nói chuyện với Ryujin-san kìa!}''

Cuộc trò chuyện giữa hai bố con vang lên rõ ràng từng chữ, với cậu Tổng của họ có dặn cho người cha của mình: ``\vie{Con lên trường, có chút việc quan trọng. Có thể sẽ về muộn, chắc phải qua sáu giờ tối.}''

Đáp lại cậu chỉ là một tiếng ``ừ'' của người đàn ông, tỏ vẻ cực kỳ bình thản như không có gì xảy ra.

Nghe đến câu đó, cả nhóm phía sau cánh cửa đều nín thở thêm mấy giây, rồi khi tiếng bước chân của Kuon vang xa dần, từng người một thở phào, cảm giác như vừa sống sót qua một bài kiểm tra tim mạch.

Ngay khoảnh khắc đó, từ tầng ba đến tầng bốn, bầu không khí như được nới lỏng. Ai nấy đều buông thõng vai, cảm giác như vừa trút được một tảng đá khỏi ngực.

Noel suýt bật ra tiếng ``\textbf{Yatta!}'' nếu không bị Subaru bịt miệng kịp thời.

Korone thả người ngồi phịch xuống sàn, ôm ngực rên rỉ: ``Hú hồn! Em cứ tưởng tiêu thật rồi chứ\~{}!''

Fubuki khẽ vỗ vai cô bạn, cười nhẹ: ``Chắc chưa tiêu đâu, nhưng lần sau mà căng vậy thì em chuẩn bị\dots\ bông gòn chèn tim trước cho chắc.''

Từ cửa sổ tầng ba, vài người - trong đó có Sora và Miko - lặng lẽ dõi mắt nhìn theo bóng Kuon sải bước ra khỏi căn nhà, rảo chân đi về phía con đường lớn ven con sông. Đến khi cậu đã khuất hẳn khỏi tầm mắt, tất cả mới dám thả lỏng vai, đồng loạt thở phào, ai nấy đều cảm giác như vừa tránh được một đòn chí mạng.

Ngay sau đó, Ryuuhou gửi dòng tin cuối cùng vào nhóm:

\txtbx{``Onii-san đã đi, chắc trước sáu giờ sẽ không về đâu. An toàn rồi! Mọi người tiếp tục thôi!''}

Nàng Sư tử trắng nhìn theo bóng cậu Tổng khuất dần, hòa chung với đường xe cộ nườm nượp. Cô nhún vai, giọng trầm đều: ``Vòng đầu thì tránh được boss, mấy vòng sau không biết có `flawless' nổi không đây.''

Nghe tới đó, nàng Yêu tinh tuyết - người nãy giờ vẫn ôm gối co ro trong góc - lập tức ngẩng đầu, khuôn mặt tái mét, miệng lí nhí: ``Có khi tôi ngất thật luôn đó\dots\ Bà nói như thể chúng ta đang ở trong game kinh dị á!!''

Cả nhóm nghe vậy đều bật cười khúc khích, sự căng thẳng ban nãy cũng tan bớt phần nào, không khí nhẹ nhõm trở lại.

Cả căn nhà gần như khôi phục lại không khí rộn ràng ban đầu, mọi người nhanh chóng quay lại vị trí và tiếp tục những phần việc dang dở, nụ cười dần quay trở lại trên gương mặt từng người. Bữa tiệc bất ngờ vẫn chưa bị bại lọ, và cuộc chuẩn bị thì lại tiếp tục, khi cơn bão vừa lặng sóng tạm thời.

\ct

Ngay khi không khí đã tạm ổn định trở lại, Aqua được Choco và Mel rủ xuống tầng hai để phụ giúp phần bánh ngọt trong bếp. Cô nàng vẫn còn chút hoang mang sau đợt ``truy sát tinh thần'' vừa rồi, nhưng rồi cũng vui vẻ gật đầu đi theo, để lại phía sau là không gian tầng ba nhộn nhịp tiếp tục bước vào giai đoạn chuẩn bị kế tiếp.

Trên tầng, tiếng bước chân di chuyển lại rộn ràng. Mio và Flare đang cùng nhau chỉnh sửa lại dải ruy-băng màu bạc treo giữa trần nhà, nơi hai dải băng cứ chực rơi xuống vì keo dán không đủ chắc. Nàng Sói đen nhón chân, một tay đỡ lấy mép băng, một tay còn lại với lấy cuộn băng dính mà nàng Dị giới đưa tới.

``Góc này lệch rồi đó Mio-senpai. Phải kéo thêm một chút nữa bên trái.''

``Còn bên phải nữa! Để chị trèo ghế lên gắn lại cho.''

Mio lẩm bẩm vài tiếng nhỏ khi cố gắng điều chỉnh góc nghiêng, mái tóc dài xõa ra phía sau lắc nhẹ theo từng động tác. Cạnh đó, Flare thì vẫn kiên nhẫn giữ đầu còn lại, đôi mắt sáng lên vì tập trung.

Ở phía đối diện, hai người nhà Hikawa Kuon cũng đã bắt đầu hòa mình vào cuộc vui một cách tự nhiên như thể họ cũng là một phần của đại gia đình ấy từ lâu. Kaien thì ngồi cùng Sora và A-chan, cùng nhau gấp từng mảnh giấy thành hoa, xếp thành từng cụm nhỏ trên bàn gỗ.

``Cái này là\dots\ gấp thành hoa sao? Nhìn khó ghê,'' A-chan nghiêng đầu thắc mắc.

``Không sao đâu. Sora-chan gấp đẹp lắm, em cứ nhìn theo là làm được ấy mà,'' người phụ nữ nội trợ, tay vẫn thoăn thoắt gấp giấy.

Sora bật cười nhẹ: ``Cháu học từ hồi cấp hai đấy ạ. Giờ vẫn còn nhớ chút chút.''

Không khí ở góc này khác hẳn với lúc trước: một sự ấm áp, nhẹ nhàng và đượm màu thân quen. Ryujin dù không nói nhiều, vẫn lặng lẽ cầm kéo cắt bìa cứng theo mẫu hình trái tim mà mấy cô gái đưa. Ông không quá khéo tay, nhưng lại rất cẩn thận, từng nét cắt đều vuông vắn, ngay ngắn một cách đáng ngạc nhiên.

``Chú Ryujin-san ơi, chú cắt giỏi thật đấy. Còn thẳng hơn cả máy nữa!'' Flare ghé mắt nhìn sang, không giấu được sự ngạc nhiên.

Ông chỉ gật đầu nhẹ, đôi môi khẽ nhếch lên thành một nụ cười kín đáo. Có lẽ đây là lần đầu tiên kể từ Tết Nguyên Đán năm ngoái, căn nhà Hikawa mới lại đông vui đến vậy.

Tiếng cười khúc khích vang khắp các tầng, đan xen cùng âm thanh lạch cạch của kéo, của giấy gói, của vật dụng được chuyền tay nhau không ngừng. Bữa tiệc bất ngờ vẫn chưa bị bại lộ, và cuộc chuẩn bị thì lại tiếp tục, khi cơn bão vừa lặng sóng tạm thời nhưng niềm vui thì đã bắt đầu lấp ló nơi ánh mắt từng người.

\ct

Một giờ trôi qua trong tiếng cười nói rộn ràng, ngôi nhà nhỏ nhà Hikawa giờ đây đã gần như lột xác hoàn toàn. Những dải ruy-băng đầy màu sắc phủ kín các bức tường, xen lẫn là những nét vẽ tươi sáng bằng cọ quét - công cụ đặc biệt mà cả nhóm đã mang về từ trung tâm thương mại hồi sáng.

Fubuki đang khéo léo vẽ những bông hoa nhỏ dọc viền tường trong khi Ayame tỉ mỉ đính từng sợi dây kim tuyến vào rìa các khung tranh.

``Ayame-chan, bên trái lệch rồi kìa! Để đó chị giữ cho, em dán lại thử xem.'' nàng Cáo trắng gọi với, tay giơ cao một đầu dây.

``Mồ\~{} sao cái này khó thế nhờ!'' nàng Quỷ sừng phồng má, dán rồi bóc, mãi đến lần thứ ba mới vừa ý.

Trong khi đó, ở một góc khác, Kanata và Watame đang quỳ xuống sàn, dùng cọ tô phần chân tường với các hoạ tiết đơn giản những hình tròn nhỏ tượng trưng cho pháo giấy, xen kẽ các ngôi sao. Pekora chạy ngang qua, tay cầm một chùm bóng bay, hô lớn: ``\textbf{Tránh ra, tránh ra nào peko!! Để Pekora ta đây bay ngang cái nào!!}''

``\textbf{Làm gì thế hả!? Đừng dẫm lên màu tụi này vừa vẽ xong chứ!!}'' nàng Thiên sứ hét lên trong khi nàng Cừu ôm lấy những hộp sơn.

``Peko\uparrow peko\downarrow peko\uparrow peko\downarrow! Tránh đường, tránh ra! Đường của tôi hết!!'' nàng Thỏ cười khoái chí, bay vụt qua như cơn gió.

Lamy và Matsuri phụ trách treo các băng-rôn phụ quanh phòng. Trên những dải vải nhỏ ghi dòng chữ như ``Hôm nay là ngày của Kuon-senpai!'' hay ``Đừng để Kuon-senpai phát hiện!''. Đôi lúc, họ còn tranh thủ chạy sang giúp Mio và Subaru treo lại các khung ảnh kỷ niệm hai bên hành lang dẫn ra ban công, mỗi khung là một hình ảnh của các thành viên chụp cùng cậu Tổng trong các hoạt động trước đây.

``Nene-chan, đưa chị cái khung tranh đi!'' nàng Cổ động viên gọi với, tay chỉ về phía chiếc khung đang nằm lăn lóc trên bàn.

``Đây\~{}'' nàng ``Hải cẩu'' nhanh nhảu nói.

Nhưng chính ở bức tường lớn nhất - bức tường trắng phía Tây của phòng khách - cả hai nhóm dường như cùng lúc nảy ra cùng một ý tưởng.

``Khoan\dots\ hay là mình vẽ cả chân dung của anh ấy ở đây luôn nhỉ?'' Haato vừa nói, vừa nhìn lên khoảng trống mênh mông trước mắt.

``Em cũng định nói thế luôn á!'' Lamy tròn mắt. ``{Nếu là cả hai nhóm cùng làm thì chắc chắn đủ sức đó.}''

Nói đoạn, nàng Song tính hô lớn tất cả mọi người lại.

Không cần nhiều lời, hai mươi người tản ra thành từng cụm nhỏ. Ai khéo tay thì phụ vẽ phác hình tổng thể bằng bút chì màu, ai vững kỹ thuật thì vào tô màu. Những người còn lại chuẩn bị bảng pha màu, dọn dẹp nền tường, hoặc truyền dụng cụ.

Kanata và Towa phối hợp khá ăn ý - một người vẽ đường nét chính, người còn lại tô các phần chuyển sắc. Một bên là nét cá tính mạnh mẽ, một bên là sự mềm mại - bức chân dung dần hiện rõ nét như thể chính Kuon đang mỉm cười nhìn họ từ nơi đó.

``Trông ngầu phết đấy chứ ha?'' nàng Ác ma nheo mắt, đứng lui lại vài bước.

``Anh ấy sẽ bất ngờ lắm cho mà xem.'' nàng Thiên sứ vừa nói vừa vẽ tiếp mảng tóc, ánh mắt long lanh như chứa cả sự chờ đợi.

Cuối cùng, khi nét vẽ cuối cùng được đặt lên - một đôi mắt nâu trầm, ánh lên chút dịu dàng đặc trưng - cả căn phòng vỡ oà trong tiếng reo vui. Pekora và Polka đã chuẩn bị sẵn băng-rôn từ trước, chỉ đợi khoảnh khắc này.

``{Lên nào! 1, 2, 3\dots\ Treo!}''

Một dải băng-rôn dài được căng lên giữa hai đầu bức tường: ``Chúc mừng sinh nhật, Kuon-senpai!'

Dưới ánh đèn vàng ấm áp, gương mặt vẽ tay cậu Tổng như sống dậy giữa không gian thân quen ấy. Từng chi tiết, từng nét vẽ - không hoàn hảo, nhưng lại mang đầy tình cảm của hai mươi con người đã cùng nhau dồn hết tâm huyết vào một bức tường.

Và ở đâu đó trong ngôi nhà ấy, không ai để ý rằng thời gian trôi nhanh đến vậy. Bữa tiệc đang đến rất gần.

\ct

Bốn giờ ba mươi chiều.

Từng bước chân dồn dập vang lên từ cầu thang ngoài hành lang, nhóm nấu ăn do Aki và Mel dẫn đầu bắt đầu lần lượt bưng các khay đồ ăn đầy ắp xuống phòng khách. Những nồi nước dùng còn lăn tăn bọt, những bát thịt được ướp thơm phức, những dụng cụ phục vụ cho bữa tối lỉnh kỉnh được mang xuống, đi kèm với những món ăn đi kèm như những đĩa rau củ tươi ngon và được cắt thẳng tắp và phần một nhỏ cơm nắm được nặn đều tay cũng được bưng xuống.

``Tránh đường nào\~{} Đồ ăn đang tới đây!'' nàng Ma cà rồng hô vang, vừa nhún nhảy vừa nâng chiếc nồi lẩu nóng trên bếp điện lên.

``Này, Mel-sama, đừng rung mạnh quá, kẻo đổ nước ra sàn kìa!'' nàng Y tá vội ngăn, mặt tái xanh.

Các thành viên còn lại tạm gác việc trang trí, chạy đến phụ giúp. Flare và Mio nhanh chóng trải khăn lên bàn dài giữa phòng, trong khi Sora phụ chia các phần ăn ra đĩa nhỏ.

``Cẩn thận nha, cái này là bánh của riêng Kuon-senpai á,'' náng Sói nói, vừa đặt chiếc bánh mousse trà xanh mà nhóm nấu bếp tầng dưới làm riêng lên chính giữa.

``Nhìn đẹp ghê\dots\ Thấy đói bụng ghê luôn đó\~{}'' Rushia liếm môi.

Trong lúc đó, những người vẫn đang tô nốt các chi tiết cuối cùng trên các bức tường cũng tranh thủ đổi ca để nghỉ tay. Họ chừa lại đúng mười chỗ trống cho những người đang bận việc khác, đặc biệt là cho Aqua.

Nàng Hầu gái trở lại, hai tay vẫn còn dính vết bột mì, đôi mắt ánh lên niềm hứng khởi: ``Xin lỗiii\~{} Tôi quay lại rồi nè! Chỗ của tôi còn không đó?''

``Còn chứ!'' Ayame giơ tay vẫy. ``Bọn này để đúng chỗ của bà nè. Mau vẽ tiếp đi, Kuon-senpai mà thấy chắc cười xỉu luôn đó!''

Ở góc phòng, A-chan ngồi bên cạnh ông Ryujin và bà Kaien -hai người lớn tuổi vẫn giữ vẻ điềm tĩnh nhưng ánh mắt thì đã sáng long lanh suốt từ đầu buổi. Họ lặng lẽ quan sát các cô gái cười đùa, trêu chọc nhau, cùng vẽ tranh và bưng bê thức ăn như một đại gia đình thực thụ.

``Không ngờ mấy đứa lại thân thiết đến thế.'' người thợ cơ khí khẽ nói, bật cười với không khí khẩn trương cho căn nhà.

``Dù đến từ nhiều nơi khác nhau, nhưng nhìn đám nhóc gắn bó như chị em vậy. Thật hiếm thấy,'' người nội trợ mỉm cười, mắt dõi theo Aqua đang đính sao giấy lên góc tường.

``Với cháu thì,'' A-chan khẽ nghiêng đầu, ``nhìn thấy mấy đứa hợp sức thế này, giống như được chứng kiến một phép màu vậy.''

Cả ba cùng im lặng thêm một lúc, như để lắng nghe rõ hơn tiếng cười của tuổi trẻ.

Khi bức vẽ cuối cùng được hoàn tất, Polka bất chợt reo lên: ``Che nó lại đi! Che lại! Chưa được cho Kuon-senpai thấy chừng nào chưa tới giờ chính thức!''

``Ủa đúng á, để bất ngờ mà!'' Shion vội bật dậy sau chút thời gian nghỉ ngơi. ``Có cái màn đen trong phòng của anh ấy đó. Để tôi đi lấy!''

Chỉ mười phút sau, với sự giúp đỡ của Shion, Towa và Nene, một tấm màn đen lớn được kéo căng che phủ kín toàn bộ bốn bức tường, nơi đặt chân dung khổng lồ và băng-rôn chúc mừng. Họ còn dùng băng dính cố định khéo léo phần rìa, bảo đảm rằng không một góc nào bị lộ ra ngoài.

``Ổn rồi! Tới lúc đó chỉ cần kéo một cái là cả thế giới lộ ra luôn\~{}'' nàng Cáo sa mạc vừa nói vừa làm động tác kéo màn đầy kịch tính.

``Ừm, giống mấy chương trình truyền hình á. Mình sẽ cho anh ấy bất ngờ đến phát khóc luôn.'' Nene nắm tay cô bạn, mắt long lanh.

Không khí trong phòng dần lắng xuống sau những tiếng cười. Một cảm giác hồi hộp nhưng cũng đầy ấm áp đang dần lan toả.

\ct

Năm giờ chiều.

Sau khi nhận chiếc bánh sinh nhật của nhóm làm bánh kem, Shion đặt nó xuống giữa căn phòng tầng ba và bắt đầu lẩm nhẩm câu chú.

``\textbf{Phóng lớn!}''

Một luồng sáng tím nhạt loé lên giữa hai tay cô bé phù thuỷ, rồi ngay sau đó\dots\ *bụp!* một vật thể khổng lồ hiện ra ngay giữa phòng: chiếc bánh sinh nhật sáu tầng được phóng đại, ngập tràn hương thơm và sắc màu.

``Ối chà! To quá, yo!'' Ayame há hốc miệng, hai tay dang rộng như muốn ôm hết chiếc bánh.

``Có tới\dots\ hai mươi mấy hương vị luôn!?'' Towa khẽ bất ngờ, vừa đếm vừa ngó nghiêng. ``Tầng trên cùng là mousse sô-cô-la với cam vàng? Đỉnh ghê!''

``Phía dưới nữa nè, có dâu tây, matcha, chanh dây, vani\dots\ còn cả tiramisù nữa!'' Nene hoa mắt liệt kê, suýt thì ngã nhào vào bánh.

Nhìn thấy mọi người đang bàn tán với nhau về chiếc bánh, nàng Phù thủy vội vàng thốt lên, chạy tới chắn trước mặt họ: ``Gượm đã! Ta phải trang trí xong đã rồi mới ăn chứ!''

``Tưởng mọi người đã trang trí rồi chứ?'' Korone nghiêng đầu. ``Tôi nhìn thấy có cả những hạt sô-cô-la với hoa quả trên kia rồi mà?''

``Thì đúng là thế,'' Okayu đáp lại, ``nhưng mà ta còn phải làm thêm một số thứ nữa. Mọi người cùng nhau làm thôi!''

Các cô gái nhóm 2 nhanh chóng bắt tay vào hoàn thiện những nét chấm phá cuối cùng cho chiếc bánh: bước hoàn thiện đúng nghĩa.

Polka thò tay vào giỏ trang trí, lôi ra một lọ kẹo lấp lánh hình sao, hí hửng rắc lên tầng ba của bánh. ``Hưm hưm\~{} Một chút ngôi sao cho Kuon-senpai!''

Nhưng, có lẽ vì cô nàng quá phấn khích mà tay bắt đầu run rẩy.

``Ôi cái đệch!''

*BỤP!*

Kẹo rơi thẳng vào tầng dưới, trộn lẫn với phần phủ sô-cô-la nguyên bản. Shion nhìn thấy thế, liền nhảy dựng lên: ``Ê! Phần đó là phần dâu tây mà!''

``Thì sô-cô-la ăn với  dâu cũng hợp mà, ha ha ha\dots'' Polka cười trừ, vội dùng thìa hớt bớt nhưng càng hớt càng lem nhem.

Cùng lúc đó, Nene đang loay hoay với hai cây nến số to tướng mà Botan đã chuẩn bị từ trước: một cây số 2 màu hồng, và một cây số 5 màu xanh.

Nàng ``Hải cẩu'' gãi đầu gãi tai, liếc qua liếc lại hai cây nến, tự vấn: ``Ờ\dots\ để bên nào nhỉ? 25\dots\ hay 52?'

``Để bên nào cũng được, miễn là có đủ số thôi!'' nàng Cáo sa mạc đáp lại, tay vẫn đang cố gắng dán cho bằng được cái nơ lên tầng bánh.

``Bà nghĩ kỹ đi nha!'' nàng Sư tử trắng ở dưới gọi vọng lên. ``Cắm sai thứ tự là tự biến mình thành dở hơi đấy!''

``S-Shishiron! Bà cứ lấp lửng thế thì làm sao biết được!''

``Thực ra thì luận dễ mà,'' nàng Yêu tinh tuyết đáp lại, tay vẫn đang cắm bánh quế. ``Chỉ cần nhìn vào diện mạo của Kuon-senpai là ra thôi!''

``Hả?''

``Thì nhìn vào cái mặt của anh ấy mà đoán thôi!'' nàng Cáo sa mạc nhún vai. ``Chẳng phải anh ấy là một người rất thích sự hoàn hảo sao?''

``Đúng rồi đó!'' nàng Quỷ sừng gật đầu. ``Nên chắc chắn là số 25 rồi!''

``Vậy thì cắm số 2 bên trái, số 5 bên phải nha!'' Nene vội vàng cắm hai cây nến xuống, rồi nhảy xuống sàn ngắm nghía.

Để rồi, tất cả mọi người đều phải ngớ người ra vì một sự cố\dots\ ngớ ngẩn. Nàng Họa sĩ (?) đã cắm ngược hai cây nến, khiến cho số 2 nằm bên phải và số 5 nằm bên trái, tạo thành số 52.

``Trời ơi!'' nàng Cáo sa mạc ôm đầu, ``Bà làm cái gì vậy Nene-chi!?''

``Thì tôi có biết đâu!'' Nene bối rối, ``Tôi tưởng là\dots\''

Trong lúc cả nhóm đang loay hoay sửa chữa tai nạn số tuổi, thì ở góc khuất sau tấm rèm, một đôi mắt đỏ và một cặp mắt tím đang lén lút rình mò\dots

``\hl{Ayamesa-senpai\dots{} Bánh kem kìa\dots}'' Luna thì thầm, ngón tay chỉ về phía phần kem bơ chưa trang trí còn sót lại trong thố.

``\hl{Nhìn ngon ghê ha\dots{} Kuon-senpai không biết đâu nhỉ?}'' Ayame nhỏ nhẹ lại, đôi mắt sáng rỡ.

*CHỤT!* Một cú bốc kem nhanh như chớp.

*CHỤT!* Một cú nữa.

``Ngonnnn quáaaa-nanoraaaa\~{}\dots'' nàng Công chúa xứ Kẹo lăn lộn trong sung sướng.

``Cứ như là ninja vâjy, không ai thấy hết,'' nàng Quỷ sừng gật gù tự hào. ``Chị cũng muốn ăn thử một miếng!''

Nhưng\dots

``\textbf{Hai đứa đang làm gì vậy!!!}'' tiếng Shion rống lên như sấm giữa trời quang.

Cả hai nhóc con lập tức đông cứng tại chỗ, môi vẫn còn dính kem, miệng há hốc chưa kịp nuốt.

``Bị bắt quở tang rồi!'' Luna thều thào.

``Chạy thôi\~{}!'' Ayame quăng thìa, kéo Luna chạy trối chết ra khỏi phòng.

Sau tất cả những pha ``tai nạn ngọt ngào'', chiếc bánh cuối cùng cũng được trang trí xong dù có hơi lệch một chút ở tầng ba vì bị đổ kẹo, và một mảng kem ở tầng bốn hơi lõm vào (dù đã cố gắng vá lại bằng một miếng cookie hình Suisei).

Tuy không hoàn hảo tuyệt đối, nhưng ai nhìn vào cũng phải công nhận: đây là một kiệt tác thực sự, một chiếc bánh không chỉ được làm bằng nguyên liệu cao cấp, mà còn đầy ắp sự vui nhộn, ấm áp và tình cảm của cả một tập thể của các cô gái nhà \hll.

Ngay khi chiếc bánh sinh nhật vừa được trang trí hoàn tất, Polka và Nene cùng nhau trải tấm bảng đặc biệt ra giữa căn phòng. Đó là một tấm bảng vải lớn, được nhóm làm bánh kỳ công chuẩn bị từ trước. Nền vải xanh đậm, nổi bật lên hàng chữ lớn ``Happy Birthday Kuon-senpai!'' được ghép từ những mẩu lụa màu sắc sặc sỡ, mỗi chữ lại do một người khác nhau thiết kế, mang đậm dấu ấn riêng biệt. Có chữ được may tỉ mỉ, có chữ vẽ tay bằng màu acrylic\footnote{Còn được hiểu dân dã là màu 3D, là màu được tạo nên những sắc tố có khoáng và hữu cơ, dùng để vẽ lên các chất liệu như gỗ, giấy hay vải. Loại màu này tan trong nước, nhưng khi khô sẽ có khả năng kháng nước.}, có chữ gắn thêm kim tuyến lấp lánh hay thậm chí có cả chữ được đính kèm hình hình bánh cá của Miko.

Ngay bên dưới hàng chữ, họ nhẹ nhàng dán thêm một tờ giấy dài, được chia làm ba mươi mốt ô vuông nhỏ, mỗi ô là một lời chúc riêng của từng người trong đại gia đình nhà Tam Giác Xanh. Mỗi nét bút mang một cảm xúc, một cá tính: có dòng nghiêm túc, có dòng hài hước, có dòng nguệch ngoạc đầy biểu cảm đáng yêu. Một vài người vẽ thêm biểu tượng quen thuộc của mình như chòm sao của Suisei, hình đầu lâu nhỏ của Towa, chiếc nơ của Miko, hay hình chiếc vương miện mini của Luna.

Cả tấm bảng như một bức tranh sống động, tuy lộn xộn mà đầy yêu thương.

Ở bên kia phòng, ông bà Hikawa đứng im lặng nhìn mọi người treo bảng lên tường. Người đàn ông chắp tay sau lưng, đôi mắt vốn sắc sảo thường ngày bỗng dịu lại.

``\vie{Không ngờ là mấy đứa làm kỹ vậy,}'' ông khẽ nói, miệng khẽ cong lên một nụ cười.

``\vie{Ừ. Tưởng là tiệc cho bạn bè thôi, ai ngờ công phu như tổ chức quốc khánh ghê,}'' Kaien tiếp lời, giọng có chút trầm nhưng ấm.

Dù chính họ cũng đóng góp một phần không nhỏ,Ryujin đã dành khoảng nửa tiếng cùng với nhóm của Mel và Choco để chuẩn bị các món nướng ngoài trời, Kaien cặm cụi với Aki và Subaru canh nồi lẩu đặc biệt mà Kuon thích từ bé, và cô em gái Ryuuhou thì âm thầm làm thêm một chiếc bánh nhân bơ trứng nhỏ xinh, đặt riêng trong hộp kính ngoài chiếc bánh sinh nhật khổng lồ mà cô làm với Aqua và Shion, nhưng khi chứng kiến sự tận tâm của những cô gái này, lòng họ vẫn không khỏi ngạc nhiên và xúc động.

Ryuuhou lặng lẽ tiến đến gần tấm bảng, chăm chú đọc từng lời chúc. Đến khi đọc tới một ô chữ viết bằng mực tím, nguệch ngoạc: ``Chúc Onii-san không còn đi ngủ lúc ba giờ sáng nữa nhé! - từ Ryuuhou'', cô khẽ phì cười.

Cô quay sang nhìn chiếc bánh sáu tầng, tấm bảng vải lung linh, tờ giấy lời chúc chen kín từng ô\dots\ rồi quay lại nhìn đám bạn bè, hay đúng hơn là đông nghiệp của anh trai mình, những người tưởng như thuộc một thế giới xa lạ, nay lại trở nên thân thuộc đến kỳ lạ.

``\vie{\dots Có lẽ Onii-san sẽ bất ngờ tới phát khóc lắm đây.}''

Kaien khẽ đặt tay lên vai con gái, vỗ nhẹ. Ryujin thì đứng khoanh tay nhìn ra cửa sổ, nơi ánh nắng chiều đang bắt đầu đổ xuống khu vườn sau nhà.

Không khí trong căn phòng tầng ba, lúc ấy, vừa rộn rã vừa ấm áp, tựa như một lời chúc sinh nhật thầm lặng mà sâu sắc nhất, không cần nói ra, nhưng ai cũng cảm nhận được bằng cả trái tim.

\ct

Sáu giờ chiều.

Như vậy là mọi công đoạn cuối cùng cũng đã hoàn tất.

Chiếc bánh sáu tầng lộng lẫy được đặt ngay chính giữa bàn dài ở tầng hai và được đặt trong hộp kính, che chắn cẩn thận bằng một màn đen, tấm bảng "Happy Birthday Kuon" được phủ màn  cẩn thận, dàn đèn led chạy theo hình trái tim gắn sát trần nhà đã được Ayame và Subaru thử nghiệm lần cuối, còn mùi đồ ăn - từ lẩu, món nướng, đến những đĩa sushi và đồ ăn kèm - cứ thế lan tỏa khắp gian phòng.

Ba mươi mốt người, không ai bảo ai, đều cùng tụ tập xuống tầng hai, nơi sẽ diễn ra bữa tiệc lớn nhất mà họ từng góp mặt\dots\ Một bữa tiệc mà nếu so với cái sinh nhật gần đây nhất của YAGOO - sếp lớn của các cô gái, chắc chắn cũng phải ``ngang cơ'' về mức độ hoành tráng và công phu.

Mio đặt tay lên hông, nhìn quanh một vòng rồi nói: ``Giờ ta chỉ cần chờ Kuon-senpai về thôi nhỉ?''

Mọi người gật đầu. Không khí bỗng chốc im lặng trong vài giây, như thể ai cũng đang tưởng tượng ra vẻ mặt ngạc nhiên của Kuon khi bước vào căn phòng ngập tràn ánh sáng và tình cảm này.

Dưới sự chỉ huy nhẹ nhàng của những người chủ nhà, các thành viên nhanh chóng vào vị trí: người thì chỉnh nến, người kiểm tra dây đèn, người gấp khăn, người tắt hết đèn trong nhà để sẵn sàng cho khoảnh khắc bật sáng đồng loạt.

Ở góc phòng, Roboco ngó nghiêng mấy món đồ lạ đang được giấu kỹ trong chiếc thùng lạnh mini.

``Ơ, Miko-chan, đống ly ở trên với túi đá khô dưới này là sao vậy?''

Cô nàng Vu nữ chỉ nháy mắt tinh nghịch: ``Bà cứ chờ đi rồi sẽ thấy! Nói trước thì mất vui lắm, ne!''

Cả phòng râm ran tiếng cười khúc khích. Một vài người còn bắt đầu đặt cược xem ai sẽ là người đầu tiên khiến Kuon bật khóc (Marine và Lamy đang dẫn đầu danh sách ứng viên). Suisei thì kiểm tra lại phần nhạc nền, còn Fubuki đứng cạnh cửa, giữ vai trò ``thám thính viên'' phòng khi Kuon bất ngờ về sớm hơn dự kiến.

``Được rồi! Chúng ta cùng bắt đầu bữa tiệc nào!'' Sora giơ tay lên cao, giọng đầy phấn khích.

Cả căn phòng đồng thanh hô lớn: ``\textbf{Được lắm!!!}''

Giữa tiếng hô dậy sóng đó, A-chan đang kiểm tra phần điện nguồn lần cuối, cũng chỉ biết thở dài, chống tay vào trán, lẩm bẩm: ``Mấy đứa này đúng là rất thích hô khẩu hiệu mà.''

Dẫu vậy, khóe môi của cô vẫn cong lên một cách rất dịu dàng, vì cô hiểu rằng, khoảnh khắc đáng mong chờ nhất vẫn còn ở phía trước.

Bữa tiệc sinh nhật lớn nhất trong đời của Kuon, một bữa tiệc mà đến bản thân cậu Tổng cũng không thể ngờ tới, một bữa tiệc không ai có thể quên được\dots

\dots sẽ bắt đầu ngay khi cánh cửa chính mở ra.

\end{document}